\section{Formalization}
\label{app:formalization}

Every theorem stated in the main text is formalized and
machine-checked in Lean~4 against \texttt{mathlib}, and each
headline result passes a \texttt{\#print axioms} check that depends
only on the three foundational axioms \texttt{propext},
\texttt{Classical.choice}, and \texttt{Quot.sound} --- no unproved
postulate, \texttt{sorry}, or \texttt{native\_decide} enters any
headline theorem. Those checks are not run by hand: each of the
65 headline theorems is wrapped in a \texttt{\#guard\_msgs}
assertion, so a changed axiom set is a compile error rather than a
line of output that a reader must notice. The kernel checks the formal
statements; the correspondence between them and the economic sentences of
the manuscript is the reviewed map recorded in this supplement, not a
kernel fact. The hypotheses of every headline
row and of every main-text result's backing declaration are enumerated from
their elaborated types, structure fields and predicate unfoldings included,
in the repository's \path{docs/HYPOTHESIS-CENSUS.md}; that document, not
this appendix, is where ``the assumption boundary we state'' is checked
against what the Lean actually binds over.

Two boundaries on that claim should be stated at the outset rather
than left to be discovered. First, machine-checking a theorem is not
the same as discharging its hypotheses: several of the results below
are \emph{conditional}, consuming an ordering, foresight premium, or existence
condition that is stated precisely in Lean but derived from the
model's primitives only under further assumptions. Which
hypotheses are primitive and which are derived is recorded in the
development's assumption-surface register rather than inferred from
the theorem statements. Second, a theorem proved at an exhibited
parameter vector certifies that the class in question is inhabited
and nothing more; the results so certified are marked as witnesses
below, and no genericity or necessity claim is made from them. This
appendix records the correspondence between the results stated in
the main text and their Lean identifiers; the identifiers are the
names under which each object is verified in the development, so a
reader with the source can locate and re-check every claim. The
main text names no Lean identifier, so the argument can be read as
economics; the full correspondence is collected here.

The development uses Lean~4, whose trusted-kernel architecture is
described by \citet{deMouraEtAl2015}, together with the mathematical
library documented by \citet{MathlibCommunity2020}. Machine-checked
economic theory has precedents such as \citet{Nipkow2009}, while
\citet{VajjhaEtAl2021} formalize the contraction-based convergence of
value and policy iteration for finite-state Markov decision processes
in Coq. The present development does not import that Coq result: it
builds the required continuous-state and two-player foundations in
Lean, then checks the sanctions-specific theorems, exact witnesses,
counterexamples, and axiom dependencies against those definitions.

\paragraph{The external theory is formalized too.} This paper rests on
standard machinery: monotone comparative statics, dynamic programming,
and the ergodic theory of the state chain. The development does not
merely cite that machinery. It states and machine-checks each result it
uses, in Lean~4, in a separate \emph{foundations} layer of 81 modules,
none of which mentions the sanctions game's parameters or payoffs; the
theorems of this paper are then built on those verified statements
rather than on their informal counterparts. Where \texttt{mathlib}
already supplies the underlying theorem --- Banach's contraction
mapping, Knaster--Tarski --- the module derives the cited statement from
it rather than re-deriving it from scratch.
\Cref{tab:formalized-theory} lists the modules the paper's claims rest
on, what each verifies, and where the paper uses it. Three qualifications belong in the same
breath. First, what is verified is the statement each module records,
which is often a \emph{specialization} of the published theorem --- the
scalar case, the one-dimensional choice, the sufficient direction ---
and in one case a strictly weaker fragment: of \citet{HeSun2017} only
the Skolemization step (per-state argmax $\Rightarrow$ stationary
policy) is formalized, \emph{not} their Theorem~1, whose
measurable-selection fixed-point argument we cite and do not re-prove.
Each such scope restriction is recorded in the row itself. Second, four
modules (marked $\dagger$) are verified but imported by no downstream
theorem: two of them support claims the paper makes only in prose, and
two are older supermodular-games modules retained for backward
compatibility and partly subsumed by the Topkis module. None of the
four is load-bearing. Third, the \emph{general-width reusable} analytic
existence package remains literature-grounded: the Brouwer/Schauder
fixed-point theorem and the \citeauthor{HeSun2017}
conditional-expectation fixed point are consumed from the literature
and are \emph{not} machine-checked
(\Cref{sec:existence:nonvacuity}); \texttt{mathlib} contains neither.
That boundary does not apply to the exact full-width $P_3$ witness,
whose quartic construction, Bellman identities, global leader optimum,
equilibrium existence, and sanction-rule monotonicity are
machine-checked directly. The formal guarantee therefore stops at the
general analytic package, not at the explicit $h=1$ equilibrium.

\begin{table}[p]
\centering
\footnotesize
\begin{tabular}{@{}p{0.30\linewidth}p{0.42\linewidth}p{0.20\linewidth}@{}}
\toprule
\textbf{Source (Lean module)} & \textbf{What is verified} & \textbf{Where used}\\
\midrule
\multicolumn{3}{@{}l}{\emph{Panel A. Monotone comparative statics and lattice fixed points}}\\
\addlinespace[2pt]
\citet[Thm.~4]{MilgromShannon1994}, sufficient direction only\newline\texttt{MilgromShannon}
& Quasi-supermodularity and single crossing give a strong-set-order monotone argmax.
& \Cref{ass:sender}; \Cref{thm:dominance}\\
\addlinespace[2pt]
\citet[Props.~5--6]{QuahStrulovici2007ext}, scalar case --- the unpublished \emph{Some Extensions} notes, whose numbering differs from \citet{QuahStrulovici2009}\newline\path{IntervalDominanceOrder}
& Interval dominance is weaker than single crossing and still yields a monotone argmax (Prop.~5); a derivative sign condition delivers it (Prop.~6).
& \Cref{rem:ido}\\
\addlinespace[2pt]
\citet[Ch.~2: \S2.4, Thms.~2.7.1, 2.4.3, 2.7.6]{topkis1998}\newline\texttt{Topkis}
& Strong set order; the argmax of a supermodular function is a sublattice; monotone selection.
& \Cref{sec:assumptions}; \Cref{sec:mechanism:forces}\\
\addlinespace[2pt]
\citet[Thm.~3.9.2]{topkis1998}, affine-transition case$^{\dagger}$\newline\texttt{TopkisDD}
& A Bellman update preserves decreasing differences with no curvature assumption.
& Not imported (legacy)\\
\addlinespace[2pt]
\citet{Topkis1978}; \citet{zhou1994}$^{\dagger}$\newline\texttt{Supermodular}
& A monotone best response on a complete lattice has least and greatest fixed points.
& Not imported (legacy)\\
\addlinespace[2pt]
\citet[Thm.~4.1]{Vives1990}; \citet[Thm.~5]{MilgromRoberts1990}$^{\dagger}$\newline\texttt{VivesLattice}
& Lattice-game equilibrium exists; a monotone joint map on a product of complete lattices has least and greatest fixed points, which bound every other fixed point (the rationalizable bound is not formalized).
& \Cref{sec:model} (greatest selection); \Cref{sec:existence:nonvacuity}\\
\addlinespace[2pt]
\citet[Thm.~3.1, Lem.~5.1]{BalbusReffettWozny2014JET}, one-dimensional choice\newline\texttt{BalbusReffettWozny}
& Greatest and least stage argmax exist under continuity alone; no strict concavity needed.
& \Cref{sec:model} (sender's myopic best response)\\
\addlinespace[2pt]
Knaster--Tarski, abstract template\newline\texttt{TarskiMPE}
& An order-preserving joint best response on a product lattice has a fixed point.
& \Cref{sec:existence:nonvacuity} (route refuted at the example economy)\\
\addlinespace[2pt]
\citet[Thm.~17.31]{AliprantisBorder2006}, the existence and upper-hemicontinuity clauses; uniqueness is elementary from strict concavity and is not in 17.31\newline\texttt{Berge}
& A continuous objective on a compact set has a maximizer, and a strictly concave one has exactly one.
& \Cref{sec:assumptions}\\
\addlinespace[2pt]
Order topology; no single external source\newline\texttt{MonotoneThreshold}
& Upward-closed condition sets are intervals; the infimum is the tipping pivot.
& \Cref{prop:tipping}\\
\addlinespace[2pt]
In-repo; the engine for \citet[\S II]{Kaempfer1988}\newline\path{MarketClearingMonotone}
& The root of a strictly decreasing excess-pressure function inherits parameter monotonicity.
& \Cref{app:fourgroup}\\
\addlinespace[2pt]
\citet[eqs.~6--9, 11, 15--16]{Becker1983}\newline\texttt{BeckerContest}
& Quadratic-cost pressure contest: unique closed-form Nash; pressure rises with the stake.
& \Cref{sec:fourgroup:contest}\\
\bottomrule
\end{tabular}
\caption{External theory and shared infrastructure re-stated and
machine-checked in the development's foundations layer, Panel~A;
Panels~B--D continue in \Cref{tab:formalized-theory-b} and
\Cref{tab:formalized-theory-cont}. Twenty-nine Lean modules in all, none
of which mentions the sanctions game's parameters or payoffs.
$\dagger$~= verified but imported by no downstream theorem, and so not
load-bearing. Where \texttt{mathlib} supplies the underlying theorem,
the module derives the cited statement from it.}
\label{tab:formalized-theory}
\end{table}

\begin{table}[p]
\centering
\footnotesize
\begin{tabular}{@{}p{0.30\linewidth}p{0.42\linewidth}p{0.20\linewidth}@{}}
\toprule
\textbf{Source (Lean module)} & \textbf{What is verified} & \textbf{Where used}\\
\midrule
\multicolumn{3}{@{}l}{\emph{Panel B. Dynamic programming and fixed points}}\\
\addlinespace[2pt]
\citet[Thms.~3.1, 3.2, 3.3, 4.6--4.7]{StokeyLucas1989}\newline\texttt{StokeyLucas}
& Completeness of the bounded continuous functions under the sup norm (3.1); Banach contraction (3.2); Blackwell's conditions (3.3, \path{blackwell_lipschitz}); the Bellman operator preserves monotonicity (4.6--4.7); the Weierstrass extreme-value bounds the module also ships are classical and not numbered in the source.
& \Cref{sec:existence}; \Cref{app:existencemachinery}\\
\addlinespace[2pt]
\citet[\S5.1, \S7.5]{krawczyk_petkov_multistage}\newline\texttt{MultistageGames}
& Verification theorem; feedback-Stackelberg structure with the leader anticipating the follower.
& \Cref{sec:model}\\
\addlinespace[2pt]
\citet[Thms.~4.2, 4.5]{Curtat1996}; \citet{Bertsekas2012} weighted sup-norm\newline\texttt{Curtat}
& Closed-form row-Lipschitz constants of the coupled Bellman pair, and the Schur conditions.
& \Cref{thm:existence}; \Cref{app:closedforms}\\
\addlinespace[2pt]
Perov's principle, after \citet{Toda2021}; the $2\times2$ spectral-radius equivalence is proved in the module itself (\citealp{HornJohnson1985} for the general statement)\newline\texttt{PerovContraction}
& Spectral radius below one equals the Schur determinant condition; unique joint fixed point.
& \Cref{thm:existence}; \eqref{eq:Schur}\\
\addlinespace[2pt]
\citet[Thm.~5.4.3]{AubinFrankowska1990}, scalar $C^{1}$ case\newline\path{ImplicitFunctionTheorem}
& First-order condition plus strong concavity give a Lipschitz argmax selector.
& \Cref{app:closedforms} (selector constants)\\
\addlinespace[2pt]
\citet[Thm.~4.2(c)]{Hinderer2005}$^{\dagger}$\newline\texttt{Hinderer}
& Closed-form Lipschitz bound for the value function of a discounted program --- in the \emph{state} direction only.
& \eqref{eq:KS-closed} (special case); \Cref{sec:existence:nonvacuity} (scope caveat)\\
\addlinespace[2pt]
\citet{HeSun2017} --- \textbf{Skolemization step only; Theorem~1 is \emph{not} formalized}\newline\texttt{HeSunMPE}
& Per-state argmax existence yields a global stationary policy (classical choice).
& \Cref{sec:existence:nonvacuity} (Thm.~1 cited, not re-proved)\\
\addlinespace[2pt]
In-repo; absent from \texttt{mathlib}\newline\texttt{SupDiff}
& Sup-difference inequality: the algebraic core of the Bellman contraction.
& \Cref{sec:existence}\\
\bottomrule
\end{tabular}
\caption{Foundations layer, Panel~B (continued from
\Cref{tab:formalized-theory}; Panels~C--D in
\Cref{tab:formalized-theory-cont}). The Brouwer/Schauder fixed-point
theorem and the \citeauthor{HeSun2017} conditional-expectation fixed
point are \emph{not} formalized (\Cref{sec:existence:nonvacuity}); of
\citet{HeSun2017} only the Skolemization step is.}
\label{tab:formalized-theory-b}
\end{table}

\begin{table}[p]
\centering
\footnotesize
\begin{tabular}{@{}p{0.30\linewidth}p{0.42\linewidth}p{0.20\linewidth}@{}}
\toprule
\textbf{Source (Lean module)} & \textbf{What is verified} & \textbf{Where used}\\
\midrule
\multicolumn{3}{@{}l}{\emph{Panel C. The long-run law of the political state}}\\
\addlinespace[2pt]
Kantorovich--Rubinstein duality; in-repo\newline\texttt{DualKantorovich}
& Dual metric; a Lipschitz-contracting adjoint gives contraction, uniqueness, geometric convergence.
& \Cref{app:existencemachinery}\\
\addlinespace[2pt]
In-repo; determining-class (CDF) route\newline\texttt{DualKantorovichS}
& The dual metric separates probability laws on $[0,1]$.
& \Cref{app:existencemachinery}\\
\addlinespace[2pt]
In-repo; uniform kernel\newline\texttt{MovingAverage}
& Smearing is nonexpansive and linear, and state-Lipschitz with modulus the continuation value's sup norm over the shock half-width $h$.
& \Cref{sec:existence:nonvacuity} (purification)\\
\addlinespace[2pt]
In-repo; box kernel\newline\texttt{BoxPullback}
& The affine box smear maps $1$-Lipschitz test functions to $(1-\delta)$-Lipschitz ones.
& \Cref{app:existencemachinery}\\
\addlinespace[2pt]
In-repo; construction\newline\texttt{BoxTransfer}
& Transfer operator built as a pushforward; its adjoint identity is proved, not assumed.
& \Cref{app:existencemachinery}\\
\addlinespace[2pt]
Krylov--Bogolyubov; \citet[Thm.~12.10]{StokeyLucas1989}\newline\texttt{KrylovBogolyubov}
& Ces\`aro averages give an invariant probability law on a compact state space.
& \Cref{sec:results} (invariant law); \Cref{app:existencemachinery}\\
\addlinespace[2pt]
In-repo; state-dependent drift\newline\path{ClosedLoopBoxTransfer}
& An invariant law exists for every continuous equilibrium drift; uniqueness only if it contracts.
& \Cref{sec:results}; \Cref{app:existencemachinery}\\
\midrule
\multicolumn{3}{@{}l}{\emph{Panel D. The domain layer (no external source)}}\\
\addlinespace[2pt]
In-repo\newline\texttt{UnitInterval}
& Notation and the state/action space $[0,1]$.
& Throughout\\
\addlinespace[2pt]
In-repo\newline\texttt{UnitIntervalProjection}
& Clipping onto $[0,1]$ is continuous and monotone; the state space is compact.
& Throughout\\
\bottomrule
\end{tabular}
\caption{Foundations layer, Panels~C--D (continued from
\Cref{tab:formalized-theory-b}).}
\label{tab:formalized-theory-cont}
\end{table}

\noindent
Two conventions in the table below repay attention, because both have
already misdirected a reader of this development.

\emph{Layer is not directory.}  A Lean file under \path{Examples/} need not
contain point checks.  The symbolically parameterized family of
\Cref{prop:symbolic-family} lives there because it fixes a functional form
for the stage payoffs, not because its theorems are instances: the
\texttt{symbolic\_*} results are quantified over \emph{every} member
satisfying the family's defining inequalities and license universal
statements about the family.  Where a result
has both a family-level and a numeric backing, the table lists the family
level first and marks the numeric one as an instance.

\emph{Instances are labeled.}  A row marked ``non-vacuity instance'' or
``numeric instance'' discharges existence at one economy and licenses
nothing beyond that.  It is never the primary backing for a universal
claim in the text.

% A longtable: the list is several pages long and a tabular cannot break.
\begingroup
\small
\begin{longtable}{@{}p{0.24\linewidth}p{0.70\linewidth}@{}}
\multicolumn{2}{c}{\textbf{Headline results and their primary Lean identifiers.}}\\[3pt]
\toprule
\textbf{Result} & \textbf{Primary Lean identifier}\\
\midrule
\endfirsthead
\multicolumn{2}{c}{\textbf{Headline results and their primary Lean identifiers} (continued)}\\[3pt]
\toprule
\textbf{Result} & \textbf{Primary Lean identifier}\\
\midrule
\endhead
\midrule
\multicolumn{2}{r}{\emph{continued on the next page}}\\
\endfoot
\bottomrule
\endlastfoot
\Cref{prop:entrenchment-welfare} (national welfare at the dual-basin economy; Online Appendix)
& (i) \path{db_political_parts_zero_at_full_trade}, \path{db}$\sigma$\path{_eq_zero_of_le};
(ii) \path{db_entrenchment_loss_after_episode};\newline
(iii) \path{dualRep_rule_eq_zero}, \path{dualRep_rule_eq_zero_of_carrier}\\
\Cref{prop:myopic-benchmark} (the myopic benchmark; main text)
& \path{myopic_dynamic_equilibrium_separation_generalF_stage}, clause (2);\newline
\path{myopicEquilibriumTopS}, \path{myopicPathS_le_geom}, \path{myopicPathS_tendsto_zero}\\
\Cref{thm:regime-separation} (foresight against the benchmark; main text)
& \path{myopic_dynamic_equilibrium_separation_generalF_stage}, clauses (3)--(4);\newline
at every full-non-compliance equilibrium \path{defiance_side_of_policy_zero}; the cap below one under\newline
interiority \path{myopic_dynamic_equilibrium_separation_generalF}\\
\Cref{thm:four-cells} (four cases; discussion)
& \path{four_cell_generalF'}, \path{four_cell_generalF'_of_response_monotone}, \path{four_cell_net_generalF}\\
\Cref{prop:lock-in} (lock-in; main text)
& (i)--(iii) from (P1)--(P3) and (P5), no (P4): \path{standardCompliance_persistence_generalF_from_sender_package};\newline
entry: \path{standardCompliance_min_drift_shock_generalF_from_sender_package_of_preBasinAscent}\\
\Cref{thm:persistence} (persistence)
& \path{standardCompliance_persistence_generalF_from_package}; (i)--(ii) without A1:\newline
\path{standardCompliance_persistence_generalF_level_star}\\
\Cref{thm:dominance} (dynamic versus myopic sanction at a fixed target policy)
& \path{theorem2_dynamic_dominates_myopic_at_same_alpha}; every stage maximizer:\newline
\texttt{stage\_max\_le\_$\sigma$\_dyn\_top}, \path{stageTopRule_le_dynamicTopRealizedRule}\\
\Cref{thm:shock} (shock hysteresis)
& \path{standardCompliance_min_drift_shock_generalF_from_sender_package_of_preBasinAscent} (no P4)\\
\Cref{cor:shock-transition-strength} (transition strength)
& \texttt{fullShockDrift\_minima\_horizon\_}\newline
\path{mono_of_top_transition_le}\\
\Cref{prop:longrun} (long-run lock-in)
& \texttt{dynamicTopLockInPathF\_tendsto\_steadyState\_full\_}\newline
\texttt{of\_strictAdjustment}\\
\Cref{thm:existence} (Schur-condition contraction)
& \path{exists_feedbackStackelbergTop_from_valueClass2_with_membership_schur};\newline
coincidence clause: \path{feedbackStackelbergTop_unique_in_carrier_schur}\\
\Cref{thm:corner-existence} (basin-corner existence)
& \path{feedbackStackelbergTop_corner_of_substituted_stage_le} (weak form; the strict form
\path{feedbackStackelbergTop_corner_of_substituted_stage_strictDecr} is its corollary)\\
\Cref{thm:lowbasin} (mirror low basin)
& \path{standardCompliance_low_basin_generalF_of_sender_level}; high orientation\newline
\path{standardCompliance_persistence_generalF_from_sender_package}\\
\Cref{prop:tipping} (derived tipping thresholds)
& \path{target_tipping_threshold}, \path{sender_tipping_threshold}\\
\Cref{prop:flip-route} (derived-basin route)
& \path{standardCompliance_three_goals_from_stage_flip_generalF_of_preBasinAscent}; non-vacuity: the family\newline
\path{sfFCert_three_goals_generalF}, its affine member \path{sfCert_three_goals_generalF} and its\newline
non-affine member \path{sfqbCert_three_goals_generalF}\\
\Cref{cor:pathwise-dominance} (pathwise ordering)
& \path{pathwise_dominance_generalF}; affine fixed-start form\newline
\path{standardCompliance_three_goals_local_regime_comparison}\\
Affine invariant law of the shocked state (\Cref{sec:results})
& \path{boxTransfer_existsUnique_invariant_attracting}\\
Uniqueness within the invariant sets (\Cref{sec:existence}; the Lean names say carrier)
& \path{feedbackStackelbergTop_unique_in_carrier_schur}\\
Distortion identity (\Cref{app:fourgroup}, \eqref{eq:planner-distortion})
& \texttt{plannerW\_distortion}\\
Planner flip boundary (\Cref{app:fourgroup})
& \path{plannerW_lockin_dominated_iff}\\
Discounted lock-in inefficiency (\Cref{app:fourgroup})
& \path{fgc_planner_inefficiency_equilibrium}\\
Initial-condition selection, general $\F$ (\Cref{sec:mechanism:bistability})
& \path{equilibriumTop_initial_condition_selects_generalF_of_sender_basins} (sender halves);
joint basins \path{equilibriumTop_initial_condition_selects_generalF}\\
\quad non-vacuity instance (affine $\F$, $\delta=1/10$)
& \path{dualParams_initial_condition_selects}\\
\Cref{prop:order-selection} (order-selected steady state)
& \path{equilibriumTop_tendsto_order_selected_fixedPoint_of_stepMonotoneOn} (monotone closed loop);
monotone-rule form \path{equilibriumTop_tendsto_order_selected_fixedPoint_of_continuousOn'}\\
\Cref{cor:order-selection-identified} (limit identified)
& \path{equilibriumTop_tendsto_order_selected_fixedPoint_identified_of_stepMonotoneOn};
monotone-rule form \path{equilibriumTop_tendsto_order_selected_fixedPoint_identified_of_continuousOn'}\\
\Cref{prop:derived-interval} (attraction on a derived interval)
& \path{equilibriumTop_attraction_on_derived_interval_of_stepMonotoneOn}; the state half, without
(F5): \path{equilibriumTop_attraction_up_to_least_steadyState_of_stepMonotoneOn};
monotone-rule form \path{equilibriumTop_attraction_on_derived_interval_of_invariant}\\
\Cref{prop:derived-dividing-line} (the dividing line is derived)
& \path{equilibriumTop_derived_dual_basins_unit} (no maps-into condition); window forms
\path{equilibriumTop_derived_dual_basins}, \path{equilibriumTop_derived_dual_basins'}\\
\quad identification of the two crossings
& \path{leastWeakAbove_spec}, \path{greatestWeakBelow_spec}\\
\quad non-vacuity instance
& \path{dualParams_derived_dividing_line}\\
\Cref{prop:derived-crossings} (general form)
& \path{derivedThresholds_eq_iff_unique_fixedPoint}, \path{leastWeakAbove_le_greatestWeakBelow},\newline
\path{threshold_below_tendsto_left_of_lowerBound},\newline
\path{threshold_above_tendsto_right_of_upperBound}; the two crossings as above\\
\Cref{thm:b2-benchmark} (positive-discount benchmark)
& \path{b2NonMono_three_goals_dominance}\\
\Cref{prop:primitive-persistence} (primitive persistence)
& \path{persistence_generalF}; on the core with $n+v\ge0$:\newline
\path{SymbolicEconomyFPersist.persistence_generalF}; semiconcave sender value\newline
under a corner curvature bound: \path{persistence_generalF_semiconcave}\\
\Cref{prop:equilibrium-patience-feedback} (patience feedback)
& \emph{family level:} \path{SymbolicEconomyF.equilibrium_feedback_decomposition},\newline
\path{lockin_amplification}, \path{lockin_strict_iff}, \path{lockin_strict},\newline
\texttt{VT\_$\beta$S\_invariant}\\
\quad affine family, closed-form values
& \path{symbolic_equilibrium_feedback_decomposition},\newline
\path{symbolic_lockin_amplification}, \path{symbolic_lockin_strict_iff}\\
\quad general-$\F$ form, no functional form assumed
& \texttt{senderValue\_patience\_decomposition\_}\newline
\path{of_bellman}, \texttt{senderValue\_lockin\_}\newline
\path{amplification_of_bellman}\\
\quad underlying algebra
& \texttt{PatienceDecomposition.}\newline
\path{difference_decomposition}\\
\quad numeric instance ($\beta_S$-interval)
& \texttt{eeCertB\_equilibrium\_feedback\_}\newline
\path{decomposition}\\
\Cref{prop:discounted-stationary-welfare} (discounted welfare)
& \path{FourGroupCoeffs.stationaryPlannerValue_lockin_dominated_iff} (every coefficient vector)\\
\Cref{prop:affine-welfare-gap} (welfare gap, general transition; affine instance)
& \path{FourGroupCoeffs.pathValue_cornerF_gap},\newline
\path{pathValue_cornerF_gap_pos_all_discounts_iff_of_params};\newline
affine: \path{pathValue_affineCorner_gap},\newline
\path{affineCornerGap_pos_iff},\newline
\path{affineCornerGap_pos_all_discounts_iff} (all over \path{FourGroupCoeffs})\\
\quad the game does not certify the sign (hawk rescaling)
& \path{exists_hawkRescale_cornerF_gap_neg}, \path{exists_hawkRescale_cornerF_gap_pos};\newline
member $\phi=5$: \path{fourGroupCertHawkScaled_gap_neg}\\
\Cref{prop:planner-optimum-lower-bound} (planner optimum, exact loss)
& \path{plannerOptimalValue_eq_of_upperChords} ($B\ge0$),\newline
\path{plannerOptimalValue_eq_of_nonpos} ($B\le0$),\newline
\path{plannerOptimalValue_eq_affineF}, \path{plannerOptimalValue_loss_of_upperChords};\newline
certificate instance \texttt{fgcPlannerOptimalValue\_eq},\newline
\path{fgc_plannerOptimalValue_gap_eq_equilibrium}\\
\Cref{def:interest-group-representation}, \eqref{eq:welfare-identity} (representation; accounting identity)
& \path{InterestGroupRep}, \path{payoffs_add_eq}\\
\Cref{prop:welfare-rescaling} (rescaling one group)
& \path{exists_rescaleS_ranks}; \path{rescaleS}, \path{rescaleS_welfare}\\
\Cref{prop:welfare-nonidentification} (welfare non-identification, benchmark against forward-looking)
& \path{ChordEconomyRelief.exists_reps_reversing_myopic_forward_ranking}\\
\Cref{prop:welfare-nonidentification-four} (welfare non-identification across the four cases)
& \path{ChordEconomyCs.exists_reps_reversing_regime_ranking}\\
\Cref{prop:zero-sanction-planner} (zero sanctions socially optimal)
& \path{zeroSanction_planner_optimal}\\
\Cref{prop:planner-efficient-persistence} (persistent play socially optimal)
& \path{ChordEconomy.planner_efficient_rep_at_policy_zero},\newline
\path{ChordEconomy.exists_planner_efficient_rep}\\
\Cref{lem:welfare-completion} (interest-group microfoundations)
& \path{govS_eq}, \path{govT_eq}, \path{welfare_eq},\newline
\path{exists_welfareCompletion},\newline
\path{exists_completion_bJ_eq_bK_iff}\\
\Cref{prop:family-welfare-nonidentification} (family-wide welfare non-identification)
& \path{exists_completions_reversing_regime_ranking_of_response_monotone} (non-decreasing response);\newline
under the increasing-differences condition \path{exists_completions_reversing_regime_ranking}\\
\Cref{cor:planner-efficient-boundary} (planner-efficient boundary, every full-non-compliance equilibrium)
& \path{exists_planner_efficient_completion_of_policy_zero},\newline
\path{exists_planner_efficient_completion},\newline
\path{exists_planner_efficient_completion_at_equilibrium},\newline
\path{flat_of_margin_of_r_eq_zero},\newline
\path{placableBoundary_planner_efficient_at_equilibrium}\\
\Cref{prop:steady-gap}, \Cref{rem:wedge-sharp} (long-run gap; the premium bound attained)
& \path{fixedPoint_ge_limit10_add_gap}, \path{fixedPoint_ge_steady10_add_gap}, \path{steadyGap_eq},\newline
\path{dynamicPath_limit_ge_limit10_add_gap_of_response_monotone},\newline
\path{dynamicPath_limit_ge_steady10_add_gap_of_response_monotone},\newline
\path{dynamicPath_limit_ge_steady10_add_gap} (form under the increasing-differences condition),\newline
\path{placableBoundary_wedge_attained},\newline
\path{wedge_coefficient_cannot_be_raised}\\
\Cref{prop:branch-invariance} (path independent of the target's parameters)
& \path{senderValue_eq_of_defiance_section_eq},\newline
\path{equilibriumRule_eq_of_defiance_section_eq},\newline
\path{defiance_branch_invariance_generalF},\newline
\path{dynamicSanction_eq_of_sender_side_eq}\\
\Cref{prop:relief-participation} (relief participation)
& \path{relief_participation_of_policy_zero_at_defiance},\newline
\path{relief_participation_of_policy_zero},\newline
\path{senderReliefValue_add_loss_le},\newline
\path{targetReliefValue_add_cost_le},\newline
\path{relief_joint_loss},\newline
\texttt{no\_date\_zero\_transfer\_supports\_}\newline
\path{relief_at_defiance},\newline
\path{no_date_zero_transfer_supports_relief},\newline
\path{relief_participation_of_margin}\\
\quad refutation outside the entrenchment condition
& \path{lossOffset_relief_period_benefits_target}\\
\Cref{prop:limit-dependence} (limit dependence)
& \path{rulePathF_limit_isFixed},\newline
\path{exists_pos_limit_sanction_of_limit_dependence},\newline
\path{limit_eq_zero_of_limit_sanction_zero} (zero-sanction descent)\\
\Cref{prop:no-cycles} (no conflict cycles under a non-decreasing rule)
& \path{rulePathF_no_cycles}, \path{rulePathF_ruleSanctionsF_tendsto},\newline
\path{rulePathF_const_of_return}, \path{rulePathF_periodic_const},\newline
\path{equilibrium_no_cycles} (\path{Theory/LockIn/NoCycles.lean}); the four\newline
cases: \path{myopic_no_cycles}, \path{senderOnly_no_cycles},\newline
\path{targetOnly_no_cycles}, \path{dynamic_no_cycles}\newline
(\path{Examples/PlacableSenderNoCycles.lean})\\
\Cref{prop:target-two-sided} (two-sided target value)
& \path{VTstar_lipschitz}, \path{VTstar_continuation_bracket},\newline
\path{VTdef_lipschitz}\\
\Cref{prop:certified-equilibria} (certified equilibria)
& \path{senderValue_lowerSlope_of_policy_chord},\newline
\path{V_S_lowerSlope_of_antitone},\newline
\path{equilibrium_sanction_ge_wedge_of_antitone}\\
\Cref{prop:bracket} (the bracket, general transition)
& \path{targetObj_strictAntiOn_of_sharp}, \path{VTdef_lowerSlope},\newline
\path{targetObj_antitoneOn_of_sharp_le}, \path{reliefBestReply_of_relief_side_le},\newline
\path{equilibriumTop_sharp}, \path{gSharp_balance},\newline
\path{zeroLevelAt_lipschitz}, \path{VTrel_lipschitz},\newline
\path{reliefEquilibriumTop},\newline
\path{defiance_not_maxOn_of_relief_side},\newline
\path{bracket_sides_disjoint}, \path{collapse_sides'},\newline
\path{exists_bracket_both_sides_nonaffine}\\
\Cref{prop:monotone-uniqueness} (uniqueness in the monotone class)
& \path{target_lowerSlope_of_rule_monotone},\newline
\path{policy_eq_zero_of_rule_monotone},\newline
\path{unique_equilibrium_of_rule_monotone}\\
\Cref{prop:patience-existence} (a more patient target keeps the non-compliance equilibrium)
& on the relief class, with no (T6):\newline
\texttt{exists\_defiance\_equilibrium\_mono\_$\beta$T},\newline
\texttt{SameBut$\beta$T.balance\_of\_balance},\newline
\texttt{balance\_of\_le\_$\beta$T}\\
\Cref{prop:security} (the target's minmax payoff)
& \path{reliefPayoff_isLeast_history},\newline
\path{reliefPayoff_isGreatest_history},\newline
\path{reliefPayoff_le_V_T}, \path{reliefPayoff_isLeast}, over any economy with the two lower
slope bounds, the three target bounds and the strong entrenchment condition (\path{Theory/Equilibrium/SecurityCore.lean})\\
\Cref{prop:identification} (uniqueness of non-compliance among monotone rules)
& \path{unique_equilibrium_of_rule_monotone},\newline
\path{policy_eq_zero_of_rule_monotone},\newline
\path{no_zero_rule_equilibrium}\\
\Cref{prop:collapse-threshold} (threshold discount factor for full non-compliance)
& \path{collapse_equilibrium_threshold_le_tight}, \path{collapse_defianceExistenceSet},\newline
\path{collapse_exact_thresholds}, \path{collapseRecert},\newline
\path{collapse_equilibrium_threshold_le} (attained; headline row 59),\newline
\path{exists_defiance_equilibrium_at_betaStar},\newline
\path{collapse_equilibrium_threshold} (strict, $\beta_T\ne\beta_T^*$), \path{collapse_sides'}\\
\Cref{prop:defiance-closed} (existence of non-compliance closed in $\beta_T$)
& \path{defianceBestReply_of_tendsto},\newline
\path{exists_defiance_equilibrium_of_tendsto},\newline
\path{defianceExistence_sInf_mem}\\
\Cref{prop:exact-thresholds} (exact thresholds in $\beta_T$, general $\F$)
& \path{exists_relief_equilibrium_iff}, \path{VTrel_monotone},\newline
\texttt{reliefBestReply\_anti\_$\beta$T}, \path{reliefBestReply_of_tendsto},\newline
\path{reliefExistenceSet_eq_Ioc}, \path{defianceExistenceSet_eq_Ico},\newline
\path{reliefExistence_lt_margin}\\
\Cref{prop:affine-window} (the exact intermediate range at the affine member; an instance of the row above, affine $\F$)
& \path{aw_exact_thresholds}, \path{aw_general_hypotheses},\newline
\path{aw_zero_rule_iff}, \path{aw_full_defiance_iff},\newline
\texttt{aw\_window\_at\_$\beta_1$}, \path{aw_continuum}\\
\Cref{prop:core-outside-chord} (the non-compliance core outside the class; non-vacuity instance)
& \path{defiance_core_outside_chord_class},\newline
\path{oc_not_chord}\\
\Cref{prop:interior-witness} (a strict interior optimum in the class; non-vacuity instance)
& \path{strict_interior_equilibrium_in_chord_class},\newline
\path{qc_growth}\\
\Cref{prop:welfare-interval} (partial welfare identification)
& \path{regimeValue11_sub_10_eq}, \path{dS_ge},\newline
\path{hawkBounded_nonempty_iff},\newline
\path{gap_mem_interval}, \path{gap_attained_withSender},\newline
\path{forall_gap_pos_iff}, \path{forall_gap_neg_iff},\newline
\path{exists_both_signs_iff},\newline
\path{exists_withSender_weightBounded_reversal}; each case-level theorem also with a\newline
non-decreasing response in place of the increasing-differences condition, e.g.\ \path{gap_mem_interval_of_response_monotone}\\
\Cref{prop:history-deviations} (pure history-dependent deviations)
& \path{leader_optimal_against_history_deviation_top},\newline
\path{follower_optimal_against_history_deviation_top},\newline
\path{follower_node_optimal_against_history_deviation_top}\\
\end{longtable}
\endgroup

\begingroup
% Lean identifiers use underscores as natural break points.  Keeping this
% redefinition local prevents long formal names from overrunning the margin.
\let\formalizationunderscore\_
\renewcommand{\_}{\formalizationunderscore\allowbreak}

\paragraph{Assumptions: interval dominance and its dischargers (\Cref{sec:assumptions}).}
\begin{itemize}[leftmargin=1.4em,itemsep=1pt]
\item \path{monotone_argmax_of_IDOFamilyOn} --- Verifies that the interval-dominance order delivers the monotone greatest argmax used by A1 and \Cref{rem:ido}.
\item \path{IDOFamilyOn.of_singleCrossingOn} --- Verifies that single crossing implies interval dominance with no side conditions, so it remains a sufficient discharger of A1 (\Cref{rem:ido}).
\item \texttt{hIDO\_on} --- hIDO\_on is the interval-dominance certificate field consumed by the persistence certificate (rem:ido).
\item \texttt{MonotoneTransferOn} --- MonotoneTransferOn is the ordinal escalation-consistency condition (rem:ido).
\item \path{IDOFamilyOn.of_quasiconcave_monotoneTransfer}, \path{hIDO_on_of_quasiconcave_monotoneTransfer} --- Verify that quasiconcavity (single-peakedness) plus escalation-consistency deliver the interval-dominance certificate field (rem:ido).
\item \path{MonotoneTransferOn.of_singleCrossingOn} --- Verifies that escalation-consistency is a theorem (not an assumption) wherever single crossing holds (rem:ido).
\item \path{b2f_not_singleCrossingOn}, \path{b2Cert_not_singleCrossingOn} --- Verify that single crossing fails on the basin at the example economy (rem:ido).
\item \texttt{b2\_hQC}, \path{b2f_monotoneTransferOn} --- Verify single-peakedness (quasiconcavity) and escalation-consistency at the concrete example payoffs (rem:ido).
\item \texttt{b2Certificate}, \path{b2NonMono_three_goals_dominance} --- Verify the full three-goals compliance certificate at the example economy (rem:ido).
\item \path{target_tipping_threshold}, \path{sender_tipping_threshold} --- The two derived-tipping-threshold packages (prop:tipping): under the derivative-free (increasing-differences) form of the stage complementarity and the boundary signs, the explicit inf-constructions \texttt{targetFlipPivot}, \texttt{senderFlipPivot} are interior and each flip set is exactly the upper interval [pivot, 1]; \texttt{targetFlipPivot\_mem}, \texttt{senderFlipPivot\_mem} verify that each pivot attains its own flip condition.
\item \path{standardCompliance_three_goals_from_stage_flip}, \path{StandardComplianceLocalCoreCertificate.of_stage_flip} --- The first, affine form of the derived-basin route (prop:flip-route): the three-goals conclusions at the derived basin [max(xbar\_T, xbar\_S), 1] with the pivots of prop:tipping, clauses (i)--(iii) entering as the hypotheses \texttt{h\_diag}, \texttt{hV\_T\_mono\_carrier}, \texttt{hMyoCeiling}; the target entrenchment-condition block is auto-discharged at c\_T = g\_T = 0.
\item \path{h_diag_of_pivot_action_frontier}, \path{hMyoCeiling_of_stage_strictAntiOn_above_floor} --- The stage-level dischargers of clause (iii) of prop:flip-route and of the diagonal on [sbar, x] at basin states x \textgreater{} sbar (app:dischargers); the second is clause (i) of the route only when xbar\_T \textless{}= xbar\_S; the sigma-difference bridge from the myopic-ceiling clause to fully primitive form remains open, as stated inline.
\item \path{targetCombinedChordOn_of_stage_flip}, \path{StandardComplianceLocalCoreCertificate_MS.of_stage_flip}, \path{standardCompliance_three_goals_from_stage_flip_generalF} --- The general-transition form of the derived-basin route (prop:flip-route): the combined slope condition of ass:target is machine-checked as a \emph{theorem} on the derived basin for any transition satisfying (F2)--(F4) --- the stage part from flip-set membership, the continuation part because s \textgreater{}= x keeps F(x,s) \textgreater{}= F(x,x) = x in the basin range where every invariant-set value V\_T is monotone --- and the three-goals conclusions follow with the trajectory evolved by the general law; the affine wrapper is the specialization retaining the closed-form shock horizon.
\item \path{StandardCompliancePersistenceCertificate.of_stage_flip}, \path{standardCompliance_three_goals_from_stage_flip_generalF_of_preBasinAscent}, \path{standardCompliance_three_goals_from_stage_flip_generalF_min_drift_shock_of_preBasinAscent} (\path{Theory/Basin/StageFlipPersistence.lean}) --- \Cref{prop:flip-route} as stated (round BH-1, 2026-10-06): the persistence certificate at the derived pivots, from the sender's stage complementarity on the floor range only, its top boundary sign, the diagonal residual, monotone sender values in the invariant set and interval dominance on the derived basin; the three goals follow with the myopic ceiling and full-sanction ascent below the basin, general in $\F$. Against the route of the previous item it drops the target combined slope condition with the target's stage complementarity, top sign and invariant-set monotonicity, the bottom sign, single crossing, and ascent above the basin; the previous wrapper is now its corollary. The example economy of the derived-basin route inhabits it; no economy is exhibited that this route reaches and the previous one misses.
\item \path{eeCert_one_not_mem_targetFlipSet} (and the \texttt{qbCert}/\texttt{fgcCert}/\texttt{lbCert}/\texttt{b2Cert} siblings), \path{targetFlipSet_zero_mem_at_sbar_one}, \path{stageFlip_ordering_forces_sbar_lt_one}, \path{eeCert_stageFlip_MS_boundary_infeasible}, \path{eeCert_stageFlip_legacy_infeasible} --- The machine-checked non-vacuity refutation of prop:flip-route at the shipped economies (the ``Non-vacuity status of this route'' paragraph): at every example economy with a complete certificate the top boundary sign $1\in\mathcal F_T$ fails for every admissible floor $\bar s<1$ (all five are entrenchment-condition designs with strictly negative target myopic sanction slope), the degenerate floor $\bar s=1$ contradicts the route's remaining clauses, and the stage increasing-differences primitives themselves hold (\path{eeCertParams_targetStageID} et al.) --- the failure is the boundary sign, not the complementarity structure. The gap is closed by the derived-basin-route example economy in the next entry.
\item \path{sign_conditions_insufficient_counterexample} (\path{Examples/SignConditionsInsufficient.lean}) --- the sentence after \Cref{ass:flip}: an economy of the class satisfying every sign condition the paper states (strict sender and target stage complementarity and their derivative-free forms, the substitution between compliance and the sanction, strict adjustment) that possesses an equilibrium and in which no equilibrium sustains sanctions above any positive floor on any non-degenerate basin; sign conditions order difference quotients and bound no magnitude, and the magnitude condition persistence needs is the sub-pivot slope condition (P5). Since round BH-5 (item BH5.3) the probe holds over any base economy, keeping its transition and discount factors (\path{probeOver}, \path{sign_conditions_insufficient_generalF}): no condition on the transition, added to the sign conditions, delivers perpetual sanctions.
\item \texttt{sfParams}, \texttt{sfCertificate}, \path{sfCert_three_goals_generalF} (headline status retired in round V, 2026-09-15; the theorem stays), \path{sfParams_one_mem_targetFlipSet}, \path{sfParams_zero_not_mem_targetFlipSet} --- The derived-basin-route example economy inhabiting prop:flip-route (the ``Non-vacuity status of this route'' paragraph): target sanction slope $5x-a-1$ flips at the derived pivot $2/5$ (top boundary sign holds with margin $4-a\ge3$; bottom sign strictly fails at $x=0$ --- the exact signs every entrenchment-condition example economy provably lacks); kinked sender payoff with state-riding stage argmax $\sigM(x,a)=\mathrm{clip}(x)$; singleton exact-fixed-point invariant sets; the certificate produced BY \path{StandardComplianceLocalCoreCertificate_MS.of_stage_flip} with derived pivots and basin $[1/2,1]$, and the three goals via the general-transition wrapper (plus the transition-independent dominance corollary \path{sfCert_three_goals_dominance}, itself routed through the general-transition headline). Since round BH-5 (item BH5.2) the witness is a family over the transition (\path{sfParamsF}, \path{sfFCertificate}, \path{sfFCert_three_goals_generalF}): its conclusions hold at every transition satisfying (F1)--(F4) with the corner section $\F(y,1)=y/2+1/2$ (\path{CornerSectionTop}) and a lower bound of at least $11/36$ on the sanction slope, which keeps the sender at full sanction; \texttt{sfParams} is its affine member by definition, and \path{sfqbParams} with \path{sfqbCert_three_goals_generalF} (\path{sfqbParams_F_ne_affineF}, \path{sfqb_lowerChord}) is the quadratic-bump member, lower bound $47/100$ on the sanction slope, the route's first non-affine inhabitant. The corner section is an equality on the line $s=1$, so the realized full-sanction path of every member is the affine one; existence only (Rule~9).
\item \path{TransitionStateLipschitzWith},
\path{V_S_valueLipschitzWith_of_transitionStateLipschitz}, and
\path{transitionStateLipschitzWith_of_affine} --- The optional quantitative general-$F$ state-regularity route (F6-Q): a fixed-sanction state modulus $\gamma_F$ and $\beta_S\gamma_F<1$ give the constructed sender value the closed-form modulus $K_S=L_{S,x}^{\alpha}/(1-\beta_S\gamma_F)$; the affine transition has the exact specialization $\gamma_F=1-\delta$.
\item
\path{senderDynObjective_strictMono_of_transition_adjusted_chord_and_V_lipschitz},
\path{stageSenderObjective_strictMono_of_transition_adjusted_chord},
\path{senderDynObjective_strictMono_of_transition_adjusted_chord_and_stateLipschitz},
and
\path{StandardComplianceLocalCoreCertificate_MS.of_transition_adjusted_chord}
--- The exact general-$F$ cardinal route of prop:F-chord-route: the stage gain is compared with $\beta_S K_S$ times the actual distance between the two successor states, so no sanction-coordinate Lipschitz modulus for $F$ is assumed; the transition-adjusted slope condition supplies strict stage monotonicity and dynamic-objective monotonicity, and packages these into the slim MS certificate. The constructor still takes the independently load-bearing fields for monotonicity on the sender invariant set and for the target combined slope condition.
\item \path{senderDynObjective_strictMono_of_stage_lowerSlope_and_V_lipschitz}, \path{sender_high_argmax_exists_of_strictMono_below_threshold}, \path{hSC_on_of_strong_rally_primitives} --- The affine specialization of the sender cardinal route: under the strict-slope margin beta\_S K\_S delta \textless{} m\_S the dynamic sender objective strictly increases in s below the pivot for every K\_S-Lipschitz continuation, locating the upper argmax at or above max(x, sbar); the Milgrom-Shannon single-crossing clause on the basin is derived from the same foresight premium primitives under a positive stage slope, and at the ordinal generality of the invariant set it is carried by the assumed certificate field (rem:ido).
\item \path{SCLowerOn}, \path{targetFF_strictSCLower_targetMM_of_switchWeak_targetStrict},
\path{alpha_order_dynamic_le_myopic_of_switchWeak_targetStrict} --- the mixed
composition after \Cref{prop:ordinal-ordering}, in
\path{Sanctions/Theory/LockIn/TwoSidedForwardLooking.lean}.  The weak relation
carries two clauses and deliberately has \emph{no} arbitrary-pair ordering
theorem: \path{argmax_meet_join_of_sCLowerOn} orders the correspondences in the
strong set order, and \path{greatest_argmax_le_of_sCLowerOn} their extremal
elements, which is exactly the caution the text states.
\item \path{sigma_dyn_top_ge_sigma_of_senderContinuationMonotoneAt},
\path{sigma_dyn_top_ge_sigma_of_responsePairContinuationOrder},
\path{SenderContinuationMonotoneAt.of_valueMonotone},
\path{SenderContinuationOrderAtResponses.of_localMonotone} --- the localization
of the step for the direct effect of the sender's foresight, same file.  The three conditions form a ladder
from global monotonicity of $\Vs$ to a condition on the two responses alone,
each implication proved, so the paper's weaker statement costs nothing.
\item \path{IncreasingChordsOn}, \path{DecreasingChordsOn},
\path{argmax_shift_ge_of_increasingChords_of_upperGreatest},
\path{argmax_shift_le_of_decreasingChords_of_lowerGreatest},
\path{argmax_shift_eq_of_chords_of_greatest} --- the chord mechanism of
\Cref{app:switch-route}, in \path{Sanctions/Foundations/ArgmaxSensitivity.lean}.
Stated over an abstract loss, perturbation and action set, with no reference to
the sanctions model.  \path{ConvexOn.increasingChordsOn} and
\path{ConcaveOn.decreasingChordsOn} record curvature as a discharger rather than
as the hypothesis.
\item \path{RecentredSCLowerAt}, \path{responseGap_le_of_recentredSCLower},
\path{recentredSCLowerAt_of_increasingChords} --- the ordinal weakening in
\path{Sanctions/Theory/LockIn/TwoSidedForwardLooking.lean}.  The third
declaration is the ladder rung establishing that the ordinal condition is
implied by the chord one, which is what makes it a weakening rather than a
different assumption.  The translated-loss hypothesis these consume is
\path{TranslatedLossSenderDifferencesAt}, the payoff-\emph{difference} form
stated in the text, which leaves an arbitrary compliance-dependent level free;
\path{TranslatedLossSenderAt.toDifferences} converts from the exact-level form,
so a model supplying the stronger identity loses nothing.
\item \path{senderSwitchWeakSCLower_of_convexContinuation_and_targetShape} ---
the assembled route to the weak switch leg, same file, carrying all five of its
inputs visibly: the gap condition and the four shape conditions on the target's
complete sanction payoff, the latter confined to feasible sanctions by
\path{switchWedge_antitoneOn_of_gap_and_shapeOn}.
\item \path{BaselineTwoSidedOrderingUndetermined}, in particular
\path{nt_ordering_reversed}, \path{baseline_does_not_determine_two_sided_order},
\path{nt_wedge_condition_fails} and
\path{nt_gap_monotoneOn_via_general_mechanism} --- the reversal economy cited in
\Cref{app:switch-route}.  The last of these discharges the four hypotheses of
the adverse branch of the mechanism at that economy, so the abstract route is
visibly instantiated rather than merely asserted to cover it.  Its sender value
is linear, so the composite continuation collapses to $\F(x,\cdot)$ and the
adverse shape is the transition's alone.
\item \path{NetChannelDominanceAt},
\path{netChannelDominanceAt_iff_strictDecreasingDifferences},
\path{alpha_order_dynamic_le_myopic_of_netChannelDominance},
\path{netChannelDominanceAt_of_switchAntitoneOn_targetStrictAntiOn} ---
\eqref{eq:net-channel} and its two boundaries, in
\path{Sanctions/Theory/LockIn/TwoSidedForwardLooking.lean}.  The second
declaration is the equivalence with the strictly-falling total gap, which is
what makes the condition an accounting restatement rather than a new assumption;
the fourth is the strictly stronger channel-by-channel route, recorded as a
corollary so the direction of the implication is visible.
\item \path{directStrictSC_of_netChannelDominance} in the same file, with
\path{strictSCLowerOn_cube_id},
\path{not_strictAntiOn_cube_sub_id} and
\path{exists_strictSCLowerOn_without_gap_strictAnti} in
\path{Sanctions/Foundations/BinaryComparativeStatics.lean} --- the two halves of
the first boundary in \Cref{app:two-channel}.  The first names the arrow from
net dominance to the ordinal relation \eqref{eq:sc-lower}, which was already in
use inside the proof of the $\alpha$-ordering but had no name of its own; the
rest refute the converse on $[0,1]$ at $f(a)=a^{3}$, $g(a)=a$.  Together they
fix the direction: the cardinal condition is strictly stronger, not
incomparable.
\item \path{CrossOrderOn},
\path{crossOrderOn_iff_le_monotoneFloorOn} in
\path{Sanctions/Foundations/CrossOrderComparison.lean}, with
\path{CrossClosedLoopStepOrderOn},
\path{rulePathF_le_on_stopped_of_crossStepOrder},
\path{crossClosedLoopStepOrderOn_of_rule_le_and_upper_mono} and
\path{ruleSanctionsF_le_on_stopped_of_crossOrders} in
\path{Sanctions/Theory/LockIn/DifferentPolicyComparator.lean} ---
\Cref{prop:cross-step-paths}, the exact
hypothesis behind \Cref{app:closed-loop-monotone}.  Ordering the two state paths
needs neither rule to be monotone: what the induction consumes is that the lower
closed loop at a state of the region never exceeds the upper closed loop at any
feasible state weakly above it, which is domination by the greatest monotone
minorant of the upper loop.  Monotonicity of the upper rule is one discharger of
that, and the earlier theorem is re-proved through the general one with its
statement unchanged.  Sanction-path order is stated separately rather than
inferred from state order, because the two rules are evaluated at different
states.
\item \path{ruleStepF}, \path{rulePathF}, \path{ruleSanctionsF},
\path{rulePathF_mem_I}, \path{rulePathF_le_of_rule_le_and_closedLoop_mono},
\path{rulePathF_le_of_rule_le_and_upper_mono},
\path{rulePathF_le_of_rule_le_and_upper_monoOn_basin} and \path{affineCopy} in
\path{Sanctions/Theory/LockIn/RulePathF.lean} --- the general-transition rule
path and its order facts, moved unchanged from
\path{DifferentPolicyComparator.lean} in round BH-6 (item BH6.9), and the affine
copy of an economy, the same economy with its transition replaced by the affine
law at its own rate $\delta$. The affine rule-path lemmas
\path{rulePath_le_of_rule_le_and_upper_mono},
\path{rulePath_le_of_rule_le_and_upper_monoOn_basin} and
\path{lockInPath_mem_I_of_mem_I} are one-line instances at that copy
(\path{rulePath_eq_affineCopy}, \path{lockInPath_eq_affineCopy}); no statement
changed and no transition hypothesis was added.
\item \path{exists_crossStepOrder_without_monotone_rule} in
\path{Sanctions/Examples/CrossStepWithoutMonotoneRule.lean} --- the Rule 9B
witness that the weakening is real.  Under the lattice-meet transition
$\F(x,s)=\min\{x,s\}$, the rule $R(x)=\max\{x,1-x\}$ is feasible and not
monotone, yet lies weakly above the state everywhere, so the closed loop is the
identity and the cross-order condition holds while the monotonicity it replaces
fails.  Rule monotonicity is a property of the rule; what the comparison needs is
a property of the rule composed with the transition, and the two come apart
wherever the transition is insensitive to the sanction.
\item \path{targetFFObjectiveAtValues},
\path{DirectTargetStrictSCLowerAtValues} and
\path{alpha_order_dynamic_le_myopic_of_directStrictSC} in
\path{Sanctions/Theory/LockIn/TwoSidedForwardLooking.lean} --- the ordinal
comparison stated at a bare value pair rather than at an equilibrium, so that
invariant-set-level and fixed-point-level results can impose it before an equilibrium
exists.  The equilibrium forms are definitional instances, so this re-indexes the
assumption surface without adding to it.
\item \path{greatestMyopicAlpha},
\path{myopicFeedbackEquilibrium_greatest} and
\path{myopic_alpha_le_greatest} in
\path{Sanctions/Theory/LockIn/DifferentPolicyComparator.lean} --- the greatest
myopic target policy, mirroring the sender-side \path{sigma_dyn_top}: the target's
stage argmax set is non-empty, bounded and closed, so it contains its supremum.
Every myopic equilibrium's policy lies weakly below it, which is what makes it the
right object for a weak comparison.  No continuity is claimed for it and the
equilibrium structure asks for none, so it must not be fed into a forward-looking
Bellman construction without a separate argument.
\item \path{generatedChord_strictly_weaker} in
\path{Sanctions/Examples/GeneratedChordStrictlyWeaker.lean} --- the Rule 9B
witness for the generated-chord rung of \Cref{app:switch-route}.  Substituting
$u=\sigM(x,a_H)+z$, the generated condition says that the $\Delta$-difference of
the composite continuation is non-decreasing, where $\Delta$ is the sanction gap:
one lag, against every lag for the all-chord form.  A continuation term periodic
with period $\Delta$ therefore satisfies it with vanishing difference while being
far from convex, and the witness is the $1/2$-periodic
$\varphi(s)=\bigl|\,|s-1/2|-1/4\bigr|$ against a clamped sender response whose gap
at $a_L=0$, $a_H=1$ is exactly $1/2$.
\item \path{alpha_order_at_equilibrium}, \path{ruleDom_at_equilibrium},
\path{SchurValueSystem}, \path{SchurValueSystem.equilibrium},
\path{SchurValueSystem.ruleOrdered_at_constructed} and
\path{exists_ruleOrdered_equilibria_schur} in
\path{Sanctions/Theory/Certificates/TwoSidedRuleComparison.lean} --- the three
levels of \Cref{app:jointiteration}, in order: conditions at a supplied
equilibrium, at the constructed fixed point, and over the whole invariant set.
\path{SchurValueSystem.equilibrium} is a \path{def} choosing from an existence
theorem proved from primitives by the in-repo Perov contraction, so the fixed
point is constructed rather than assumed.  The invariant-set-level theorem retains its
statement and is re-proved by discharging the point-specific conditions at the
constructed equilibrium, which is what exhibits the ordering between levels.
\item \path{sCLowerOn_iff_pairwise_meet_join} --- the converse to
\path{argmax_meet_join_of_sCLowerOn} in
\path{Sanctions/Foundations/BinaryComparativeStatics.lean}: on a
one-dimensional choice set the weak relation holds \emph{if and only if} the
argmaxes are strong-set ordered on every feasible pair, so it is what the
strong-set conclusion is rather than one sufficient condition among others.
Alongside it, \path{alpha_leastFF_le_anyMM_of_twoWeakChannels} and
\path{alpha_anyFF_le_greatestMM_of_twoWeakChannels} record what an extremal
selection buys when neither channel can be signed strictly.
\item \path{SemiconvexOnSet} (the mirror of \path{SemiconcaveOnSet}, the two
forming one derivative-free-form pair in \path{Sanctions/Foundations/SemiconcaveSup.lean}),
\path{not_semiconvexOnSet_of_linearDefect},
\path{not_semiconcaveOnSet_of_linearDefect} in
\path{Sanctions/Foundations/SemiconvexityDefect.lean}, with
\path{minMeetTransition_closedLoop_not_semiconvex} and
\path{maxJoinTransition_closedLoop_not_semiconcave}
(\path{Sanctions/Examples/TransitionShapeNotPreserved.lean}) and
\path{clampedPeakSender_closedLoop_not_semiconvex} together with its
\path{..._not_semiconcave} twin
(\path{Sanctions/Examples/SelectorShapeNotPreserved.lean}) --- the two
independent obstructions of \Cref{app:existencemachinery}.  The Foundations
lemmas are stated over an abstract function and refute $\exists c$, so they say
no finite modulus works rather than that some particular one fails.
\item \path{selfMatchingMyopic_not_antitoneOn},
\path{selfMatchingMyopic_strictMonoOn},
\path{exists_myopic_policy_not_antitone} in
\path{Sanctions/Examples/MyopicPolicyNotAntitone.lean} --- the refutation behind
\Cref{app:closed-loop-monotone}.  \path{rulePathF_le_of_rule_le_on_stopped} and
its wiring \path{rulePathF_myopic_le_dynamicTop_stopped} are the positive
counterpart: path comparison at every date through which the lower path has
in fact stayed in the region, with no invariance assumed and the region required
to be neither closed nor an interval.
\end{itemize}

\paragraph{Existence and the Schur condition (\Cref{sec:existence}).}
\begin{itemize}[leftmargin=1.4em,itemsep=1pt]
\item \path{feedbackStackelbergTop_corner_of_substituted_stage_le}, \path{V_T_dyn_on}, \path{targetDynBCF_congr_on}, \texttt{propext}, \texttt{Classical.choice}, \texttt{Quot.sound} --- The Lean theorem verifying the Basin-corner feedback-Stackelberg existence result (thm:corner-existence) in its weak form (round BH-2): continuity of the response on feasible states only, comparison with full non-compliance only, and a weak decrease of the substituted stage payoff; the strict form \path{feedbackStackelbergTop_corner_of_substituted_stage_strictDecr} is its corollary. Axiom-clean.
\item \path{b2_leaderArgmax_not_antitone_in_V_T} --- The Lean-checked order-reversal refuting antitonicity of the leader argmax in V\_T (Knaster-Tarski route obstruction, sec:existence:nonvacuity).
\item \texttt{cycleParams}, \path{jointIteration_two_cycle}, \texttt{jointIteration\_cycle\_\{T01,S01,T10,S10\}}, \texttt{cycle\_pair\_ne}, \path{cycle_pair0_not_fixed}, \path{cycle_pair1_not_fixed} --- The kernel-checked exact rational period-two orbit of the joint Bellman iteration (the ``machine-checked cycling'' of sec:existence:nonvacuity and the Schur-condition discussion): at an economy with strictly concave sender payoff ($\delta=\beta_T=\beta_S=\tfrac12$), the iteration exchanges two distinct affine value pairs, every argmax along the orbit is a pinned singleton (both selection conventions agree, with exact rational corner margins $\tfrac1{20}$, $\tfrac1{40}$), the leader flipping $a{=}1\leftrightarrow a{=}0$ through the D3b order-reversal mechanism; upgrades the Doraszelski--Satterthwaite \S3 numerical cycling to an exact certificate. Cluster index: \path{Examples/DeterministicInteriorObstruction.lean}.
\item \texttt{eatonEngersParams}, \texttt{F\_mono\_s}, \path{Polity.inducedF_mono_s}, \path{eatonEngersPolity_lockIn} --- Lean verification that the entrenchment sign (F3) of the Eaton-Engers example economy, i.e. F\_mono\_s, is derived from the polity's lock-in mechanism.
\item \texttt{eeCertCertificate} --- The Lean certificate object for the certificate-complete Eaton-Engers example economy, which discharges every field of the local invariant-set certificate.
\item \path{eeCert_three_goals_dominance} --- The Lean theorem that the three goals (persistence/dominance/shock) hold unconditionally at the eeCert example economy.
\item \path{fixedActionIDO_not_preserved_by_continuous_feasible_policy} --- The exact counterexample in \Cref{rem:policy-substitution-obstruction}: fixed-action interval dominance, continuity, and feasibility do not imply interval dominance after substituting the realized target policy.
\item \path{substParams_DAU_counterexample}, \path{substParams_certificate_senderSide_infeasible} --- The two Lean impossibility artifacts showing the compliance--sanction cross term is necessary (forcing the Eaton-Engers sign d\textasciicircum{}2U\_S/ds da \textless{} 0).
\item \path{mechanismWitnessParams}, \path{mechanismWitnessParams_three_goals_conditional}, \texttt{propext}, \texttt{Classical.choice}, \texttt{Quot.sound} --- The Lean objects for the P3 mechanism-active example economy and its conditional three-goals result, plus the axiom-cleanliness aside.
\item \texttt{lowDeltaParams}, \texttt{lowDelta\_window\_C1}--\texttt{C4}, \path{lowDelta_chain_Llead}, \path{lowDelta_contraction_margin} --- The low-$\delta$ window diagnostic (Table tab:witness-roster): the $P_3$ stage payoffs at $\delta=1/100$ with purified kernel half-width $h=1/10$, at which all four double-box kernel-width inequalities hold as exact rational facts (contraction constant $L_{\mathrm{lead}}\approx0.144<1/2$, robust to the chain-form ambiguity). The instantiated shocked economy still uses the one-box moving average. The exact scalar arithmetic is optional checking at the example economy: the conditional existence wrapper that consumed the leader selector and objective-geometry package was retired on 2026-07-27 to commented audit evidence, and no low-$\delta$ existence declaration remains. The double-box arithmetic is not a proof of the one-box mixed objective-difference estimate.
\item \path{movAvgReal_deriv_lipschitz_of_valueLipschitz}, \path{movAvgReal_strongConcaveOn_of_valueLipschitz}, \path{coupled_sender_strongConcave}, \path{coupled_sigmaDynTop_continuous_subtype}, \path{coupled_sigmaDynTop_clamped_continuous} --- The repaired sender regularity row at the low-$\delta$ example economy: a one-box moving average of a Lipschitz continuation has a Lipschitz derivative and the corresponding semiconcavity bound; combined with the primitive sender curvature this proves modulus \(37877/32000\) on the whole sender invariant set. Uniqueness plus Berge gives continuity of the physical selector and its clamped extension. These facts are proved rather than consumed by the since-retired coupled templates.
\item \path{argmax_lipschitz_two_coordinates_of_quadraticGrowth_of_coordinate_geometry}, \path{argmax_lipschitz_two_coordinates_of_quadraticGrowth}, \texttt{coupled\_L\_TS}, \texttt{coupled\_L\_ST}, \path{coupledSchur_margin_robust}, \path{coupled_rowSum_fails}, \path{coupled_lead_lipschitz}, \path{coupled_isMaxOn_of_quadraticGrowth} --- The canonical Foundations theorem starts from quadratic growth with modulus \(q>0\) and separate objective-difference bounds \(L_x\,d(x,x')|a-a'|\) and \(L_y\,d(y,y')|a-a'|\), then derives the selector bound \((L_x/(2q))d(x,x')+(L_y/(2q))d(y,y')\). Thus the selector moduli come from objective geometry rather than from numerical policy hypotheses. The coupled low-$\delta$ declarations are only exact-rational diagnostics. Both cross constants are positive, the Schur inequality holds, and the target row-sum fails, but the current value set does not inhabit the required realized-choice objective-geometry premise.
\item \path{not_coupledLeaderStrongConcave} (headline status retired in round V, 2026-09-15; the theorem stays), \path{coupledLeaderObj_zero_zero_quadratic_growth}, \path{coupledGrowthRamp_mem}, \path{coupledGrowthRamp_obj_plateau}, \path{not_coupledTargetQuadraticGrowth}, \path{coupledZeroPair_target_bellman_zero}, \path{coupledZeroPair_target_bellman_one_tenth}, \path{coupledZeroPair_target_image_not_mem}, \path{coupledSteepPolicy_sender_image_not_mem} --- The disprove-first coupled frontier. The zero-pair slice has a unique maximizer at \(35/1392\) and a valid local modulus-one loss bound, but modulus-\(2\) curvature is false. More strongly, a monotone target continuation in the value set creates a global-maximizer plateau and refutes quadratic growth uniform over the value set. At the endogenous zero pair the actual target Bellman image falls from \(0\) at \(x=0\) to \(-7/600\) at \(x=1/10\), directly refuting invariance of the present monotone target value set. Continuity and feasibility alone likewise do not preserve the sender Lipschitz value set.
\item \path{affine_mixed_rectangle_of_deriv_lipschitz}, \path{targetDynBCF_norm_le_of_cube}, \path{targetDynBCF_valueLipschitz_of_integrand}, \path{senderPrimitiveValLipCarrier}, \path{targetPrimitiveValLipCarrier}, \path{movAvg_bellmanS_shock_mem_senderPrimitiveValLipCarrier}, \path{movAvg_targetDynBCF_shock_mem_targetPrimitiveValLipCarrier_of_integrand} --- The generic toolkit for alternative value sets: a four-point mixed estimate for affine shifts of a function with Lipschitz derivative and primitive-derived norm--Lipschitz value sets. Their norm radii are discounted stage-payoff bounds; the sender state radius is the Bellman closed form from the composed stage and transition moduli, while the target state radius is the derived modulus of the realized target integrand.
\item \path{lowDeltaTargetCarrierNorm}, \path{lowDeltaSenderCarrierNorm}, \path{lowDeltaSenderCarrierLip}, \path{coupledCarrierTLip}, \path{coupledSenderPrimitiveCarrier}, \path{coupled_sigmaDynTop_state_lipschitz}, \path{coupled_target_Bellman_mem_lip_carrier}, \path{coupled_sender_Bellman_mem_primitive_carrier}, \path{coupled_sender_Bellman_mem_carrier_of_policy_lipschitz} --- The low-\(\delta\) redesign of the value sets for invariance in \path{Sanctions/Examples/CoupledRoute/CoupledCarrierRedesign.lean}. The moduli of the value sets are now defined by primitive closed forms: stage cube bounds and discounts determine norm radii, while payoff slopes, curvature, the moving-average width, transition geometry, and a symbolic realized-policy modulus determine state radii. The familiar rationals \(20\), \(523/160\), \(57377/3200\), \(14320323/3787700\), and \(25683423000/283976393\) survive only as exact equalities at the example economy. The actual greatest sender response proves target-row invariance. For every nonnegative symbolic policy modulus \(L_\alpha\), the sender Bellman row preserves its \(L_\alpha\)-indexed primitive invariant set whenever the realized policy has that modulus. The old numerical-budget theorem for a fixed value set is retained only as a diagnostic; the formalization does not assert that the endogenous target selector satisfies the symbolic policy premise.
\item \path{LowDeltaCarrierLocalEquilibriumFrontier.lean}, \path{lowDeltaCarrierLocalSenderPolicySlack_eq} --- The former six-premise existence wrapper local to the value sets has been retired. The file retains only local-set wiring and exact residual/slack diagnostics; none is promoted as an equilibrium-existence theorem or as an additional model assumption.
\item \path{coupledRedesignTarget_two_global_maxima}, \path{not_coupledTargetQuadraticGrowthLipCarrier}, \path{coupledRedesignTarget_nearby_unique_maxima}, \path{coupledRedesign_every_optimal_selector_violates_budget}, \path{not_coupledRedesign_optimal_selector_state_lipschitz_budget}, \path{coupledRedesign_reduced_branches_exact_quadratic_growth} --- The exact obstruction for the redesigned value sets in \path{Sanctions/Examples/CoupledRoute/CoupledCarrierRedesignCounterexamples.lean}. One continuation pair from the value sets produces two global target maxima at \(x=5/12\), refuting modulus-\(1/2\) quadratic growth uniform over the value sets even though each smooth response branch retains exact positive quadratic growth. At the two states \(5/12\pm10^{-8}\), the unique optimal actions force every optimal selector for this pair to have secant modulus about \(94.04\), above the sender-invariance budget \(4675077/640000\approx7.3048\). The result refutes this exact discharge uniform over the value sets; it does not prove equilibrium nonexistence, selector discontinuity, or failure of every finite Lipschitz constant.
\item \path{coupled_sigmaDynTop_action_lipschitz}, \path{coupled_sigmaDynTop_value_lipschitz}, \path{coupledLeaderObj_diff_lip_in_VT}, \path{coupledLeaderObj_sender_advertised_bound_fails}, \path{not_coupledLeaderObj_sender_uniform_advertised}, \path{coupledLeaderObj_pair_pointwise_bound_on_lip_carrier}, \path{coupled_lead_sq_holder_of_quadraticGrowth_on_lip_carrier} --- The coupled objective-difference audit in \path{Sanctions/Examples/CoupledRoute/CoupledObjectiveDifference.lean}. The realized sender response is action-Lipschitz with constant \(32000/37877\) and continuation-value Lipschitz with constant \(3200/37877\). An explicit value pair refutes the advertised sender-value coefficient as a uniform all-actions estimate, but the redesigned value set supports the pointwise two-value bound with coefficients \(9/10\) and \(43520000/283976393\). Conditional on modulus-\(1/2\) growth at both realized choices, this yields the squared H\"older bound with coefficients \(9/5\) and \(87040000/283976393\). It does not discharge the literal product-form realized-choice difference bound required by the Perov capstone, which remains open.
\item \path{lowDelta_realized_IDO_slope_threshold_eq}, \path{lowDelta_realized_IDO_slope_threshold_pos}, \path{lowDelta_senderDynObjective_realized_IDOFamilyOn_of_policy_slope} --- The realized-policy interval-dominance implication in \path{Sanctions/Examples/CoupledRoute/LowDeltaRealizedIDO.lean}. For every sender continuation in the approved value set, the full-state sender objective satisfies interval dominance if the target policy obeys the one-sided slope bound \(\alpha(y)-\alpha(x)\le(2959677/3200000)(y-x)\) for \(x\le y\). The exact threshold is positive. Its discharge for the endogenous equilibrium policy remains open, as does an unconditional explicit low-\(\delta\) coupled-equilibrium witness.
\item \texttt{eeVTstarC1}, \path{eeCert_shadow_price_target/sender}, \path{fgc_shadow_price_target/sender}, \path{*_envelope_recursion_*}, \path{eeCert/fgc_target_hardline_gain}, \path{b2_interior_envelope_identity}, \path{eeCert/fgc_three_goals_local_C1} --- The C\textsuperscript{1}/envelope layer at the example economies (the shadow-price paragraph of sec:compstat): first inhabitants of the C\textsuperscript{1} value set at the pinned exact fixed points; the exact shadow prices $d\Vt^{\ast}/dx$, $d\Vs^{\ast}/dx$ with the Rinc\'on-Zapatero--Santos boundary-envelope recursion verified as algebra and on the model objects; Rule-16 sign corollaries; a machine-checked interior envelope (Danskin) identity at the non-monotone example economy; and the C\textsuperscript{1} three-goals endpoint discharged at both example economies. The abstract envelope layer stays frozen (would require differentiability primitives).
\item \emph{Family level, and this is the primary backing:} on the general-transition family since round BJ (item BJ4.6), \path{SymbolicEconomyF.equilibrium_feedback_decomposition}, \path{lockinValue_eq}, \path{lockin_amplification}, \path{lockin_strict_iff}, \path{lockin_strict}, \texttt{VT\_$\beta$S\_invariant}, at any equilibria whose sender values lie in the certified invariant sets (no uniqueness within the invariant sets is used; \path{corner_bellman_unique} gives the target's invariance for members that share $\F$, $\beta_T$, $-c$ and $n+v$, whatever $\ell$); on the affine family, for its closed-form values, \path{symbolic_equilibrium_feedback_decomposition}, \path{symbolic_lockin_amplification}, \path{symbolic_lockin_strict_iff}, \path{symbolic_lockin_strict}, \texttt{symbolic\_VTstar\_\{beta\_S\}\_invariant} --- the equilibrium-feedback comparative statics (\Cref{prop:equilibrium-patience-feedback}) quantified over \emph{every} member of the symbolic certified family that satisfies its defining inequalities, not over a calibration. \emph{General-$\F$ level:} \path{senderValue_patience_decomposition_of_bellman}, \path{senderValue_lockin_amplification_of_bellman}, \path{senderValue_lockin_capitalised_of_bellman}, \path{senderValue_lockin_strict_iff_of_bellman} --- the same two identities and the sign equivalence with \emph{no functional form for the transition}, the absorbing state derived from \texttt{F\_diag} rather than assumed; the underlying algebra is \path{PatienceDecomposition.difference_decomposition} and its siblings in \path{Foundations/PatienceDecomposition.lean}. Since round BH-6 (item BH6.3) the lock-in identities hold for either government's stage payoff (\path{cornerValue_lockin_amplification_of_bellman}, \path{cornerValue_lockin_capitalised_of_bellman}, \path{cornerValue_lockin_strict_iff_of_bellman}), and the target trio \path{targetValue_lockin_amplification_of_bellman}, \path{targetValue_lockin_capitalised_of_bellman}, \path{targetValue_lockin_strict_iff_of_bellman} says that, for every transition, the target's lock-in value is its corner flow capitalized, $\Ut(1,a,1)/(1-\beta_T)$, and responds to its own patience through that flow alone; the $\delta$-invariance lemmas of the example families are coefficient algebra consistent with it. \emph{Numeric instance:} \path{eeCertBEquilibrium}, \path{eeCertB_equilibrium_feedback_decomposition}, \path{eeCertB_equilibrium_lockin_strict}, \path{eeCertB_lockin_amplification}, \path{eeCertD_lockin_value_invariant}, \path{fgcD_VS/VT_lockin_invariant} --- every parameter in the exact sender-patience interval packaged as an equilibrium in $\mathcal E$ rather than a pinned candidate value; the two identities there are now one-line corollaries of the general algebra rather than separate closed-form computations. Ordering matters in this entry: a 2026-08-16 audit read the numeric names as the only backing and wrongly reported the proposition as a witness-backed universal.
\item \path{epsIsolation_of_schur_rows}, \path{epsilonIsolation_feedbackStackelbergTop_schur}, \path{coupledEpsMargin_frontier_bracket}, \path{eeCert_epsilon_isolation}, \path{eeCert_global_isolation} --- The $\varepsilon$-isolation layer (the local-isolation passage, moved in round T to the notes' moved-exposition section): the parametric neighborhood-isolation theorem (sup-norm closeness + shape pins $\Rightarrow$ coincidence with the fixed point in the invariant set, via the Schur condition on the enlarged invariant set), and its instantiation at the Eaton--Engers example economy, together with the global form there: the non-compliance-policy convention alone pins both values through the raw discount contractions. At the coupled low-\(\delta\) economy what is machine-checked is the Schur slack on the enlarged value set itself --- positive at \(\varepsilon=24\), negative at \(\varepsilon=25\) --- and not an instantiated isolation conclusion; the conditional wrapper that would have drawn one was withdrawn on 2026-07-28 because its modulus-2 leader premise is refuted on those value sets (\path{not_coupledLeaderStrongConcave}, \path{coupledSchurWide_sc_T_contradictory}).
\item \texttt{NNMat2.specRadius}, \path{specRadius_lt_one_iff_schurDet}, \path{perov2_existsUnique_fixedPoint}, \path{exists_fixedPoint_in_carrier_schur}, \path{exists_feedbackStackelbergTop_from_valueClass2_with_membership_schur} --- The Schur-condition sufficiency chain (thm:existence): for non-negative coefficients with diagonal entries below one, rho(M) \textless{} 1 is machine-checked equivalent to the elementary inequality L\_TS L\_ST \textless{} (1-beta\_T)(1-beta\_S,row); the Perov weighted-max contraction yields the unique joint fixed point, and on the invariant set the fixed point delivers an equilibrium in $\mathcal E$; the row-Lipschitz estimates enter as hypotheses, with no sign condition on their four coefficients since round BJ (item BJ1.1: a negative coefficient forces a one-point invariant set, a case the Perov step handles itself, and both diagonal bounds are read), matching the theorem's explicit caveat that the Schur condition is an input: it is discharged at the analytic example economies, including the non-affine quadratic-bump economy of the next item, and remains open at the coupled low-$\delta$ economy, where the present value set does not inhabit the realized-choice objective-geometry premise the Perov capstone needs. Since round BJ (item BJ4.3) the value class asks for the sender's argmax to be unique only at feasible states and feasible compliance levels $a\in[0,1]$ (the field \texttt{sender\_unique\_on} of \path{ValueClass2}, \path{Params.DynamicArgmaxUniqueOnI}), which uniqueness at every real compliance level implies (\path{DynamicArgmaxUnique.toOnI}); the joint target operator keeps its values, its continuity now coming from the integrand with compliance projected to $[0,1]$ (\path{jointTarget_BCF_on_S}).
\item \path{existsUnique_fixedPoint_in_carrier_schur}, \path{fixedPoint_of_feedbackStackelbergTop_of_selector_coincidence}, \path{feedbackStackelbergTop_unique_in_carrier_schur} (headline status retired in round V, 2026-09-15; the theorem stays) --- The chain for uniqueness within the invariant sets (sec:existence): under the Schur condition the joint fixed point in the invariant sets is unique; every equilibrium with values in the invariant sets and selector coincidence is such a fixed point; hence any two such equilibria agree in both value functions and the policy --- the machine-checked form of ``uniqueness holds within the carrier value class,'' with global uniqueness neither claimed nor established. Since round BJ (item BJ1.1) the chain carries no sign binder, and \path{feedbackStackelbergTop_unique_in_carrier_schur} backs the coincidence clause of \Cref{thm:existence}.
\item \path{StandardComplianceLocalCoreCertificateSchur.ofRowSum}, \path{StandardComplianceLocalCoreCertificateSchur_MS}, \path{StandardComplianceLocalCoreCertificate_MS.unique_in_carrier}, \path{eeCert_unique_in_carrier}, \path{b2_unique_in_carrier}, \path{dbCert_unique_in_carrier}, \path{fourGroupCert_unique_in_carrier}, \path{qbCert_unique_in_carrier}, \path{sfCert_unique_in_carrier}, \path{lbCert_unique_in_carrier} --- Row-sum certificates map into the weaker Schur form, including the general-\(F\) Milgrom--Shannon package, and the conclusion of uniqueness within the invariant sets is instantiated with all numerical inequalities discharged at every example economy with a certificate, including the non-affine quad-bump economy.
\end{itemize}

\paragraph{Existence and the Schur condition (\Cref{sec:existence}), continued.}
\begin{itemize}[leftmargin=1.4em,itemsep=1pt]
\item \texttt{ShockedParams}, \path{bellmanS_shock_contracting}, \texttt{V\_S\_shock} --- The shocked sender Bellman operator remaining a beta\_S-contraction and the sender value being constructed by the Banach engine (shocked-existence Banach route; sibling of thm:corner-existence).
\item \texttt{HeSunShockParams}, \path{smearBCF_nonexpansive}, \path{heSun_shocked_corner_existence} --- The shocked corner equilibrium existing from the He-Sun atomless-kernel primitive (corner shocked existence).
\item \path{ShockedFeedbackStackelbergEquilibriumTop}, \path{V_T_dyn_shock_bellman_eq}, \path{shockedFeedbackStackelbergTop_of_constant_corner} --- The shocked equilibrium structure, both shocked Bellman fixed-point equations, and a non-vacuous shocked corner instance, with the two value functions constructed as Banach fixed points (sibling of thm:corner-existence).
\item \path{bellmanS_BCF_policy_pointwise_dist}, \path{shockedFeedbackStackelbergTop_interior} --- The interior shocked equilibrium via the slaved-operator contraction beta\_S + L\_\{U\_S,a\} L\_alpha \textless{} 1.
\item \path{argmax_lipschitz_of_strongConcave_lipschitzDiff} --- The general lemma: c-strongly-concave objective with action-Lipschitz maximizer-difference gives Lipschitz (not Holder-1/2) argmax bound c|a\_f-a\_g| \textless{}= 2L.
\item \texttt{window\_shift\_bound}, \texttt{movAvg\_lipschitz} --- The moving-average smear being state-Lipschitz with constant ||U||/h via the elementary window-shift bound (discharges smear\_lipschitz).
\item \path{lead_lipschitz_via_core}, \path{shockedFeedbackStackelbergTop_interior_discharged} --- The argmax-core reducing lead\_lip to a single action-direction estimate on the leader objective difference; interior existence discharged.
\item \path{sigma_top_lipschitz_in_V}, \path{VTdyn_shock_lipschitz} --- The shock's sigma- and V\_T\textasciicircum{}dyn-Lipschitz bounds that control the magnitude of the action-direction estimate.
\item \path{eatonEngersMovAvgShock} --- The concrete uniform-kernel (Eaton-Engers moving-average shock) inhabiting the shocked-existence class (non-vacuity).
\item \path{movAvgShockedParams_smear_lipschitz} --- The smear's state-Lipschitz constant being the closed form ||V||/h rather than a named primitive.
\item \path{mechanismWitnessParams_purified_existence_canonical_leader_symbolic}, \path{mechanismWitnessParams_purified_existence_canonical_leader_one_local_symbolic}, \path{mechanismWitnessParams_purified_existence_canonical_leader}, \path{mechanismWitnessParams_purified_existence_canonical_leader_one_local} --- The reusable conditional Banach construction of the purified P\_3 equilibrium with the leader derived from the genuine target objective. The canonical theorems are symbolic in any positive composite-curvature modulus \(c_T\) and nonnegative coupling modulus \(L_{\mathrm{lead}}\), under the direct condition \(\beta_S+2L_{\mathrm{lead}}/c_T<1\). Joint follower continuity plus composite strong concavity derive the leader map, feasibility, Berge continuity, and pointwise optimality; at full width, exact affine smoothing also discharges follower continuity. The fixed modulus-\(2\), \(L_{\mathrm{lead}}<1/2\) declarations are compatibility corollaries and are absent from the headline surface.
\item \path{mwStaticLeaderComposite_not_strongConcave_of_nonneg}, \path{mwStaticLeaderComposite_not_strongConcave_two} --- The exact clipped-response counterexample at state \(5/12\) and actions \(9/10,37/40,19/20\). The middle static composite payoff lies below the endpoint chord, so the primitive static target/sender composite is not even concave and hence cannot be strongly concave at any nonnegative modulus. This refutes deriving the dynamic composite-curvature premise merely from the target stage Hessian; it does not refute every continuation-adjusted shocked objective.
\item \path{mwEqQuartic_root_exists}, \texttt{mwEqSenderValue\_fix}, \path{mwEqLeaderPolicy_actual_isMaxOn}, \texttt{mwEqTargetValue\_fix}, \texttt{mwEqEquilibrium}, \path{mechanismWitnessParams_purified_existence_unconditional_one}, \path{mechanismWitnessParams_purified_existence_unconditional_one_monotone} (headline status retired in round V, 2026-09-15; the theorem stays) --- The separate exact full-width P\_3 construction. The endpoint-spread system is reduced to a quartic with an exact rationally bracketed root; the induced piecewise-affine leader policy, clipped greatest sender response, and both bounded-continuous values are constructed; global target optimality and both Bellman equations are verified through the clipping regions. The capstone has no hypotheses and packages an actual shocked equilibrium with monotonicity of its equilibrium-derived sanction rule on $[0,1]$.
\item \path{mwEqEquilibrium_sender_value_nonconstant},
\path{mwEqEquilibrium_target_value_nonconstant},
\path{mwEqEquilibrium_sender_continuation_strict_at_top},
\path{mwEqLeaderPolicy_endpoint_strict},
\path{mwEqDynamicTopRealizedRule_endpoint_strict},
\path{mechanismWitnessParams_fullWidth_state_varying_frontier},
\path{mechanismWitnessParams_fullWidth_mechanism_active_frontier} (headline status retired in round V, 2026-09-15; the theorem stays)
--- The full-width mechanism frontier. Both equilibrium values have a
strictly positive endpoint spread, and at state one the dynamic sender
optimum strictly exceeds the myopic best response evaluated at the same
equilibrium target action. The stronger state-varying capstone proves that
the target action rises strictly from state zero to state one and that the
realized sanction rises from zero to a strictly positive, strictly interior
top-state value. These are equilibrium state comparative statics. The
companion frontier combines the same activity facts with equilibrium-rule
monotonicity and the impossibility of any non-empty high-basin dynamic floor
at this specialization; neither declaration is a cross-economy
policy-feedback result.
\item \path{mechanismWitnessParams_purified_persistence} --- The conditional composition of purified P\_3 existence with the separate high-basin inputs for dynamic-over-myopic dominance: forward-looking realized sanction at least the myopic one on the basin (persistence, thm:persistence/thm:dominance). It is not an application to the explicit $h=1$ equilibrium, whose dynamic floor is impossible.
\item \path{shocked_dynamicTopSanctions_ge_myopic_on_basin} --- The shocked Goal-2 theorem once sanction-rule monotonicity, the dynamic floor, and the myopic basin/ceiling inputs are supplied.
\item \path{movAvgBCF_valueLipschitz} --- The moving-average smear preserving the value's Lipschitz constant K (same K), the key mechanism for the persistence lift.
\item \path{senderDynObjective_shock_strictMonoOn_of_transition_adjusted_chord},\\
\path{senderDynObjective_shock_singleCrossingOn_of_transition_adjusted_chord},\\
\path{movAvgShockedParams_senderDynObjective_shock_singleCrossingOn_of_transition_adjusted_chord} --- Consequence (i), exact general-$\F$ route: if each stage-payoff increase strictly exceeds $\beta_S K_S$ times the actual clipped transition displacement, the shocked sender objective is strictly increasing and therefore satisfies Milgrom-Shannon single crossing on the basin.  The moving-average witness discharges the smeared-value Lipschitz premise from the base value.  No affine pin or sanction-coordinate modulus for $\F$ is used.  The affine theorem \path{senderDynObjective_shock_singleCrossingOn_of_strong_rally} remains a sufficient specialization.  The independent transition-free cross-partial route is \path{senderDynObjective_shock_singleCrossingOn_of_cross_partial_generalF}.
\item \path{V_S_shock_valueLipschitzWith_of_contraction} --- The shocked sender Bellman fixed point being K\_S-Lipschitz with the same closed-form modulus (Lipschitz-cone preservation at the constructed equilibrium). Transition-general: the modulus is stated with the transition's state-Lipschitz constant, which is $1-\delta$ at the affine law.
\item \path{targetSanctionObjective_shock_monotoneOn_of_loss_offset} --- Consequence (ii): shocked target objective non-decreasing in the sanction on the basin from the entrenchment condition (Goal 1, thm:persistence; affine form plus the transition-free sibling under the combined slope condition,\\
\path{targetSanctionObjective_shock_monotoneOn_of_smear_mono_generalF}, using the box average's monotonicity preservation \texttt{movAvgBCF\_mono}).
\item \path{dynamicTopRealizedRule_shock_monotoneOn_of_transition_adjusted_chord},\\
\path{movAvgShockedParams_dynamicTopRealizedRule_shock_monotoneOn_of_transition_adjusted_chord},\\
\path{shockedFeedbackStackelbergTop_sanctionsF_ge_myopic_on_basin_of_IDO}
--- Consequence (iii), exact general-$\F$ route: the first two declarations
derive state monotonicity of the shocked sanction rule from the
transition-adjusted slope condition, including the box-kernel discharge stated on the
base value.  The third declaration instead substitutes the realized target
policy of a shocked feedback-Stackelberg equilibrium and derives the same rule
monotonicity from interval dominance.  With the separate dynamic-floor and
myopic-basin/ceiling inputs (no sign on the basin pivot since round BG-2), it proves that the
\path{ruleSanctionsF} sequence lies weakly above the independently optimized
myopic sequence in every period (Goal 2, thm:dominance).  The affine
compatibility route remains
\path{shocked_dynamicTopSanctions_ge_myopic_on_basin}, with
\path{dynamicTopRealizedRule_shock_monotoneOn_of_strong_rally} as its
discharger under a positive stage slope; the independent transition-free monotonicity sibling
is \path{dynamicTopRealizedRule_shock_monotoneOn_of_cross_partial_generalF}.
\item \path{mwCanonicalLeaderOne_adjusted_antitoneOn},\\
\path{mwSenderDynObjective_canonical_one_IDOFamilyOn},\\
\path{mwDynamicTopRealizedRule_canonical_one_monotoneOn_of_IDO} --- At the full-width P\_3 specialization, the adjusted canonical leader action $\alpha(x)-(27/10)x$ is antitone, the induced sender objective satisfies interval dominance on $[0,1]$, and the sanction rule is monotone there without assuming a positive stage slope.  This discharges the monotonicity component, not the separate dynamic floor required for the pathwise ordering.
\item \path{mwEqLeaderPolicy_monotone}, \path{mwEqSenderVertex_along_leader_monotone}, \path{mwEqSenderRule_along_leader_monotone}, \path{mwEqDynamicTopRealizedRule_monotoneOn} --- The direct equilibrium-level monotonicity proof for the explicit full-width witness: the optimal leader policy is increasing, the sender vertex remains increasing after endogenous compliance is substituted, clipping preserves the order, and the actual greatest-response rule of \texttt{mwEqEquilibrium} is therefore monotone on the physical state interval.
\item \path{mwSigmaDynTop_eq_clip_vertex_one},
\path{mwDynamicTopRealizedRule_canonical_one_top_eq_one_iff},\\
\path{mwCanonicalLeaderOne_endpoint_spread_of_floor},\\
\path{mwDynamicTopRealizedRule_canonical_one_no_uniform_floor_test} --- The exact full-width P\_3 sender response is the clipped quadratic vertex with endpoint correction $(V_S(1)-V_S(0))/24$.  Therefore every non-empty-basin floor forces $V_S(1)-V_S(0)\ge10\alpha(1)-1$ at the top state.  A bounded-continuous non-equilibrium test continuation makes the canonical top response zero, proving that the floor cannot follow uniformly from monotonicity or canonical selection.
\item \path{mwShockedOne_sender_endpoint_spread_le},\\
\path{mwShockedOne_target_endpoint_spread_le},\\
\path{mwShockedOne_no_top_floor},
\path{mwShockedOne_no_dynamic_floor}, \path{mwEqEquilibrium_realizedRule_one_lt_one}, \path{mwEqEquilibrium_no_dynamic_floor} --- The equilibrium fixed points close the stronger question in the negative. The sender spread bound makes a putative top floor imply $\alpha(1)\le17/50$; the target spread bound then controls the continuation loss from deviating to action $1/2$. The deviation is strictly profitable, so no full-width shocked equilibrium satisfies the top-state floor or any non-empty-basin floor ending at state $1$. The last two declarations instantiate that impossibility at the explicit unconditional equilibrium: its top sanction is strictly below one and its monotone rule cannot satisfy the dynamic floor.
\item \path{basin_forward_invariant_iff_le_fullShockDrift}, \path{basin_forward_invariant_under_box_shock_generalF}, \path{basin_escape_of_box_shock_wide_generalF}, \path{basin_forward_invariant_of_closedLoop}, \path{basin_escape_of_closedLoop} (\path{Theory/Shocked/ShockedHysteresis.lean}) --- Consequence (iv) for every transition (round BH-6, item BH6.7): for $0<x_T\le1$ the high basin $[x_T,1]$ absorbs every shock in $[-h,h]$ after full sanctions if and only if $h\le d_{\F}(x_T)=\F(x_T,1)-x_T$, the one-period push full sanctions give the floor state, and a wider shock from the floor leaks below it (Goal 3 hysteresis, thm:shock). The affine pair \path{basin_forward_invariant_under_box_shock}, \path{basin_escape_of_box_shock_wide}, with the threshold $\delta(1-x_T)$, are its instances (\path{Examples/AffineInstance/LockInShocked.lean}).
\item \path{movAvgReal_affine_lipschitz_contraction} --- The affine box-smeared transfer operator being a $(1-\delta)$-contraction in the dual Kantorovich-Rubinstein metric (uniqueness of and convergence to the invariant law).
\item \path{dualKR_le_of_adjoint_lipschitz}, \path{dualKR_eq_zero_of_invariant}, \texttt{dualKR\_geometric} --- The abstract dual-KR contraction and its two consequences: a unique invariant law and geometric convergence.
\item \texttt{ext\_of\_Iic} --- The measure-separation (two measures with equal Lipschitz-test integrals are equal) via the cumulative-distribution route.
\item \path{eq_of_dualKR_eq_zero}, \texttt{eq\_of\_invariant} --- Genuine measure equality from dual-KR distance zero (uniqueness of the invariant law is real measure equality, not merely d=0).
\item \path{boxPullback_lipschitz} --- The affine box kernel's dual pullback being $(1-\delta)$-Lipschitz on $1$-Lipschitz test functions (discharges the abstract Lipschitz hypothesis).
\item \texttt{integral\_map}, \texttt{integral\_prod},
\texttt{integral\_bind}, \texttt{boxTransfer},\\
\texttt{boxTransfer\_adjoint}, \path{boxTransfer_dualKR_contraction},\\
\path{boxTransfer_unique_invariant},
\path{boxTransfer_dualKR_geometric} --- The box transfer operator T* and its adjoint identity being constructed (not cited), with contraction/uniqueness/geometric convergence; integral\_map/integral\_prod/integral\_bind are mathlib lemma names.
\item \path{exists_invariant_of_adjoint}, \path{boxTransfer_exists_invariant} --- Existence of the invariant law via the Krylov-Bogolyubov compactness fixed point, with the adjoint identity supplying linearity and weak continuity.
\item \path{boxTransfer_existsUnique_invariant_attracting} --- The headline result: the unique, globally attracting invariant law of the affine box-transfer chain (end-to-end derivation).
\item \path{closedLoopBoxTransfer},
\path{closedLoopBoxTransfer_exists_invariant},\\
\path{closedLoopBoxTransfer_existsUnique_invariant_attracting},\\
\path{equilibrium_closedLoop_exists_invariant},\\
\path{closedLoop_dualKR_route_blocked},
\path{closedLoopBoxTransfer_affine_const} --- The general closed-loop law (sec:existence:machinery): the state-dependent-drift transfer operator with its Feller adjoint; Krylov--Bogolyubov existence for every continuous drift; and, conditionally on a drift Lipschitz modulus $L\in(0,1)$, a unique invariant law with the geometric bound $L^n$.  The equilibrium instantiation gives a stationary law when the equilibrium-induced drift is continuous.  The example economy with a positive stage slope then marks the honest uniqueness boundary: its closed-loop factor is exactly $1+\delta/8>1$, so the contraction route is provably unavailable.  The affine compatibility lemma makes the affine box-transfer operator the drift $g_\delta(x)=(1-\delta)x+\delta$ specialization by definitional equality.
\item \path{shockedFeedbackStackelbergTop_interior_discharged_canonical} (headline status retired in round V, 2026-09-15; the theorem stays) --- The interior-discharged canonical theorem carrying the consumed action-Lipschitz coupling hypothesis over the genuine target payoff, with a concrete atomless-kernel witness for non-vacuity.
\item \path{ratchetParams_strictAdjustment}, \path{ratchetParams_feedbackStackelbergTop_corner_exists} --- The non-affine up-fast/down-slow ratchet (one-quarter downward adjustment, one-half upward adjustment) satisfies the maintained strict-adjustment condition and supports the constructed corner equilibrium, confirming that the general-F construction is not silently affine-only.
\item \texttt{SymbolicEconomy}, \path{symbolicEconomy_exists_strict},
\path{symbolic_MS_certified_family_nonvacuous} --- The affine symbolic
family of \Cref{prop:symbolic-family}: nine free payoff coefficients with
membership cut out by finitely many inequalities (strict: the implacability
and sender stage conditions, $0<\delta<1$, $\beta_S,\beta_T<1$, $q,u_0,\ell>0$
and $0<\bar s<1$ for the sanction floor $\bar s$; weak: $\beta_S,\beta_T\ge0$,
$q\le r$, $u_0+r\le k$, $n\ge0$, $\ell\le c-n$, $0\le\bar x\le\bar s$ and
the entrenchment condition), non-emptiness by
backwards Fourier--Motzkin elimination rather than by naming a point, and
the certificate instantiated at a member satisfying the two strictness
side conditions. \path{certificate} is argument-free; \path{certificateMS}
takes a strictly interior basin bottom only (the strict substitution condition it
also took until round BG-3 is not needed, item BG3.7), which is why clause
(ii) of the proposition scopes the two forms differently.
\item \texttt{SymbolicEconomyF}, \path{SymbolicEconomyF.persistence_generalF},
\path{persistence_generalF_corner}, \path{SymbolicEconomy.toF} --- The
general-$\F$ family and \Cref{prop:primitive-persistence}. The transition
is a structure field constrained only by slope bounds, never by functional
form; \path{persistence_generalF} takes \emph{no arguments beyond the
structure}, which is the formal content of ``from primitives''.
\path{persistence_generalF_corner} pins the realized sanction at one at
every date, and is the declaration the corner qualification cites --- the
monotone conclusion is flat, not a state-dependent ratchet.
\path{SymbolicEconomy.toF} embeds the affine family as the
$g_{\mathrm{lo}}=g_{\mathrm{hi}}$ subfamily, so the general family
extends \Cref{prop:symbolic-family}'s rather than competing with it;
the affine proposition is still proved on its own certificates, not
re-derived through the embedding. Since round BJ (items BJ4.1 and BJ4.2) the
general family carries no sender stage condition and no sanction-floor
parameter: the basin bottom only satisfies $\bar x<1$, and the sender's strict
stage monotonicity on the sub-pivot range follows from $u_0>0$, $q\le r$ and
$\bar x\ge0$; feasibility of $\F$ (\path{SymbolicEconomyF.F_mem}) and
$g_{\mathrm{lo}}\le g_{\mathrm{hi}}$ are theorems, the field $0<d_{\mathrm{lo}}$
is weakened to $0\le d_{\mathrm{lo}}$ with strictness derived from implacability
(\path{dlo_pos}), and the positive corner payoff (\path{cornerStagePayoff_pos})
rests on implacability alone (\path{wedge_lt_stageSlope}). Since round BJ
(item BJ4.3, 2026-10-08) the family carries no curvature block: the upper
bound $d_{\mathrm{hi}}$ on the sanction slope, the curvature ceilings $M_F$ in the sanction
and $M_{\mathrm{state}}$ on the corner section, their signs and the curvature
condition are deleted, and with them the Riccati semiconcavity clause of the sender
invariant set, which is now the gain ball alone (\path{SymbolicEconomyF.carrierS}).
The block bought the sender's argmax uniqueness at every real compliance level;
the value class reads it only at feasible compliance, and there it follows from
strict monotonicity of the sender's dynamic objective in the sanction
(\path{senderObj_strictMonoOn}, \path{SymbolicEconomyF.DAUOnI_on_cone}): full
sanction is the unique maximizer, however curved $\F$ is. The statement of
\path{SymbolicEconomyF.persistence_generalF} is unchanged, but its conjunct that
the equilibrium sender value lies in the invariant set no longer includes
semiconcavity; \path{persistence_generalF_semiconcave} adds it back, at the
Riccati modulus $\beta_SK_SM_{\mathrm{state}}/(1-\beta_Sg_{\mathrm{hi}}^2)$,
for any corner-section modulus $M_{\mathrm{state}}\ge0$ supplied as a hypothesis
(\path{eq_V_S_semiconcave}, \path{carrierS_semiconcave_invariant}), and needs
neither the sanction curvature nor the upper bound on the sanction slope. \path{SymbolicEconomy.toF}
needs nothing beyond the affine family's own fields. The class is strictly
larger than before: it contains the wobble file's own published transition (next
item), which the deleted cap $M_F<d_{\mathrm{lo}}/2$ excluded. The persistence
theorem is proved on a core (item BJ4.4; Rule~3 item~6):
\texttt{SymbolicEconomyFPersist}, the family with the target's three conditions
$0\le n$, $\ell\le c-n$ and the entrenchment condition replaced by the single sign
$n+v\ge0$ (\path{SymbolicEconomyFPersist.persistence_generalF}).
\path{SymbolicEconomyF.persistence_generalF} (headline row 30) is its corollary
through \path{SymbolicEconomyF.toPersist}, with its statement unchanged, and the
core is strictly larger: \path{wobbleEconPersist}, the wobble economy with
sanction pain $c=100$, fails the entrenchment condition and so is no member of the family
(\path{wobbleEconPersist_not_toPersist}), and persistence holds there. Clauses
(ii)--(iii) of \Cref{prop:symbolic-family} on the general family (item BJ4.6):
\path{SymbolicEconomyF.symbolic_family_generalF} and
\path{regime_separation_symMyopic}, with no side condition against the
family's own myopic equilibrium; against the canonical myopic equilibrium,
\path{SymbolicEconomyF.regime_separation} takes $0<\bar x$ and $u_0+r<k$.
\item \path{wobbleEconMS}, \path{wobbleEconF},
\path{wobbleEconF_not_strict_placation},
\path{wobbleEconF_persistence}, \path{wobbleEconMS_F_ne_affineF},
\path{wobbleEconFHalfTenth}, \path{wobbleF_half_tenth_realisable},
\path{wobbleEconFHalfTenth_curvature}, \path{wobbleEconFHalfTenth_persistence} --- The
non-affine example economies. \path{wobbleF} deletes the transience factor that
made earlier bump transitions affine at the corner, so its corner section
is a genuine parabola and no adjustment rate makes it affine.
\path{wobbleEconF} sits exactly on the substitution boundary $u_0+r=k$, so
it fails the strict substitution condition
(\path{wobbleEconF_not_strict_placation}). Since round BG-3 (item BG3.7) the
\path{certificateMS} constructor no longer takes that condition, and the member
carries the certificate (\path{wobbleEconF_certificateMS},
\path{wobbleEconF_MS_certified_without_strict_placation}), so dropping the
condition was a strict generalization. The condition is still a hypothesis of the
separation of the myopic and forward-looking cases against the canonical myopic
equilibrium, which asserts a zero myopic sanction at the top state. Prose citing
the refutation should say ``fails the strict substitution condition'' rather than
``has no certificate''.
Since round BJ (item BJ4.3) the family also contains
\path{wobbleEconFHalfTenth}, whose transition is the one the wobble file
publishes, with bump coefficient $\kappa=1/10$ and realized recursion
$x'=x/2+1/2+(1-x)^2/10$ (\path{wobbleF_half_tenth_realisable},
\path{wobbleEconFHalfTenth_realised}). Its transition is semiconcave in the
sanction and along the corner section with modulus $1/5=2\kappa$ and with no
smaller modulus (\path{wobbleEconFHalfTenth_curvature}), above the cap
$M_F<d_{\mathrm{lo}}/2=3/20$ that the deleted curvature condition implied, and
persistence holds there with the full sanction at every date
(\path{wobbleEconFHalfTenth_persistence}). The earlier refutation that no member
carries this transition followed from the deleted condition, not from persistence,
and is deleted with it. Existence only (Rule~9).
\end{itemize}

\paragraph{Equilibrium dynamics (\Cref{sec:results}).}
\begin{itemize}[leftmargin=1.4em,itemsep=1pt]
\item \path{GeneralFPersistencePackage} and
\path{standardCompliance_persistence_generalF_from_package} --- The
paper-facing interface for \Cref{thm:persistence}. Its fields are
exactly the two value-set memberships, identification of the
equilibrium target policy with the selector the value sets define, Lipschitzness
of the sender value set, state-indexed interval dominance after target-policy
substitution, the separate transition-adjusted sub-pivot slope condition (since round BJ, item
BJ1.3, on the basin only and in pivot-star form: each sub-pivot sanction
against the pivot), and
the target combined slope condition. The conclusion is basin invariance, the
basin sanction floor, and a non-decreasing realized-sanction sequence.
Neither stage concavity nor sender-value monotonicity is a hypothesis.
\item
\path{GeneralFPersistencePackage.toJointPreferenceBasinTop},
\path{jointPreferenceBasinTopF_long_run_with_equilibrium_evidence},
and
\path{standardCompliance_persistence_generalF_with_equilibrium_evidence_from_package}
--- The stronger companion interface makes the role of P4 explicit.
The package is first converted into the joint equilibrium basin, using
the transition-adjusted sender slope condition for the level condition, A1 for
monotonicity of the sanction rule, and A2 for target
sanction preference.  Along every resulting general-$\F$ path, the
theorem returns the three conclusions of \Cref{thm:persistence}
together with feasibility of the target policy, target Bellman
optimality, sender dynamic optimality, and monotonicity of the
target's complete sanction objective on the certified high-action
range.  Thus A2 is load-bearing for the extended equilibrium-evidence
conclusion even though the shorter kinematic projection records only
invariance, the floor, and sanction monotonicity.
\item \path{SenderPersistencePackage},
\path{SenderPersistenceBasinTop}, and
\path{standardCompliance_persistence_generalF_from_sender_package} ---
The P4-free counterpart, which is what makes the division of labor
asserted in \Cref{sec:results} formally exact rather than a reading of
how the package happens to be assembled. \path{SenderPersistencePackage}
is \path{GeneralFPersistencePackage} with the field for the target combined
slope condition deleted and nothing else changed;
\path{GeneralFPersistencePackage.toSenderPersistencePackage} is the
projection, and an \path{example} beside it re-derives the full
package's persistence conclusion through the slim route, so no content
is lost. \path{SenderPersistenceBasinTop} is the corresponding
basin object. It is free-parametric in the sender value and the target
policy, exactly as \path{JointPreferenceBasinTop} is, so the object
itself is a kinematic condition rather than an equilibrium statement;
the equilibrium anchoring lives in the \emph{trajectory} theorems
\path{equilibriumTop_senderPersistenceF_long_run_nonDecreasing} and
\path{equilibriumTop_senderPersistenceF_long_run_with_evidence}, which
are what every consumer in the development actually calls. The equilibrium-local form (round I) is \path{standardCompliance_persistence_generalF_local}: the same conclusions from a Lipschitz modulus of the equilibrium sender value, the sub-pivot slope condition at the equilibrium policy and interval dominance of the substituted objective, with \path{SenderPersistencePackage.toLocal} the bridge from the package over the value sets; a weakening of the same conditional theorem, constructing no equilibrium. Its pivot-star form (round M) is \path{standardCompliance_persistence_generalF_local_star}: one comparison per sub-pivot sanction against the pivot, through \path{senderDynObjective_lt_pivot_of_star_chord_and_V_lipschitz}, \path{sender_high_argmax_exists_of_pivot_beats_subpivot} and \path{SenderPersistenceBasinTop.of_local_star}, the pairwise form being its corollary (\path{star_chord_of_pairwise_chord}, \path{standardCompliance_persistence_generalF_local_of_star}); since round BJ (item BJ1.3) the value-set-wide (P5) of \Cref{thm:persistence} is itself the pivot-star slope condition on the basin, in both packages, and the pairwise form on $[\bar x_S,1]$ is a sufficient condition for it. Its level part needs no interval dominance (round BH-2): \path{level_dyn_top_of_local_star} and \path{standardCompliance_persistence_generalF_level_star} give basin invariance, the floor, a non-decreasing state path and the target's optimality from the Lipschitz modulus and the pivot-star slope condition alone, through \path{level_dyn_top_of_high_argmax}, so A1 is read only for sanction monotonicity.
The basin object carries no hypothesis on how the target ranks sanction levels, which
is what makes P4 dispensable for (i)--(iii); the target's equilibrium
policy is of course still present, since it is an argument of the
sender's substituted objective. The level condition comes from the sub-pivot
slope condition against Lipschitzness of the sender value set and the monotonicity condition
from interval dominance, exactly as \Cref{lem:single-crossing} splits
them.
\item \path{equilibriumTop_senderPersistenceF_long_run_nonDecreasing} and
\path{equilibriumTop_senderPersistenceF_temporary_shock_permanent_high} (the
headline row is its sibling \texttt{\ldots\_of\_preBasinAscent}, with
full-sanction ascent only below the basin; round BC)
--- The equilibrium-anchored forms, and the ones to cite for
\Cref{thm:persistence} conclusions (i)--(iv) without P4. Each takes an
actual \path{FeedbackStackelbergEquilibriumTop} rather than a
value pair, so the target policy driving the trajectory is feasible
and Bellman-optimal by construction, and each attaches the per-period
target-argmax certificate --- conclusion (iv) --- along the whole
realized path. The pairing
of a free-parametric basin theorem with an anchored sibling in the same
file follows
\path{equilibriumTop_jointPreferenceBasinTop_long_run_nonDecreasing};
the parametric layer is retained deliberately, since it is what makes
the Rule-13 recovery check meaningful at full generality.
\item \path{equilibriumTop_senderPersistenceF_long_run_with_evidence} ---
The declaration that makes \Cref{sec:results}'s division of labor
formally exact. Along the equilibrium path it returns basin invariance,
both sanction floors, sanction monotonicity, feasibility and Bellman
optimality of the target action, and optimality of the realized sanction
for the sender: that is conclusions (i)--(v) of \Cref{thm:persistence}
with the target combined slope condition \emph{deleted from the hypotheses},
leaving P4 responsible for conclusion (vi) alone --- exactly what
\Cref{sec:results} asserts. No hypothesis is added to obtain (iv) and
(v): target feasibility and optimality come from the equilibrium's own
$\alpha$\path{_mem} and $\alpha$\path{_opt_top} fields, and sender optimality from
$\sigma$\path{_dyn_top_isMaxOn}, the greatest selection being a maximizer by
construction.
\item \path{StandardCompliancePersistenceCertificate} and
\path{exists_persistent_equilibrium_generalF} --- The persistence-only
certificate, which additionally \emph{constructs} the equilibrium
rather than receiving it. Relative to the Milgrom--Shannon core
certificate it drops the target combined slope condition, the myopic ceiling, and
strictness of the basin pivot. The target condition pays for conclusion (vi)
and the myopic ceiling for \Cref{thm:dominance}; finite shock entry in
\Cref{thm:shock} needs only full-sanction ascent below the basin, not a
strict pivot (round BH-1: \path{StandardCompliancePersistenceCertificate.two_goals_generalF} and
\path{StandardCompliancePersistenceCertificate.three_goals_generalF_of_preBasinAscent}, headline
row 1, with its minimum-drift sibling, row 2); the strict pivot returns in
\path{exists_persistent_equilibrium_after_shock_generalF}.
\path{StandardComplianceLocalCoreCertificate_MS.toPersistence} is the
projection.
\item \path{standardCompliance_min_drift_shock_generalF_from_sender_package_of_preBasinAscent}
--- The paper-facing interface for \Cref{thm:shock} (round BH-1). It takes the
P4-free persistence package of the preceding items and adds full-sanction ascent on
$[0,\bar x]$, from which $\bar x<1$ follows; continuity and
\textup{(F1)--(F4)} are already primitive fields. The full-package form
\path{standardCompliance_min_drift_shock_generalF_from_package_of_preBasinAscent}
(round BC) is its corollary through the projection, and the predecessor
\path{standardCompliance_min_drift_shock_generalF_from_package}, with strict
adjustment \textup{(F5)} and $\bar x<1$ in place of the ascent, is a
corollary of that. It returns an attained positive minimum
drift, the uniform ceiling duration, and the complete post-entry
persistence conclusion.
\item \path{senderStage_isMaxOn_unique_of_strictConcaveOn} and
\path{stage_strictConcavity_does_not_force_dynamic_uniqueness} ---
A3 implies uniqueness only for the sender's stage argmax. The second
declaration is an explicit counterexample: a strictly concave stage
payoff plus a continuation term yields a constant complete dynamic
objective with two distinct maximizers.
\item \path{eeCert_sender_continuation_reversal},
\path{eeCert_target_continuation_reversal}, and
\path{eeCertMS_mechanism_active_integrated_application} --- At one
certificate-produced equilibrium both values are strictly increasing,
the sender and target each strictly reverse a named stage ranking when
continuation is added, and the complete general-transition persistence
and dynamic-over-myopic dominance conclusions hold. This is the
two-sided mechanism-active full-package application. Its equilibrium
policies are at corners and its cross-value rows are decoupled.
\item \path{b2_integrated_generalF_persistence_witness},
\path{B2IntegratedGeneralFOutcome}, and
\path{b2_integrated_mechanism_active_generalF_witness} ---
\Cref{thm:b2-benchmark}: one
constructed equilibrium has positive discount factors for both
players, a strictly interior realized sender sanction at $x=1/2$, the
complete general-$F$ persistence floor for every basin start, and the
minimum-drift shock conclusion for every feasible shock start. The
strengthened corollary proves on the same equilibrium that the target
and sender values have discounted changes $9/2$ and $1/4$ between
explicit feasible states. It also records
$\alpha^E(1/2)=0$, $s^E(1/2)=7/8$, and the myopic-equilibrium sanction
$s^M(1/2)=5/14$. Thus the full-package application is
mechanism-active, not a constant-value artifact. Its cross-value rows
are nevertheless decoupled, and its concrete transition is the affine
member of the general-$F$ class; it does not establish a
nonlinear-transition witness with an interior target policy.
\item \path{niEquilibrium}, \path{niJointBasin},
\path{niParams_F_ne_affine}, and
\path{ni_integrated_nonaffine_interior_persistence} --- The
deliberately minimal integrated witness.  The displayed structure is
an exact equilibrium in $\mathcal E$ with
$\beta_T=9/10$, $\beta_S=1/2$, target policy $1/2$, sanction
rule $(x+1)/2$, and the quadratic-bump transition, which is proved
unequal to the affine law.  The joint basin supplies the full
general-$F$ path floor, path invariance, sanction monotonicity, both
players' on-path optimality, and strict target sanction preference.
Both equilibrium values are identically zero, a limitation stated in
the text: the declaration closes logical feature integration but does
not give the forward-looking incentives any economic effect. Since round BH-5 (item
BH5.4) the witness is stated over the transition of any economy, with no
hypothesis, because the successor state is never priced (\path{niParamsOf},
\path{niEquilibriumOf}, \path{ni_integrated_interior_persistence_generalF}).
\item \path{maEquilibrium}, \path{maV_T_strictMono},
\path{ma_target_continuation_changes_pivot_choice}, and
\path{ma_integrated_nonaffine_mechanism_active_persistence} --- The
continuation-adjusted nonlinear interior witness. It preserves the
quadratic-bump transition, target policy $1/2$, sanction rule
$(x+1)/2$, and the complete general-$F$ basin conclusions, while
replacing the target value by $V_T(x)=x$. At the pivot, action $9/16$
strictly improves the target's stage payoff relative to $1/2$, but the
complete Bellman objective strictly rejects the deviation. Thus the
target's forward-looking incentive changes behavior in an exact integrated
economy. The sender value remains zero, so the declaration is not
presented as a two-sided coupled mechanism witness. Over the transition of any
economy (round BH-5, item BH5.4) the equilibrium, the basin and the persistence
evidence hold with no hypothesis, and the stage payoff's preference for the
deviation at the pivot, which the complete objective rejects, holds exactly when
$\F(\tfrac12,\tfrac34)-\F(\tfrac12,\tfrac{47}{64})>5/1152$, which a lower bound
above $5/18$ on the sanction slope implies (\path{maParamsOf},
\path{maOf_pivot_choice_reversal_iff}, \path{maOf_pivot_gap_of_lowerChord},
\path{ma_integrated_mechanism_active_persistence_generalF}).
\item \path{standardCompliance_three_goals_local_MS_generalF_uniform_shock} --- the three headline dynamics results (Persistence, Dominance, Shock) under the general transition law F, with one finite full-sanction duration working uniformly for every initial state; thm:persistence, thm:dominance, thm:shock.
\item \path{eeCertCertificateMS}, \path{qbCertCertificateMS} ---
non-vacuity of the general-$F$ standard-compliance certificate via two
complete instances: the affine witness (sec:existence:nonvacuity) and
the non-affine transition.
\item \path{eeCertMS_generalF_three_goals}, \path{qbCertMS_generalF_three_goals}, \path{quadBumpF_ne_affineF}, \path{quadBumpF_corner_isAffine} --- the three-goals conclusions firing with the trajectory evolved by the general transition, and the genuine non-affineness of that transition, now refuted at \emph{every} adjustment weight rather than only at the economy's own \(\delta\) (\path{qbParams_F_ne_affine} is the \(\delta=1/2\) corollary). The scope of this particular witness is also pinned: the quad-bump's transience factor \((1-s)\) vanishes at full sanction, so its corner section \(y\mapsto\F(y,1)\) is affine and the realized recursion at this economy is the affine one. This row therefore certifies the \emph{incentive} layer under a non-affine transition, not the realized dynamics; for those see \path{quadBumpLowF_corner_not_affine} and \path{wobbleEconF10_corner_not_affine}.
\item \path{eeCert_myopic_argmax_floor} --- the myopic argmax floor at the example economy (every myopic best response \textgreater{}= min(1,(5x+10)/11)), discharging the myopic-ceiling certificate field.
\item \texttt{hV\_S\_mono\_carrier} --- the monotonicity input for the sender invariant set (each sender value in the invariant set non-decreasing in the political state) of the general-F certificate.
\item \path{monotone_of_mem_valLipGainBall}, \path{hV_S_mono_carrier_of_gainBall_subset}, \path{theorem2_dynamic_dominates_myopic_at_same_alpha_of_mem_valLipGainBall}, \path{eeCert_same_alpha_dynamic_dominance_from_gain_carrier} --- the derivation that every member of a zero-based gain-ball invariant set with nonnegative slope is monotone, followed directly by the dynamic-versus-stage comparison at a fixed target policy, including the concrete certificate witness.
\item \path{Params.StrictAdjustment}, \texttt{F\_zero\_lt\_self}, \texttt{hF\_top}, \texttt{F\_eq\_self\_iff}, \texttt{between\_of\_s\_lt}, \texttt{between\_of\_x\_lt} --- the maintained strict-adjustment condition and its machine-checked consequences: zero sanctions strictly lower every positive feasible state; full sanctions strictly raise every state below one (which gives the ascent below the basin that thm:shock uses); feasible transition fixed points lie exactly on $s=x$; and adjustment cannot overshoot the current sanction.
\item \path{targetDynBCF_shock_row_lipschitz} --- the finite target-row cross-Lipschitz bound L\_TS under the purified/atomless shock.
\item \texttt{h\_obj\_diff} --- the isolated sender-row cross-value/action coupling assumption (the unsettled half of the cross-Lipschitz question); it is one member of the general interior theorem's three-condition analytic package and, together with composite leader strong concavity, one of the two inputs left in the \emph{abstract} full-width state-local specialization. The exact $P_3$, $h=1$ construction bypasses rather than discharges this reusable-theorem input.
\item \path{standardCompliance_three_goals_local_MS}, \path{eeCertMS_three_goals_dominance} --- The affine-transition three-goals theorem, specialized to the affine law (which retains the closed-form shock entry horizon); and, separately, the unconditional dominance corollary at the Eaton-Engers certificate witness, whose conclusion names no trajectory and is therefore transition-independent --- it is proved from the general-transition sibling \path{standardCompliance_three_goals_local_MS_generalF}, not from the affine theorem (thm:persistence, thm:dominance).
\item \path{fullShockDrift}, \path{fullShockDrift_exists_pos_isMinOn},
\path{shockPathF_linear_lower_bound_before_basin},
\path{shockPathF_uniform_natCeil_entry_of_drift},
\path{shockPathF_exists_min_drift_uniform_natCeil_bound_of_strictAdjustment},
and \path{standardCompliance_three_goals_local_MS_generalF_min_drift_shock} (its
sibling \texttt{\ldots\_of\_preBasinAscent} and the three-goal theorem
\path{standardCompliance_three_goals_local_MS_generalF_of_preBasinAscent} were the
headline rows from round BC; since round BH-1 rows 1 and 2 are the composites over
the persistence certificate, which read no target condition,
\path{StandardCompliancePersistenceCertificate.three_goals_generalF_of_preBasinAscent} and
\path{StandardCompliancePersistenceCertificate.three_goals_generalF_min_drift_shock_of_preBasinAscent})
--- The quantitative general-$F$ construction behind thm:shock.
Continuity and (F5) produce an attained positive minimum drift
$\eta_{\F}(\bar x)$ on the compact pre-basin interval. The worst-case
path from zero grows by at least this amount before entry, giving the
uniform bound $\widehat N_{\F}(\bar x)=
\lceil\bar x/\eta_{\F}(\bar x)\rceil$ and equilibrium persistence after
entry. \path{shockEntryHorizon_bound} separately retains the sharper
affine formula $N_*(\bar x,\delta)$, stated without initial-state bounds, which
its proof never read (\path{shockEntryHorizon_bound_of_pow}, round L, is now an
alias).
\item
\path{fullShockDrift_minima_horizon_mono_of_top_transition_le},
\path{fullShockDrift_minima_horizon_strict_of_top_transition_lt},
\path{affine_fullShockDrift_minimum_closed_form},
\path{affine_fullShockDrift_minima_horizon_mono_of_delta_le}, and
\path{affine_quarter_half_fullShockDrift_minima_horizon_strict}
--- The machine-checked content of
\Cref{cor:shock-transition-strength}. Given full-sanction ascent of the
first law below the threshold, and no adjustment hypothesis on the second
(round BG-2), pointwise full-sanction
transition dominance orders the attained minimum drifts and weakly
orders their ceiling durations; strict dominance strictly orders the
drifts but not necessarily the integer durations. Under the affine law
the minimum is $\delta(1-\bar x)$, and the final declaration supplies
an exact directional witness comparing $\delta=1/4$ with $\delta=1/2$.
\item \path{dynamicTopLockInPathF_tendsto_steadyState_full_of_strictAdjustment} --- \Cref{prop:longrun}, general in $\F$: under (F5), at a continuous target policy and a sanction rule at or above the state on the basin (the level condition, with no floor), the state path from every basin start converges to some $x^*$ in the basin, the entire realized-sanction sequence converges to the same $x^*$, and $x^*$ is an equilibrium-consistent steady state, an optimal sender choice at the limit state weakly below the greatest selection. The argument applies to every sanction subsequence: compactness supplies a convergent further subsequence, continuity gives $F(x^*,s^*)=x^*$, and (F5) forces $s^*=x^*$. The base general-transition theorem \path{dynamicTopLockInPathF_tendsto_steadyState} gives state convergence, an optimal subsequential sanction limit $s^*$, and $F(x^*,s^*)=x^*$; \path{dynamicTopLockInPathF_tendsto_steadyState_of_strictAdjustment} uses (F5) to add $s^*=x^*$. The affine-path case is \path{dynamicTopLockInPath_tendsto_steadyState} (affine lock-in path \eqref{eq:affine}), which needs no (F5), because affine inversion recovers the sanction.
\item \path{senderDynObjective_lt_iff_continuation_gain_gt_stage_loss}, \path{targetSanctionObjective_lt_iff_continuation_gain_gt_stage_loss} --- The exact continuation-ranking identity used in the mechanism discussion: a higher sanction is dynamically preferred if and only if its discounted continuation gain exceeds its current stage-payoff loss.
\item \path{no_upper_basin_floor_of_top_lt_one}, \path{mwShockedOne_no_dynamic_floor} --- The structural endpoint obstruction in \Cref{cor:endpoint-obstruction} and its full-width \(P_3\) specialization.
\item \path{dynamicTopLockInStepF_exists_least_fixedPoint},
\path{dynamicTopLockInStepF_exists_least_fixedPoint_of_strictAdjustment},\\
\path{dynamicTopRealizedRule_eq_self_of_stepF_fixedPoint_of_strictAdjustment},
\path{dynamicTopLockInStepF_fixedPoint_iff_realizedRule_eq_self_of_strictAdjustment},
and, on the affine path,
\path{dynamicTopLockInStep_exists_least_fixedPoint},
\path{dynamicTopRealizedRule_eq_self_of_step_fixedPoint} ---
Under monotonicity of the sanction rule on the basin, the general-transition
step map $G(x)=F(x,\sigma_{\mathrm{real}}(x))$ has a least basin fixed point,
with the transition flat on the band between the state and the sanction there.
Under (F5), the pointwise implication
$F(x,\sigma_{\mathrm{real}}(x))=x\Rightarrow
\sigma_{\mathrm{real}}(x)=x$ strengthens to a full equivalence, and at the
least fixed point the realized sanction equals the state.  Hence
strict-adjustment laws paired with the same sanction rule have the same
fixed-point set, including the same least basin fixed point when it exists.
On the affine path $G(x)=(1-\delta)x+\delta\sigma_{\mathrm{real}}(x)$ the same
holds without (F5), because affine inversion recovers the sanction.
\item \path{dynamicTopLockInPathF_tendsto_below_least_fixedPoint_of_rule_mono},
\path{dynamicTopLockInPathF_tendsto_below_least_fixedPoint_of_strictAdjustment},
\path{monotone_selfMap_orbit_limit_need_not_be_fixed}, and
\path{ratchetParams_generalF_lowerOrbit_below_least_fixedPoint} ---
Monotonicity of the sanction rule makes the general-$F$ selected step
monotone. The path from the lower basin endpoint is bounded above by its least
fixed point, so the equilibrium-consistent orbit limit is weakly below that
fixed point; single crossing supplies the rule monotonicity. Equality is not
claimed: the kernel-checked monotone-map counterexample shows that an orbit can
converge to a jump point that is not fixed. The non-affine ratchet instantiates
the valid bound. Its derived corollary
\path{ratchetParams_generalF_lowerOrbit_and_least_fixedPoint_eq_one}
uses the identically-one sanction rule to sharpen the result at that example economy to
$x_L=x_F=1$ without changing the general theorem or adding an assumption.
\item
\path{equilibriumTop_tailF_no_relief_of_ratchet},
\path{jointPreferenceBasinTopF_tail_no_relief},
\path{equilibriumTop_tailF_sender_optimal_against_relief}, and the
state-local uniqueness upgrades in
\path{PostEntryNonTermination.lean} --- the post-entry
non-termination result of \Cref{cor:nontermination,app:nontermination}:
every cutoff strictly below the inherited floor has empty first-passage
set, although relief remains feasible and is compared directly with the
realized dynamic best response. The empty first-passage set needs only
the ratchet $x\le R(x)$ on the basin at an equilibrium (round BH-2); the
joint-basin forms, which carry the target condition, are its corollaries.
Selector independence and strict
suboptimality of every distinct relief action require the existing
state-local dynamic-argmax uniqueness condition.
\item \path{pathwise_dominance_generalF}, \path{standardCompliance_three_goals_local_regime_comparison} --- The conditional pathwise-ordering corollary (cor:pathwise-dominance), general in $\F$ since round BH-1: from the core certificate, which since round BJ (item BJ1.4) carries no target entrenchment-condition field, so that (P4) is not maintained, and the myopic-basin-preservation hypothesis eq:myopic-basin-preservation, at every basin state the myopic realized sanction is at most the forward-looking one, and from every common basin start the myopic sanction path evolved by $\F$ stays weakly below the forward-looking one at every date. The second declaration is the earlier affine, fixed-start form. \path{eeCert_hMyoBasin_counterexample} machine-checks that the hypothesis fails at the flagship example economy (under the myopic equilibrium the state leaves every high basin), confirming its non-genericity; \path{ruleSanctions_le_of_ceiling_and_path_floor} and \path{eeCert_pathwise_dominance} recover the pathwise conclusion there without basin preservation, via a global myopic ceiling and a corner floor under the dynamic path. General-transition forms: the cross-policy comparator and the enrichment variants now carry generalF siblings on the F-evolved paths (\texttt{ruleStepF}/\texttt{ruleSanctionsF}; \path{standardCompliance_three_goals_local_MS_regime_comparison_generalF}, \texttt{\_with\_path\_optimality\_generalF}, \texttt{\_with\_selector\_equivalence\_generalF}, \texttt{\_weak\_persistence\_generalF}, and the low-basin path-optimality twin) via the MS certificate; the pathwise comparison under the entrenchment-condition certificate remains affine because its condition $c_T\le\beta_T g_T\delta$ prices the affine slope (the restriction already noted for the cardinal form).
\item \path{totalForwardWedge_eq_targetContinuation_add_senderSwitch}, \path{targetContinuationWedge_eq}, \path{alpha_order_dynamic_le_myopic_of_strictSC}, \path{targetFF_strictSCLower_targetMM_of_two_legs} --- The exact two-channel decomposition of the forward-looking gap (\Cref{prop:two-channel}) and the ordinal form of the compliance ordering (\Cref{prop:ordinal-ordering}). The three target objectives are named \path{targetMMObjective}, \path{targetTMObjective}, and \path{targetFFObjective}; the last is definitionally the equilibrium's own Bellman integrand \path{joint_T_top_integrand}, and no maximizer of the middle one is ever constructed. The ordinal relation \path{Sanctions.BinaryComparativeStatics.StrictSCLowerOn} orders every pair of maximizers with no continuity, concavity, compactness, or uniqueness, and \path{StrictSCLowerOn.trans} is what licenses discharging the two channels separately. That the sender-switch channel cannot be signed from the sign conditions alone is the Rule-11 refutation \path{senderSwitch_sign_not_determined}.
\item \path{rulePathF_myopic_le_dynamicTop_of_strictDecreasingDifferences}, \path{ruleSanctionsF_myopic_le_dynamicTop_of_strictDecreasingDifferences} --- The composition of the cross-policy rule ordering with the general-\(F\) path comparators: from a common feasible start the myopic equilibrium's political state and realized sanction stay weakly below the forward-looking ones in every period, with no affine restriction.
\end{itemize}

\paragraph{Mechanism and dual-basin path dependence (\Cref{sec:mechanism}).}
\begin{itemize}[leftmargin=1.4em,itemsep=1pt]
\item \path{standardCompliance_low_basin_generalF_of_sender_level} (the headline row since round BC; no target condition) and \path{standardCompliance_three_goals_local_low_basin_generalF} --- The persistence theorem thm:lowbasin under the general transition law F: the low-basin (epsilon=-1) instantiation, the sign-mirror of the high-basin (epsilon=+1) headline thm:persistence. Its pivot-star sibling \path{standardCompliance_low_basin_generalF_of_sender_star} compares the ceiling with heavier sanctions only (round BH-2). The high orientation reads no target condition either (\path{standardCompliance_persistence_generalF_from_sender_package}), and initial-condition selection runs on the two sender halves (\path{equilibriumTop_initial_condition_selects_generalF_of_sender_basins}, round BH-2). Dropping the target condition enlarges the class (round BH-5, item BH5.6): over any base economy, with its transition and discount factors, a constructed equilibrium that never sanctions satisfies the sender-only low-basin hypotheses at every basin $[0,\bar x]$ and every ceiling, while the joint low-basin condition fails at every basin containing $0$ (\path{lowBasin_senderOnly_strictly_weaker_generalF}, \path{senderLowBasin_not_implies_joint}, \path{Examples/LowBasinSeparationWitness.lean}); a class separation at a zero-sanction equilibrium, silent on persistence while sanctions are active.
\item \path{jointPreferenceBasinLow_long_run_with_equilibrium_evidence}, \path{jointPreferenceBasinLowF_long_run_with_equilibrium_evidence}, \path{StandardComplianceLocalCoreCertificateSchur_MS}, \path{three_goals_local_MS_generalF} --- The low-basin path theorem re-emits target feasibility and Bellman optimality, sender dynamic optimality, and the target's orientation-reversed forward-looking sanction preference at every path state; the general-\(F\) Milgrom--Shannon route is simultaneously available under the weaker Schur package.
\item \path{lbLow_generalF_three_goals}, \path{quadBumpLowF_ne_affineF}, \path{quadBumpLowF_corner_not_affine} --- The constructed dovish-sender low-basin example economy discharging thm:lowbasin at epsilon=-1 (the low political absorbing basin), and the non-affineness of its transition law F at every adjustment weight (\path{lbParams_F_ne_affine} is the \(\delta=1/2\) corollary). Here the bump is carried by \(s\) rather than \((1-s)\), so it survives at full sanction and the \emph{corner section} is non-affine as well --- this example economy bends the realized recursion, not only the incentive layer.
\item \texttt{dualParams}, \texttt{dbCertCertificate}, \path{dbCert_three_goals_high}, \path{dbCert_three_goals_low}, \path{dualParams_initial_condition_selects} --- The dual-basin selection example economy (Table tab:witness-roster): one economy, one equilibrium, carrying the full high-basin certificate on $[19/20,1]$ and the low-basin conclusions on $[0,1/10]$; the target side sits exactly at the entrenchment-condition boundary $c_T=\beta_T\,g_T\,\delta$ (weak-inequality field, equality admissible), which makes the leader integrand independent of the sanction response (\path{dbCert_joint_T_integrand_eq}) and pins the derived leader policy at the corner; a new semiconcave sender invariant set in derivative-free form (\texttt{SemiconcaveOn}, \texttt{dbCert\_bellmanS\_mem}) discharges state-local argmax uniqueness; the selection theorem packages the two certified paths by starting basin.
\item \path{KaempferLowenberg.Polity.xstar_mono_of_lockIn}, \path{xstar_anti_of_washOut}, \path{Examples.KaempferLowenbergWitness} --- The Kaempfer-Lowenberg microfoundation of the transition sign (F3) and basin assignment: lock-in =\textgreater{} x* non-decreasing, erosion of entrenchment by sanctions =\textgreater{} x* non-increasing, plus concrete examples for both basins.
\item \path{Examples.FourGroupWitness} --- The interest-group example in
app:fourgroup, which gives the Grossman--Helpman weighted-sum
aggregation microfoundation of the stage complementarities. Separate
finite-grid value iteration supplies the exploratory dual-basin
diagnostic reported there; no Lean theorem depends on that numerical
calculation.
\item \path{standardCompliance_three_goals_local_low_basin} --- The affine-transition low-basin headline: thm:lowbasin at the affine law, of which the general-F form above is the strengthening.
\item \path{fourGroupCert_three_goals_dominance} --- The interest-group microfounded certificate headline (unconditional three-goals-with-dominance at the interest-group-derived payoffs); the full identifier map for the interest-group construction is in app:fourgroup.
\item \texttt{eqPath}, \texttt{altPath}, \path{eqPath_step_certified}, \texttt{eqPath\_mem\_basin}, \texttt{fgc\_perPeriod\_gap}, \texttt{fgc\_plannerW\_abs\_le}, \path{fgc_plannerValue_summable}, \path{fgc_plannerValue_eqPath}, \path{fgc_plannerValue_altPath}, \path{fgc_planner_stationary_le}, \path{fgc_planner_stationary_lt}, \path{fgc_planner_lockin_val} --- The period-by-period welfare accounting behind \Cref{sec:policy:welfare} and its derivation in \Cref{sec:fourgroup:planner}: the two comparison state paths (the equilibrium orbit under the \emph{derived} corner actions --- \path{eqPath_step_certified} ties the step to the certificate rather than to an assumed policy, and \texttt{eqPath\_mem\_basin} keeps it in the certified basin --- and the feasible zero-sanction orbit); the uniform cube bound $|W|\le 13/2$ and the geometric-comparison summability it buys; the two discounted closed forms \eqref{eq:planner-values}; the start-independent per-period gap $23/10 + 2^{-t}$; and the stationary-point comparison, in the exact weak-plus-strict-above-zero form the Lean proves (every stationary point with positive sanctions is strictly worse than the origin --- but the origin is \emph{not} claimed to be the unique maximizer, since every zero-sanction stationary point $(0,a,0)$ also attains $0$).
\item \texttt{FourGroup.plannerW}, \texttt{plannerW\_distortion}, \path{plannerW_lockin_dominated_iff}, \path{stationaryPlannerValue_eq}, \path{stationaryPlannerValue_lockin_dominated_iff}, \path{stationaryPlannerValue_relief_gain}, \path{fgc_planner_inefficiency} (+ strict and $\beta_S$/$\beta_T$ endpoint corollaries), \texttt{FGCPlannerPath}, \path{fgcPlannerFeasibleValues_bddAbove}, \path{fgc_plannerOptimalValue_lower_bound}, \path{fgc_equilibrium_lt_plannerOptimalValue}, \texttt{fgc\_VS\_star\_fixed}, \texttt{fgc\_VT\_star\_fixed}, \path{eeCert_VS_star_fixed}, \path{eeCert_VT_star_fixed} --- The planner/welfare layer (sec:fourgroup:planner): aggregate welfare (political weights set to one), the ring identity that the political weights are exactly the gap between government objectives and welfare, the parametric flow and discounted stationary flip boundary $(g_J-q_K)+(b_J-b_K) < k_S+d_s$, the stationary discounted relief gain expressed in existing primitives, and the discounted lock-in inefficiency identity at the example economy with exact margin $23/(10(1-\beta))+1/(1-\beta/2)$. The feasible-path structure and cube bound make the set of discounted feasible values bounded above; its supremum exceeds the certified equilibrium value by at least the same exact margin. Since round S (2026-09-15) the supremum is attained (\path{fgc_planner_bellman_upper}, \path{fgc_plannerPath_value_le}, \path{fgcPlannerOptimalValue_eq}, \path{fgc_plannerOptimalValue_gap_eq_equilibrium}; the next two items), so the margin is the exact loss. The four exact Bellman fixed-point pins upgrade the closed forms at the example economy from ``orientation only'' to fixed points.
\item \path{FourGroupCoeffs}, \path{FourGroup.toCoeffs}, \path{toCoeffs_U_S}, \path{toCoeffs_U_T} (\path{Theory/Microfoundations/FourGroupPolity.lean}), \path{toCoeffs_plannerW}, \path{toCoeffs_pathValue} (\path{Theory/Microfoundations/FourGroupPlanner.lean}) --- round BH-6 (2026-10-07, item BH6.2): the eighteen coefficients without the four region-(R) fields. Aggregate welfare, the distortion identity, the stationary comparison, the path accounting, the cube bound, the affine and general-transition gap identities with every sign characterization of the next two items, and the two state-slope bounds are proved over \path{FourGroupCoeffs}, under the same names in that namespace (for example \path{FourGroupCoeffs.pathValue_cornerF_gap}); each polity form keeps its statement and is the coefficient form at \path{FourGroup.toCoeffs}. The welfare layer therefore holds for every coefficient vector, including polities outside region~(R), in which opposition empowerment dominates the state--sanction cross term. The contest polity now extends the coefficients (\path{ContestPolity}) and no longer carries the sign conditions of region~(R).
\item \path{FourGroup.pathValue}, \path{plannerWBound}, \path{plannerW_abs_le_cube}, \path{pathValue_summable}, \path{affineCornerPath}, \path{affineZeroPath}, \path{plannerW_affineCorner_gap}, \path{pathValue_affineCorner_gap}, \path{affineCornerGap_pos_iff}, \path{affineCornerGap_pos_all_discounts_iff}, \path{cornerPathF}, \path{zeroPathF}, \path{residualF}, \path{cornerPathF_affineF}, \path{pathValue_cornerF_gap}, \path{pathValue_cornerF_gap_ge_of_nonneg}, \path{pathValue_cornerF_gap_ge_of_neg} (\path{Theory/Microfoundations/FourGroupPlanner.lean}); \path{plannerValue_eq_pathValue}, \path{fgc_affineCornerGap_eq}, \path{fgc_planner_bellman_upper}, \path{fgc_plannerPath_value_le}, \path{fgcPlannerOptimalValue_eq}, \path{fgc_plannerOptimalValue_gap_eq_equilibrium} (\path{Examples/FourGroupPlannerWitness.lean}) --- the welfare accounting in the primitives (round S, 2026-09-15; doc 87 W1--W2, on the user's word): the discounted path value general in the polity; the coefficient bound that replaces the cube bound $13/2$ of the example economy as the summability certificate; the affine corner and zero paths (affine-labeled, Rule 22) with the per-period gap $H+B(1-\delta)^t$ and the discounted identity $H/(1-\beta)+B/(1-\beta(1-\delta))$ of \Cref{prop:affine-welfare-gap}, where $B=b_J-b_K$ and $H=(k_S+d_s)-(g_J-q_K)-B$; its exact fixed-discount and all-discount sign conditions (no region sign is read by any proof); the recursive endpoint paths under any transition satisfying (F1)--(F4), with the residual-sum identity and the two sign-oriented lower bounds; the evaluation of the identity at the example economy; and the planner's optimum at the certificate economy: the Bellman upper bound with the candidate value $x/(1-\beta/2)$, the telescoped bound on every feasible path through \path{Foundations/InfiniteHorizonVerification}, the attained supremum and the exact loss of equilibrium play relative to the specified feasible zero-sanction path (\Cref{prop:planner-optimum-lower-bound}; the headline inequality \path{fgc_planner_inefficiency_lt} is unchanged and the equality is its sibling). Round U (2026-09-15; doc 87 W6) adds \path{Foundations/DiscountedSignCharacterization.lean} --- \path{DiscountedSign.summable_discounted}, \path{DiscountedSign.tsum_pos_of_lt_of_le}, \path{DiscountedSign.lt_and_le_of_tsum_pos_all}, \path{DiscountedSign.tsum_pos_all_discounts_iff}, \path{DiscountedSign.not_necessary_without_tendsto}: a discounted series $\sum_t\beta^t(C-Bd_t)$ with $d_0=0$, $d_t\in[0,1]$ and $d_t\to1$ is positive at every $\beta\in[0,1)$ iff $C>0$ and $B\le C$, and the convergence cannot be dropped --- and, in the planner module, \path{sanctionLowerChord_le_one}, \path{one_sub_cornerPathF_le}, \path{zeroPathF_le}, \path{separationF_tendsto_one}, \path{pathValue_cornerF_gap_eq_tsum}, \path{pathValue_cornerF_gap_pos_all_discounts_iff}, \path{pathValue_cornerF_gap_of_identity}, \path{fourGroupIdentityProbe}, \path{not_necessary_without_strict_adjustment}: under a positive lower bound on the sanction slope (the existing predicate \path{SanctionLowerChord}) the endpoint paths converge geometrically, their separation tends to one, and the general-transition gap has the same all-discount sign condition as the affine one; the identity transition, which keeps the cube and fixes the diagonal without adjusting, is the Rule 11 boundary at which necessity fails. Round X (2026-09-17; doc 89 D10, the reader's P1, P2 and P4 of 2026-09-16) adds \path{cornerPathF_tendsto_one_of_strictAdjust}, \path{zeroPathF_tendsto_zero_of_strictAdjust}, \path{separationF_tendsto_one_of_strictAdjust}, \path{pathValue_cornerF_gap_pos_all_discounts_iff_of_strictAdjust} and \path{pathValue_cornerF_gap_pos_all_discounts_iff_of_params}: the same all-discount condition under the qualitative (F5) strict adjustment with continuity --- the corner path is monotone and bounded, its limit is a fixed point of $\F(\cdot,1)$, strict adjustment leaves only $1$, symmetrically for the zero path --- the two directions supplied by \path{Params.StrictAdjustment} and every other input by the structure's own fields, so the general-transition welfare sign holds on the paper's maintained class and the slope bound is only its quantitative form; and, in \path{Examples/FourGroupWelfareSign.lean}, \path{fourGroupCertPositiveCost_gap_sign_iff} (the positive-cost gap equals $(3-4\beta)/(5(1-\beta)(1-\beta/2))$: positive iff $\beta<3/4$, zero at $3/4$, negative above) and \path{fourGroupCert_welfare_not_identified_by_game} (the two polities induce \texttt{fgcParams}, and along the equilibrium-derived paths their aggregate welfare ranks the same two paths oppositely at every discount and every physical start).
\item \path{residualF_succ_ge}, \path{pow_le_residualF}, \path{inv_le_tsum_residualF}, \path{cornerChords_of_sectionChords}, \path{pathValue_cornerF_gap_ge_of_nonneg_of_cornerChords}, \path{pathValue_cornerF_gap_ge_of_nonneg_of_sectionChords} (\path{Theory/Microfoundations/FourGroupPlannerStateChord.lean}) --- the sharpened lower bound of \Cref{app:fourgroup} (round BD, 2026-10-01; formerly marked not formalized): under any transition satisfying (F1)--(F4) whose state rises at rate at least $g\ge0$ toward the fixed corners along the two endpoint sections, the residual separation of the endpoint paths obeys $z_{t+1}\ge g z_t$, so $\sum_t\beta^tz_t\ge1/(1-\beta g)$ and, for $B\ge0$, the discounted gap is at least $H/(1-\beta)+B/(1-\beta g)$; general in $\F$, with no closed form, no strict adjustment and no continuity.
\item \path{hawkRescale}, \path{hawkRescale_U_S}, \path{hawkRescale_U_T}, \path{hawkRescale_sigmaSel}, \path{hawkRescale_toParamsF} (\path{Theory/Microfoundations/FourGroupPolity.lean}), \path{hawkRescale_plannerW}, \path{hawkRescale_pathValue}, \path{hawkRescale_cornerF_gap}, \path{abs_cornerF_gap_le}, \path{exists_hawkRescale_cornerF_gap_neg}, \path{exists_hawkRescale_cornerF_gap_pos} (\path{Theory/Microfoundations/FourGroupPlanner.lean}) --- round BH-6 (2026-10-07, item BH6.1): the hawk-rescaling family. Multiplying $g_J$ by $k\ne0$ and dividing the hawks' weight schedule by $k$ leaves both government payoffs, the myopic best response, the region fields and the induced game unchanged under every transition, and lowers the zero-versus-corner planner gap by exactly $(\varphi-1)g_J/(1-\beta)$ under every transition that keeps the cube. For $g_J>0$ some $k>0$ makes the gap strictly negative, and for $g_J<0$ some $k>0$ makes it strictly positive, at every planner discount, every start and every transition that keeps the cube and is non-decreasing in both arguments; each statement is one-directional. The polity of the next item is the member $k=5$ (\path{fourGroupCertHawkScaled_eq_hawkRescale}, \path{fourGroupCertHawkScaled_gap_shift}).
\item \path{fourGroupCertHawkScaled}, \path{fourGroupCertHawkScaled_U_S}, \path{fourGroupCertHawkScaled_U_T}, \path{fourGroupCertHawkScaled_sigmaSel}, \path{fourGroupCertHawkScaled_toParams}, \path{fourGroupCertHawkScaledEquilibriumTop}, \path{fourGroupCertHawkScaled_plannerW}, \path{fourGroupCertHawkScaled_gap_neg_paths}, \path{fourGroupCertHawkScaled_gap_neg}, \path{fourGroupCertPositiveCost}, \path{fourGroupCertPositiveCost_gap_eq}, \path{fourGroupCertPositiveCost_gap_at_zero}, \path{fourGroupCertPositiveCost_gap_at_nine_tenths}, \path{fourGroupCert_variants_fail_all_discounts} (\path{Examples/FourGroupWelfareSign.lean}) --- the sign conditions of region~(R) do not determine the sign of the welfare gap (round S; doc 87 W3): the hawk-scaled polity keeps both government payoffs, hence the game \texttt{fgcParams} and the exact equilibrium of the certificate economy, and reverses the planner's ranking at every discount; the positive-cost polity's ranking turns with the planner's patience. Existence certificates (Rule 9), built and checked in exact arithmetic before the characterization they refute was proved (Rule 11).
\item \path{PlannerPath}, \path{plannerOptimalValue}, \path{zeroValue}, \path{zeroStateSum}, \path{complianceFlow}, \path{exists_complianceFlow_max}, \path{optCompliance}, \path{bellman_le_of_politicalMargin}, \path{plannerOptimalValue_eq_of_politicalMargin}, \path{zeroStateSum_sub_le}, \path{plannerOptimalValue_eq_of_upperChords}, \path{plannerOptimalValue_eq_of_nonpos}, \path{plannerOptimalValue_eq_affineF}, \path{zeroValue_sub_cornerValue}, \path{plannerOptimalValue_loss_of_upperChords}, \path{plannerOptimalValue_loss_of_nonpos}, \path{plannerOptimalValue_loss_affineF} (\path{Theory/Microfoundations/FourGroupPlannerOptimum.lean}); \path{fgcPlannerFeasibleValues_eq_general}, \path{fgcPlannerOptimalValue_eq_general} (\path{Examples/FourGroupPlannerWitness.lean}) --- round BI (2026-10-07, item BI.2; the hypotheses approved under Rule~3): \Cref{prop:planner-optimum-lower-bound}(i) for every polity and every transition that keeps the cube and is non-decreasing in both arguments. A feasible path is any cube-valued triple obeying the transition from $x_0$, the optimum is the supremum of discounted aggregate welfare over them, and the best constant compliance $a^\ast$ is constructed by compactness, not assumed. The verification core says that the zero-sanction path with $a^\ast$ attains the supremum whenever the one-step political term $(g_J-q_K-d_s)s+\beta B\bigl(Y(\F(x,s))-Y(\F(x,0))\bigr)$ is nonpositive, and it reads of the polity only $k_S>0$ and $w_{\mathrm{pl}}\ge0$. The corollaries discharge that term: for $B\ge0$ from an upper bound $L$ on the sanction slope, an upper bound $\lambda$ on the state slope of the zero-sanction section with $\beta\lambda<1$ ($\lambda\ge0$ derived) and $g_J-q_K+\beta BL/(1-\beta\lambda)\le d_s$; for $B\le0$ from $g_J-q_K\le d_s$ alone; and under the affine law for either sign of $B$ (affine-restricted). The certificate optimum \path{fgcPlannerOptimalValue_eq} is the affine instance at $\delta=\tfrac12$, where the condition reads $\beta/(2-\beta)\le14/5$, and the verification at the example economy is kept. The corner path of full sanctions against full non-compliance in the loss corollaries is a feasible comparison path, and it is equilibrium play only where an equilibrium is shown to play it.
\item \path{BeckerContest.pStar}, \path{pStar_isMaxOn_contest},
\path{contest_nash}, \path{contest_nash_unique},\\
\path{ContestPolity.toPolity}, \path{toPolity_lockIn},
\path{contestWitness}, \path{contestWitness_inducedF_eq_affineF} ---
The Becker contest layer (sec:fourgroup:contest): dominant-strategy Nash
pressures with completed-square argmax certificates (existence + uniqueness,
closed form), the constructor deriving the Kaempfer--Lowenberg schedules as
contest objects and obtaining lock-in by applying the existing market-clearing
theorem, and the example instantiation at which the contest-induced transition
literally equals the affine law.  The parametric constructor is isolated in
\path{Theory/Microfoundations/KaempferLowenbergContest}; its concrete
non-vacuity instantiation is in
\path{Examples/KaempferLowenbergContestWitness}, so the Theory layer does not
import an example economy.
\end{itemize}

\paragraph{Comparative statics (\Cref{sec:compstat}).}
\begin{itemize}[leftmargin=1.4em,itemsep=1pt]
\item \path{dynamicTopSanctions_ge_myopicEquilibrium_different_alpha} --- Verifies the different-alpha equilibrium-vs-equilibrium comparison sigM(x,alphaM(x)) \textless{}= sigDynTop\_\{Vs\}(x,alpha(x)) (the two-force composition), the different-alpha corollary of the comparison at a fixed target policy, thm:dominance.
\item \path{eeCert_patience_response_witness}, \path{eeCert_patience_band_witness}, \path{eeCert_weight_band_witness} --- The lock-in comparative-statics witnesses (sec:compstat:lockin): patience -\textgreater{} response monotonicity and patience -\textgreater{} band/lock-in monotonicity at beta\_S=1/2\textless{}3/5.  The positive-weight comparison at w=1/4\textless{}3/4 is retained as a regression witness, but it binds with equality and is not evidence that adjustment speed moves the steady state.
\item \path{theorem2_dynamic_dominates_myopic_at_same_alpha} --- The kernel of the comparison at a fixed target policy: at any state and any fixed target policy, the forward-looking sender's optimal sanction weakly exceeds the myopic sender's (the fixed-policy form eq:thm-dominance-same-alpha of thm:dominance). Since round BH-6 (item BH6.8) the comparison is selection-free: every stage maximizer, not only the selected one, lies weakly below the dynamic best response (\path{exists_dynamic_argmax_ge_of_stage_max}, \texttt{stage\_max\_le\_$\sigma$\_dyn\_top}, and the local form \texttt{stage\_max\_le\_$\sigma$\_dyn\_top\_of\_senderContinuationMonotoneAt}); as rules, the greatest stage maximizer lies below the dynamic best response (\path{stageTopRule_le_dynamicTopRealizedRule}); and at $\beta_S=0$ the two coincide with no uniqueness hypothesis (\texttt{$\sigma$\_dyn\_top\_eq\_sSup\_stageArgmax\_of\_$\beta$S\_zero}). Every comparison is against the dynamic best response; that every stage maximizer lies below every maximizer of $J_S$ is false under a non-decreasing value alone.
\item \texttt{$\sigma$\_dyn\_top\_le\_of\_beta\_le},
\path{dynamicTopLockIn_band_le_of_beta_le},
\path{dynamicTopLockIn_band_mono_in_weight},\\
\path{dynamicTopLockIn_least_fixedPoint_le_of_rule_le},
\path{dynamicTopLockInPathF_le_of_rule_le_on_basin},\\
\path{dynamicTopLockInPathF_limit_le_of_rule_le_on_basin}, and
\path{dynamicTopLockInPathF_limit_le_of_beta_le_generalF} ---
The parametric comparative-statics results behind
\Cref{sec:compstat:lockin}.  The older fixed-point results order
extremal basin steady states.  The new general-$F$ path results separately
show that, under a shared transition and a monotone upper rule, a pointwise
ordering of the rules orders every realized-path iterate and hence the convergent
selected limits.  The patience specialization holds continuation value
and target policy fixed, matching the direct-effect scope stated in the
text.  The affine adjustment-weight theorem is only a weak comparison:
for positive weights the fixed-point set is invariant, and the theorem
may bind with equality.  It is not a general-$F$ selected-limit theorem.
\item \path{ratchet_dichotomy}, \path{rule_top_eq_one_of_floor},
\path{successor_top_eq_one_of_floor} --- the ratchet frontier
(\Cref{prop:ratchet-dichotomy}) in
\path{Sanctions/Theory/Equilibrium/AbsorbingTopObstruction.lean}.  For any
feasible rule on any upper interval ending at the top state, either the top
sanction is interior and the ratchet floor fails on every such interval, or it
is at the corner and the realized successor at the top is the top itself.  No
slope, curvature or payoff condition, and no hypothesis on either value
function: feasibility plus the transition's diagonal property.
\path{no_equilibrium_of_slope_and_top_gain} supplies the quantitative half
(\Cref{prop:absorbing-obstruction}); it is a sufficient-condition obstruction,
not a general impossibility theorem.
\item \path{no_dynamic_floor_of_antitone_state_chord}, \path{dynamic_stake},
\path{dynamic_stake_equilibrium} --- \Cref{prop:dynamic-stake} in
\path{Sanctions/Theory/LockIn/DynamicStake.lean}.  At a continuous target policy
and the sender value it determines, a stage payoff non-increasing in the state
along the policy keeps the dynamic best response at or below the
stage ceiling at every state; so a dynamic sanction above the stage ceiling at
some state forces the composed stage payoff to rise between two feasible states
at some feasible sanction, at a supplied policy or at an equilibrium.  The
scope theorems \path{raw_chord_ne_composed_chord},
\path{existential_rise_does_not_give_uniform_chord} and
\path{dynamic_stake_conclusion_is_x_free} record that the condition concerns
the composed payoff and not the partial effect of the state, that one rise
gives no uniform slope bound, and that the conclusion does not place the rise at the
state where the sanction exceeds the ceiling.
\item \path{weakFloor_dynamicTopStepF_mem_basin},
\path{WeakFloorBasinTop.of_jointPreferenceBasinTop_reindexed} --- level-floor
invariance and its hypothesis-free containment
(\Cref{prop:level-floor}) in \path{Sanctions/Theory/Basin/WeakFloorBasin.lean}.
Invariance needs only monotonicity in each argument, feasibility and the
diagonal condition; the containment shows the weak structure is a
\emph{weakening} of the basin conditions rather than a new axis.
\path{weakFloor_ratchet_counterexample} refutes the \emph{implication} from the
level floor to a monotone sanction path --- over arbitrary kinematic data, with
a constant witness rule --- and does not refute existence of a monotone-sanction
equilibrium.
\item \path{mwD_ge_seventeen_tenths}, \path{mwSpread_pinned},
\path{mwAlphaOne_pinned}, \path{mwAlpha_pinned},
\path{mwSenderRule_pinned} --- the $P_3$ spread bound and the pinning results in
\path{Sanctions/Examples/P3WeakFloorSpreadBound.lean} and
\path{P3SpreadPinned.lean}.  The sender endpoint spread, the target's top
action, the target policy on $[2/5,1]$ and the whole sanction rule agree
across \emph{every} equilibrium of the specialization.  This is
basin-relevant behavioral agreement, \emph{not} global equilibrium
uniqueness: the target policy away from that range is not pinned.
\item \path{p3_monotone_sanctions_from_rational_start},
\path{p3_exists_ratchet_invariant_subinterval_uniform},
\path{dynamicTopLockIn_subinterval_forces_fixed_endpoint} --- strict ratcheting
on a sub-basin (\Cref{sec:weak-floor}) in
\path{Sanctions/Examples/P3SubBasinMonotonicity.lean} and
\path{Sanctions/Theory/LockIn/SubintervalLockIn.lean}.  One interval serves
every equilibrium; its right endpoint is forced to be a fixed point of the
closed loop, which is why the ratchet terminates at an interior steady state
rather than at the ceiling.
\item \path{exists_params_with_perpetual_sanctions},
\path{exists_perfect_information_perpetual_sanctions},
\path{b2_transition_is_deterministic},
\path{exists_perpetually_escalating_sanctions} --- the existence form
(the paragraph of \Cref{sec:existence:nonvacuity} that follows the
statement of \Cref{thm:b2-benchmark}) in
\path{Sanctions/Examples/PerpetualSanctionExistence.lean}.  The transition's
determinism is exhibited in the conclusion by \path{rfl}, and the transition
rate is existentially bound, so no witness magnitude appears in the statement.
\item \path{perpetual_sanctions_on_beta_interval},
\path{eeCertB_sanctions_eq_one}, \path{beta_window_nondegenerate} ---
\Cref{prop:perpetual-interval} in
\path{Sanctions/Examples/PerpetualSanctionInterval.lean}.  The conclusion holds
at \emph{every} sender discount factor in $[6/13,4/7)$, with the realized
sanction maximal at every date.  Robustness is in that one direction only; the
remaining primitives are held at the \path{eeCert} values.
\item \path{dynamicTopLockInSanctionsF_strictMono_of_gt_self},
\path{dynamicTopLockInPathF_strictMono_of_gt_self},
\path{dynamicTopLockInStepF_gt_self}, \path{strict_escalation_needs_both} ---
\Cref{prop:strict-escalation} in
\path{Sanctions/Theory/LockIn/SubintervalLockIn.lean}.  Stated over an arbitrary
economy, an arbitrary equilibrium sanction rule and an arbitrary sub-interval:
the one-step core uses only the transition's diagonal property, and the last
declaration exhibits that the two conditions are logically independent, so
neither is redundant and neither follows from the basin conditions.  The
transition hypothesis is local upward adjustment on $[a,b)$ --- a feasible
sanction above the state strictly raises it there (round BG-2) --- since the
parameter structure supplies only weak monotonicity; strict monotonicity in the
sanction and (F5) with $a\ge0$ discharge it
(\path{dynamicTopLockInSanctionsF_strictMono_of_gt_self_of_strictMono_s},
\path{dynamicTopLockInSanctionsF_strictMono_of_gt_self_of_strictAdjustment}),
and so does any affine law with positive adjustment.
\end{itemize}


\paragraph{The class under the lifting condition and the equilibrium comparison across cases (\Cref{sec:existence:placable}, \Cref{prop:activation-not-hysteresis}).}
\begin{itemize}[leftmargin=1.4em,itemsep=1pt]
\item \texttt{ChordEconomy}, \texttt{ChordEconomyRelief}, \texttt{ChordEconomyCs} (\path{Sanctions/Examples/ChordEconomy.lean}) --- The class of \Cref{sec:existence:placable} (round AU, 2026-09-25; doc 104): the transition with its slope-bound fields, both stage payoffs abstract under the rate conditions (S1)--(S11) and (P1)--(P3) of panel~A0 of \Cref{app:proofs:placable} (the fields \path{hUS_stateLip}, \path{hUS_stateSlope}, \path{hUS_concave}, \path{hUS_convex}, \path{hUS_joint}, \path{hUS_stateConvex}, \path{hUS_crossLo}, \path{hUS_crossHi}, \path{hUS_placLo}, \path{hUS_placHi}, \path{hUS_appetite}, \path{hUS_floor}, \path{hUT_costLo}, \path{hUT_costHi}, \path{hUT_painHi}, \path{hUT_painLo}, \path{hUT_stateLo}, \path{hUT_stateHi}), and the conditions \texttt{hCurv}, \texttt{hSigV}. Since round BC (2026-09-29) \path{hUS_crossHi}, \path{hUT_costHi}, \path{hUT_painLo}, \path{hUT_stateHi} and their rates form the relief extension \texttt{ChordEconomyRelief}, which \texttt{ChordEconomyCs} extends, and interiority (\path{ChordEconomy.Interior}) and the myopic relief condition (\path{ChordEconomyRelief.StageRelief}), until then fields, are predicates; the lifting condition (\path{ChordEconomy.Placable}), the entrenchment condition (\path{ChordEconomy.LossOffset}) and the three forms of compliance neutrality (\path{ZeroSanctionNeutral}, \path{ZeroSanctionAntitone}, \path{ZeroSanctionNoGain}) are predicates on the structure and hypotheses of the results that read them. Since round BH-3 (2026-10-06) the results that read interiority only through the myopic interiority take $u_{\mathrm{hi}}<q$ (\path{uhi_lt_q}), and the myopic equilibrium takes the myopic lifting condition $u_{\mathrm{hi}}<k_{\mathrm{lo}}$ (\path{uhi_lt_k}) in place of the lifting condition; interiority stays where the cap $A/q<1$ is read. The bounds $M_U$, $M_{T,0}$ are suprema over the cube, not fields. Every declaration of this paragraph after this item is stated on the class, except where it names the family; the myopic best response is the game's field $\sigma$, the lower bound $\underline s$ (\path{sfloor}) its value at the origin. The base class is strictly larger than the relief extension (round BH-5, item BH5.6): the transition, sender payoff and discount factors of every economy in the class are shared by an economy in the class that no relief economy extends, its target's per-unit sanction pain being fully offset by entrenchment at the top state (\path{exists_not_relief_same_sender}, \path{chordEconomy_not_all_relief}, \path{Examples/ChordEconomyBaseOnlyWitness.lean}). Since round BJ (2026-10-07) the base class is slimmer: \path{hUS_stateConvex}, \texttt{hFcx} and \texttt{hFmix}, with their moduli $\tau_x$, $M_{\mathrm{cx}}$ and $M_{\mathrm{mix}}$, form the two-sided extension \texttt{ChordEconomyTwoSided}, which \texttt{ChordEconomyRelief} now extends (item BJ3.1; read by the two-sided invariant set \path{carrierS2} and the policy-class invariant set \path{carrierPol}); \path{hUS_placHi} and $k_{\mathrm{hi}}$ moved to the relief extension, and on the base class the bound is the predicate \path{PlacationCeiling}, a hypothesis of the results that read it (item BJ2.2; \path{klo_le_khi}, \path{ChordEconomyRelief.placationCeiling}); $q>0$ is the theorem \path{ChordEconomy.hq}, from the concavity condition (item BJ2.1); and the joint continuity of the myopic best response, no longer part of the basic game (item BJ5.1), is the class's own field \texttt{ChordEconomy.$\sigma$\_continuous}, which supplies the instance \path{ChordEconomy.staticSelectorContinuous}, so no statement on the class changes. Every member exhibited lies in the two-sided extension and has an upper substitution rate, so this further generality is unwitnessed (Rule~9B).
\item \texttt{PlacableEconomyF} (\path{Sanctions/Examples/PlacableSenderFamilyF.lean}) --- The quadratic--linear family, a family within the class, embedded through \path{toChord}, \path{toChordCs} (the rate identifications of panel~A0; \path{toChord_lossOffset}, \path{toChord_placable}, \path{toChord_zeroSanctionNeutral}, \path{toChord_}$\sigma$\path{_zero}): the transition with its slope-bound fields (\texttt{StateChordBandAll}, \texttt{SanctionLowerChord}, \texttt{SanctionUpperChord}, \texttt{JointSemiconcaveOn}, \texttt{SemiconvexOnSet}, \texttt{MixedChordBound}), the payoff coefficients with the signs \texttt{hk}, \texttt{hv}, \texttt{hcn} ($n<c$: the sanction hurts the target at every state), and the conditions \texttt{hCurv}, \texttt{hSigV}; since round BJ (item BJ2.3, Rule~3 item~5) interiority and the myopic relief condition, until then the fields \texttt{hInterior} and \texttt{hStageRelief}, are the predicates \path{PlacableEconomyF.Interior} and \path{PlacableEconomyF.StageRelief}, hypotheses of the results that read them, and each bridge to the class takes what it reads (\path{toChord_interior}, \path{toChord_stageRelief}, or only $u_0\le q$, \path{toChord_sfloor}); the two conditions read by the relief and non-compliance cases are predicates on the structure and hypotheses of the results that read them, not fields (round AP): \path{PlacableEconomyF.Placable} ($A\le k$, hypothesis \texttt{hP}) and \path{PlacableEconomyF.LossOffset} (the entrenchment condition in its refined form: the slope \path{plGTn} carries the state--sanction cross term $n$ at the lower bound $u_0/q$ on the response; \path{hTMargin_of_old} shows the strong form implies it; hypothesis \texttt{hTM}); the closed forms \texttt{plKS}, \texttt{plAppetite}, \texttt{plGT}, \texttt{plGTn}, \texttt{plSigC}. Members: \texttt{placableAffine} and \texttt{placableQuadDiag} (\path{quadDiagF_ne_affineF}); further members \path{placableSharpRelief} (\path{Examples/PlacableSenderSharpReliefMember.lean}: sharp relief, not two-branch), \path{placableFloor} (\path{Examples/PlacableSenderFloorMember.lean}: the entrenchment condition and not the strong one, \path{exists_floor_member_not_old}) and \path{placableRallyBinding} (\path{Examples/PlacableSenderRallyBindingMember.lean}: the transition \path{quadDiagFSwap} with concave state sections, \path{quadDiagFSwap_state_not_convex}, at which the increasing-differences condition holds with equality and a strictly positive left side, \path{exists_binding_rally_member}). The enriched subclass \texttt{PlacableEconomyFcs} (\texttt{extends PlacableEconomyF}) adds the sanction-section modulus \texttt{Mcs} with \texttt{hFcs}, condition (T6) of \Cref{app:proofs:placable}; its members are \texttt{placableAffineCs} ($0$) and \texttt{placableQuadDiagCs} ($1/5$), and only the sender-only case and the theorem on the four cases are stated over it.
\item \path{VSstar}, \path{VSstar_fix}, \path{VSstar_unique}, \path{sender_channel_active}, \path{target_channel_active_on_carriers}, \path{equilibriumTop_of_margin}, \path{both_channels_active_of_margin}, \path{both_channels_active_of_margin_stage} (\path{PlacableSenderEquilibrium.lean}) --- \Cref{prop:placable-family}: the constructed sender value, both behavioral predicates (under the myopic interiority $u_{\mathrm{hi}}<q$ in the last declaration, of which the headline one is a corollary), and the composed equilibrium given the increasing-differences condition \eqref{eq:rally-margin}, a condition on members of the two-sided extension \texttt{ChordEconomyTwoSided}, on which \path{equilibriumTop_of_margin} is stated since round BJ (item BJ3.1); \path{responsePath_monotoneOn_of_margin}, also stated there, is the monotone response, through the general $\sigma$\path{_dyn_top_monotoneOn_x_of_singleCrossing} (whose uniqueness binder is unnecessary: $\sigma$\path{_dyn_top_monotoneOn_x_of_IDO} gives the same monotonicity of the dynamic best response from interval dominance alone, round I); \path{margin_of_flat} is the condition holding trivially at flat transitions; \path{responsePath_ge_sfloor} is the lower bound $u_0/q$ on the response, on which the refined target invariant set rests (\path{targetOp_mem_of_monotone}, \path{gT_mul}).
\item \path{dynamicArgmax_le_cap}, \path{response_lt_one}, \path{responsePath_gt_stage} (\path{PlacableSenderRegimes.lean}) --- the interior bounds of \Cref{prop:placable-family}, through the abstract \path{argmax_le_cap_of_chord_slope_lipschitz} (a $q$-strongly concave stage plus an $L$-Lipschitz continuation has every maximizer at most $(L_0+L)/q$; no derivative).
\item \path{myopicPolicy}, \path{myopic_target_unique}, \path{myopicEquilibriumTop}, \path{myopicRule_eq_zero}, \path{myopicPath_tendsto_zero} --- \Cref{prop:myopic-benchmark}: the myopic equilibrium as an inhabitant of \path{RegimeEquilibrium p00}, the selector identity at zero sender discount from strict stage concavity ($\sigma$\path{_dyn_top_eq_}$\sigma$\path{_of_}$\beta$\path{S_zero_of_strictConcaveOn}, \path{Theory/Mechanism/ZeroDiscountSelector.lean}), and convergence to zero; the stage forms \path{myopicEquilibriumTopS}, \path{myopicPathS_tendsto_zero} read only $u_{\mathrm{hi}}<q$ and $u_{\mathrm{hi}}<k_{\mathrm{lo}}$ (round BH-3).
\item \path{dynamicPath_ge_approach}, \path{myopicPath_lt_dynamicPath}, \path{dynamicPath_tendsto}, \path{myopic_dynamic_equilibrium_separation_generalF} --- \Cref{thm:regime-separation} and the capstone (a headline theorem; under the stage forms, without the cap clause, \path{myopic_dynamic_equilibrium_separation_generalF_stage}, of which it is a corollary); the general-$F$ path bounds under any sanction rule are \path{rulePathF_le_geom_of_zero_rule}, \path{rulePathF_ge_approach_of_rule_ge}, \path{rulePathF_gt_zero_rule_of_pos}, \path{rulePathF_tendsto_fixedPoint_of_monotoneOn}, \path{rule_eq_self_of_closedLoop_fixed} (\path{Theory/LockIn/RegimePathBounds.lean}), on the sequences of \path{Foundations/GeometricApproach.lean}.
\item \path{unique_attracting_state_of_sharp_lt_one}, \path{LloopSharp}, \path{closedLoop_lipschitz_sharp}, \path{closedLoop_iterate_dist_le_sharp}, \path{closedLoop_unique_attracting_of_sharp_lt_one} --- \Cref{prop:activation-not-hysteresis} as one class theorem (round BH-3): on every member of the relief class with the myopic interiority $u_{\mathrm{hi}}<q$, the entrenchment condition, the increasing-differences condition and the sharp loop constant below one, the equilibrium has exactly one steady state, at which the sanction equals the state, and every start converges to it; neither the myopic relief condition nor the lifting condition is read. Both members are instances: \path{placableAffine_unique_attracting_state} (a headline theorem, now proved at the sharp constant $14/23$, \path{placableAffine_LloopSharp}) and \path{placableQuadDiag_unique_attracting_state} ($35/47$, \path{placableQuadDiag_LloopSharp}), with \path{placableAffine_finite_shock_asymptotic_irrelevance} and its non-affine twin at the rates $(14/23)^n$ and $(35/47)^n$. The factor-four forms of the first formalization, through \path{Foundations/IntervalContraction.lean}, remain: \path{closedLoop_lipschitz}, \path{closedLoop_iterate_dist_le}, \path{closedLoop_unique_attracting_of_lt_one} (now a corollary of the sharp form), \path{placableAffine_Lloop} ($43/46$) and \path{placableQuadDiag_Lloop} ($277/235$, inconclusive).
\item \path{placableQuadDiagLoose}, \path{placableQuadDiagLoose_not_margin}, \path{placableQuadDiagLoose_responsePath}, \path{placableQuadDiagLoose_responsePath_monotone}, \path{exists_member_response_monotone_not_margin} (\path{Examples/PlacableSenderLooseModulusMember.lean}) --- the increasing-differences condition \eqref{eq:rally-margin} is a certificate condition, not necessary on the class (round L): \path{placableQuadDiag} re-certified with the state-section modulus $M_{\mathrm{cx}}=4$ violates it ($270/287>1/2$) while its response at non-compliance, the same function, is monotone; existence only, and silent on the tight-modulus question at \path{placableRallyBinding}.
\item \path{equilibriumTop}, \path{both_channels_active_of_response_monotone} (\path{PlacableSenderEquilibrium.lean}), \path{both_channels_active_of_relief_response_monotone} (\path{PlacableSenderReliefActivation.lean}), \path{dynamicPathOf}, \path{dynamicPath_eq_dynamicPathOf}, \path{myopic_dynamic_equilibrium_separation_generalF_stage_of_response_monotone}, \path{dynamicPath_tendsto_of_response_monotone} (\path{PlacableSenderRegimes.lean}), \path{target_only_regime_generalF_of_response_monotone}, \path{targetOnly_limit_lt_dynamic_limit_of_response_monotone}, \path{targetOnly_steady_lt_dynamic_steady_of_response_monotone} (\path{PlacableSenderRegime10.lean}), \path{four_cell_generalF'_of_response_monotone} (\path{PlacableSenderRegime01.lean}), \path{targetOnlyPath_le_dynamicPath_of_response_monotone}, \path{dynamicPath_limit_ge_limit10_add_gap_of_response_monotone}, \path{dynamicPath_limit_ge_steady10_add_gap_of_response_monotone} (\path{PlacableSenderWedge.lean}), \path{flow11Of}, \path{exists_completions_reversing_regime_ranking_of_response_monotone} (\path{PlacableSenderWelfareCompletion.lean}), \path{sanction11Of}, \path{dXOf}, \path{dSOf}, \path{dS2Of}, \path{gapOf}, \path{gapAtOf} and the \texttt{\_of\_response\_monotone} form of every case-level theorem of the welfare interval (\path{PlacableSenderWelfareInterval.lean}) --- round BH-4 (2026-10-06, item BH4.1): the increasing-differences condition \eqref{eq:rally-margin} was read only to make the response at full non-compliance, at the constructed sender value, non-decreasing in the state (\path{responsePath_monotoneOn_of_margin}). Each of these takes that monotone-response condition in place of the increasing-differences condition, is general in $\F$, and has the form under the increasing-differences condition as a one-line corollary; \path{dynamicPath} is \path{dynamicPathOf} at the discharged condition, by definition. The condition is on an endogenous object, so it is not checkable from the primitives; the increasing-differences condition implies it and is strictly stronger (\path{exists_member_response_monotone_not_margin}). \Cref{prop:activation-not-hysteresis} keeps the increasing-differences condition (\path{unique_attracting_state_of_sharp_lt_one}), and no member of the enriched subclass that fails it is exhibited for the four-case or welfare forms.
\item \path{steady10_lt_of_r_lt}, \path{vertex_lt_of_r_lt} (\path{PlacableSenderRegime10.lean}), \path{lossOffset_iff_v_ge}, \path{rally_softening_interval_iff} (\path{PlacableSenderFamilyF.lean}), \path{four_cell_rule_chain}, \path{reliefSharp_uniform_iff} (\path{PlacableSenderRegime01.lean}) --- parametric corollaries of the closed forms (round L; staged for the manuscript as N68--N71 in the notes): the vertex $(u_0+rx)/q$, which is the target-only rule under interiority, rises strictly with the state--sanction cross term $r$ at every positive state (no interiority needed), and so does the steady state $u_0/(q-r)$ under interiority of the member with the larger slope (a hypothesis since round BJ, item BJ2.3, of which only $r<q$ is read); the entrenchment condition as the valuation threshold $v\ge c(1-\beta_Tg_{\mathrm{lo}})/(\beta_Td_{\mathrm{lo}})-nu_0/q$; the four restrictions on the cross term $n$ as one interval; the common-state rule chain $R^{00}=R^{01}=0<R^{10}<R^{11}<1$; and $u_0M_{\mathrm{tot}}\le kq$ as the largest condition under which the one-period ranking condition discharges (RL) for every compliance cost.
\item \path{riccati_feasible_iff}, \path{plSigImg_le_iff_cleared}, \path{sigV_feasible_iff}, \path{exists_SigV_of_discriminant} (\path{PlacableSenderFamilyF.lean}), \path{placableAffine_riccati_feasible}, \path{placableQuadDiag_riccati_feasible} (\path{PlacableSenderWitnesses.lean}) --- the concavity condition and the Riccati closure as one condition on the primitives (round M; promoted as N73 in round N: panel A of \Cref{app:proofs:placable} and the half-clause of \Cref{sec:four-regimes}): a supplied modulus $\sigma_V\ge0$ satisfying both exists iff $Q>0$, $0<L<2Q$ and $L^2\ge4DN$, with $a=\beta_Sg_{\mathrm{hi}}^2$, $D=\beta_Sd_{\mathrm{hi}}^2$, $Q=q-\beta_SK_SM_s$, $B_0=\beta_SK_SM_x$, $C_0=r+\beta_SK_SM_m$, $H_0=\beta_Sg_{\mathrm{hi}}d_{\mathrm{hi}}$, $L=(1-a)Q+B_0D-2C_0H_0$, $N=B_0Q+C_0^2$, the vertex $L/(2D)$ then certifying; at the two members $(Q,L,N,D)=(3,5/2,1/4,1/8)$ and $(3,1623/700,25/28,9/50)$. The structure keeps its supplied modulus.
\item \path{mem_Icc_of_mono_of_diag} (\path{Foundations/DiagonalMonotone.lean}), \path{StateChordBand.rates_le} (\path{Foundations/SemiconcaveRiccati.lean}), \path{StateChordBandAll.rates_le}, \path{lower_rates_add_le_one} (\path{Foundations/SemiconcaveJointComposition.lean}), \path{sanctionChord_rates_le} (\path{Foundations/TransitionChords.lean}), \path{MixedChordBound.nonneg} (\path{Foundations/CrossDifferenceSemiconvex.lean}), \path{plKS_pos_of}, \path{plq_pos_of_interior}, \path{plSigV_nonneg_of_closure}, \path{plv_pos_of_lossOffset}, \path{PlacableEconomyF.derived_fields} (\path{PlacableSenderFamilyF.lean}) --- the seven derived fields of the family (round P, 2026-09-11; the reader's 2026-09-11 package, document 01; promoted as N75: panel C of \Cref{app:proofs:placable} and the half-clause of \Cref{sec:four-regimes}): feasibility of $\F$, $g_{\mathrm{lo}}\le g_{\mathrm{hi}}$, $d_{\mathrm{lo}}\le d_{\mathrm{hi}}$, $0\le M_{\mathrm{mix}}$, $q>0$, $\sigma_V\ge0$ and $v>0$ follow from the remaining fields, each helper stated over raw data with no member binder. Through round P they were proof fields filled by a raw-data constructor; round Q1 deleted the fields (nine contract rows rebaselined on the user's Rule-18 word) and the seven facts are theorems of the structure under their old names (\path{PlacableEconomyF.hq}, \path{hv}, \path{F_mem}, \ldots), bundled as \path{derived_fields}. Since round BJ (item BJ2.3) $q>0$ comes from the concavity condition and $\sigma_V\ge0$, interiority no longer being a field; \path{plq_pos_of_interior} is kept as a raw-data fact.
\item \path{hookF}, \path{hookF_lowerChord}, \path{hookF_cornerTop}, \path{hookF_ne_affineF}, \path{exists_nonaffine_attaining_cornerChord_bound} (\path{Examples/Counterexamples/CornerChordAttained.lean}) --- the corner bound $d\le1-\gamma$ on the lower sanction slope of \path{sanctionLowerChord_le_of_cornerTop} is attained off the affine law (round Q1; the reader's C03--C04): the law $(x+s)/2-\max(x-s,0)^2/10$ has corner weight $1/2$, lower bound on the sanction slope exactly $1/2$, fixes the diagonal, is monotone in each argument on the line and feasible on the square, and is not the affine law at rate $1/2$ ($F(1,0)=2/5$); so equality in the bound does not force affineness, which the docstring once claimed.
\item \path{senderObj_semiconvex_carrier}, \path{response_shift_ge_linear_carrier}, \path{response_dist_le_of_compliance}, \path{targetStage_small_gain}, \path{target_channel_active_of_relief}, \path{both_channels_active_of_relief} (a headline theorem), \path{placableReliefOnly}, \path{placableReliefOnly_not_flip}, \path{placableReliefOnly_relief}, \path{placableReliefOnly_both_channels}, \path{placableReliefOnly_four_cells} (\path{Examples/PlacableSenderReliefActivation.lean}), \path{rally_softening_interval_iff'} (\path{PlacableSenderFamilyF.lean}) --- both forward-looking incentives change behavior on the enriched subclass under \eqref{eq:rally-margin} and \textup{(RL)} alone (round Q2, 2026-09-11; doc 85 §3, on the user's Rule-18 word; promoted as N76 after \Cref{thm:four-cells}): the family's field became the myopic relief condition $\ell q<k(c-n)$ (through round Q1 the one-period ranking condition $\ell q<u_0(c-n)$, which implies it through the lifting condition; since round BJ, item BJ2.3, that condition is a hypothesis too, and the family keeps only its sign $n<c$); the one-period ranking condition is a hypothesis of \path{targetStage_one_gt_zero}, \path{both_channels_active_of_margin} (a headline theorem; fingerprint moved) and the two dischargers \path{reliefSharp_of_scaled_Mtot_le}, \path{reliefSharp_of_Mtot_eq_q}; \path{myopic_dynamic_equilibrium_separation_generalF} (a headline theorem; fingerprint moved) no longer carries the clause that the forward-looking incentives change behavior, so the separation needs no one-period ranking condition; the target's forward-looking incentive changes behavior without it, through a small compliance at the equilibrium value; the affine member \path{placableReliefOnly} satisfies (RL) as $6<12$ and fails the one-period ranking condition as $6<4$.
\item \path{star_premise_at_one_false} (\path{PlacableSenderRegimes.lean}) --- the pivot-star premise of \Cref{thm:persistence}, which since round BJ (item BJ1.3) is the form of its condition (P5) itself, cannot hold at a member of the family that satisfies interiority, on any basin ending at the top state (round P; the reader's document 03 \S1): at $x=1$ the pivot is $1$ and the premise makes the full sanction strictly best in the sender's objective, while every family maximizer is at most $A/q$, which interiority puts below one; a Rule-11 refutation of a hypothesis combination, so the family's persistence content is the theorem on the four cases and the temporary-shock corollary, not \Cref{thm:persistence}.
\item \path{placableAffineVertex}, \path{placableAffineVertex_LloopSharp} (\path{PlacableSenderWitnesses.lean}) --- certificate existence is not certificate interchangeability (round P; the reader's document 03 \S6): the affine member re-certified at the Riccati vertex $\sigma_V=10$ passes curvature and closure, and its sharp loop constant is exactly $1$, where the supplied $\sigma_V=1$ gives $14/23$; every quantitative bound uses one consistently chosen modulus.
\item \path{tsParams}, \path{tsClass}, \path{tsEquilibrium}, \path{tsPackage}, \path{persistence_package_not_unique_steady_state} (\path{Examples/PersistencePackageNonuniqueWitness.lean}) --- persistence does not imply attraction, at the package level (round M; promoted as N74 in round N: the sentence after \Cref{thm:persistence}): an economy ($U_S=-(s-x)^2$, $U_T=xs-a^2$, the affine transition at rate $1/2$, both discounts $1/2$), singleton value sets $\{2y^2\}$ and $\{0\}$, the constructed equilibrium $(2y^2,0,\alpha\equiv0)$ with both optima unique, every field of \path{SenderPersistencePackage} at $K_S=0$ and pivots $1/2$ --- so \Cref{thm:persistence} applies --- while the closed loop fixes every state of $[1/2,1]$ and no state is a common limit of the basin's paths; the abstract form is \path{attraction_iff_unique_fixedPoint}. Since round BH-5 (item BH5.1) the witness holds over any transition: \path{tsParamsOf} keeps the payoffs and discount factors and takes the transition of any given economy, the proof reading it only through $\F(x,x)=x$, and \path{persistence_package_not_unique_steady_state_generalF} gives the same failure of a common limit at that transition; the affine names are its instances.
\item \path{FeedbackStackelbergEquilibriumTop.toAbstractStackelbergTop}, \path{leader_optimal_against_markov_deviation_top}, \path{follower_optimal_against_markov_deviation_top}, \path{leader_discountedPayoff_eq_value_top}, \path{follower_discountedPayoff_eq_value_top} (\path{Theory/Equilibrium/TopLifetimeVerification.lean}) --- lifetime verification for the equilibrium concept of \Cref{sec:model}: each government's discounted lifetime payoff along the equilibrium equals its value, and each equilibrium policy beats every time-dependent Markov deviation against the other's fixed strategy; bounds are derived from continuity on the compact cube, none assumed. Deviations are Markov plans $\mathbb N\to[0,1]\to[0,1]$; the leader's deviation is not a stationary-policy commitment with a follower re-solve.
\item \path{TransitionExtension.extend}, \path{extend_eq_on_square}, \path{extend_continuous}, \path{extend_mono_x}, \path{extend_mono_s}, \path{extend_diag} (\path{Foundations/TransitionExtension.lean}) --- the globality remark of \Cref{sec:model}: a transition continuous, coordinatewise nondecreasing and diagonal-fixing on the square extends to the plane with the same properties and values on the square, so the formal statements lose no cube-defined economy. The payoff and selector extension by clipping, and the assembly of a formal economy from cube data, are \path{CubeData.toParams} below.
\item \path{ControlledRecursion.controlledPath}, \path{heldPath}, \path{heldPath_le_controlledPath}, \path{heldPath_natCeil_entry}, \path{controlledPath_ge_of_high_inputs}, \path{controlledPath_ge_of_inputs_ge} (\path{Foundations/ControlledRecursion.lean}) --- comparison for a scalar recursion driven by an input sequence against the recursion under a held input: with the map nondecreasing in both arguments, inputs at or above $h$ before $T$ keep the driven path above the held one through $T$; with a diagonal-fixing map the held path from $[0,h]$ is monotone, gains at least the drift floor per period below a threshold $b$, and is at or above $b$ from $\lceil b/\eta\rceil$ on. No game is named.
\item \path{episodePath}, \path{splicedPath}, \path{splicedSanction}, \path{splicedPath_eq_episode}, \path{splicedPath_add}, \path{splicedSanction_add} (\path{Theory/LockIn/FiniteEpisodePersistence.lean}) --- the sanction path of a temporary shock evolved by $\F$, and the path that follows it before date $T$ and the sanction rule $\sigDynTop_{\Vs}(x,\alpha(x))$ from $T$; the splice identity says the post-$T$ path and sanctions are the general-$\F$ equilibrium path and sanctions restarted at the state at which the shock ends. The constant-one schedule is \path{shockPathF}, the constant-zero schedule is \path{reliefPathF}.
\item \path{affine_fullReliefMinDrift_eq}, \path{affine_fullReliefExitHorizon_eq}, \path{affine_entryBound_lt_exitBound_iff}, \path{affine_entryHorizon_le_exitHorizon_of_le_half}, \path{affine_exitHorizon_le_entryHorizon_of_half_le} (\path{Examples/AffineInstance/EntryExitComparison.lean}; affine-labeled, Rule 22) --- the affine entry/exit comparison of \Cref{app:supporting-results} (round BD, 2026-10-01; formerly marked not machine-checked): $\eta^{R}_\delta(\bar x)=\delta\bar x$, the exit horizon $\lceil(1-\bar x)/(\delta\bar x)\rceil$, the exit bound above the entry bound exactly when $\bar x<\tfrac12$ at every rate, and the weak comparison of the integer horizons on each side of the midpoint. Each is a one-line instance of a statement in \path{Theory/LockIn/EntryExitComparison.lean} about transitions with $\F(1-x,1-s)=1-\F(x,s)$ on the square, of which the affine law is one.
\item \path{equilibriumTop_post_episode_nonexit_generalF}, \path{equilibriumTop_post_episode_floor_no_relief_of_ratchet}, \path{equilibriumTop_post_episode_no_relief_of_ratchet}, \path{equilibriumTop_post_episode_no_relief_of_jointBasin}, \path{equilibriumTop_fullShock_episode_nonexit_generalF_of_preBasinAscent}, \path{equilibriumTop_fullShock_episode_nonexit_generalF} --- conditional non-exit after a temporary shock (\Cref{thm:finite-episode}): given a full equilibrium and a weak-floor basin $[\bar x,1]$ with floor $\bar s$, with no sign on $\bar x$, a temporary shock ending inside it is followed by a state that stays inside, a sanction never below $\bar s$, and no relief cutoff strictly below $\bar s$; given only the ratchet $x\le R(x)$ on the basin the sanction also stays at or above the inherited state at the end of the shock, and with a monotone rule it is nondecreasing (round BH-2; the joint-preference form is the corollary); a temporary full-sanction shock of length at least \path{fullShockEntryHorizon} ends inside from every start, under ascent below the basin only. The proofs consume only the equilibrium's value and policy; the equilibrium object is what makes the resumed rule an equilibrium rule.
\item \path{highEpisodeDrift}, \path{highEpisodeMinDrift}, \path{highEpisodeHorizon}, \path{highEpisodeMinDrift_pos_attains_isMinOn}, \path{episodePath_natCeil_entry_of_drift}, \path{episodePath_uniform_highEpisodeHorizon_entry}, \path{equilibriumTop_highEpisode_nonexit_generalF}, \path{episodePath_entry_of_lowerChord} --- the non-maximal temporary shock: for a sanction level $h$ that raises every state of $[0,\bar x]$ (held-sanction ascent, which forces $\bar x<h$ and which (F5) implies, \path{lt_of_heldAscent}, \path{heldAscent_of_strictAdjustment}; round BH-2), the drift $\eta_h(\bar x)=\min_{[0,\bar x]}(\F(x,h)-x)$ is attained and positive (\path{highEpisodeMinDrift_pos_attains_isMinOn_of_heldAscent}, \path{episodePath_uniform_highEpisodeHorizon_entry_of_heldAscent}, \path{equilibriumTop_highEpisode_nonexit_generalF_of_heldAscent}; the (F5) forms above are corollaries), and for $\bar x>0$ the ascent is also necessary for entry at that horizon from every start (\path{heldAscent_of_uniform_highEpisodeHorizon_entry}); every feasible schedule held at or above $h$ for at least $\lceil\bar x/\eta_h(\bar x)\rceil$ periods ends in $[\bar x,1]$ from every start; at $h=1$ the objects are the full-shock objects by definition (\path{highEpisodeHorizon_one}); a lower bound $d$ on the sanction slope gives the looser explicit duration $\lceil\bar x/(d(h-\bar x))\rceil$.
\item \path{heldPath_affine_eq_pivot_never_enters_counterexample}, \path{no_uniform_short_duration_counterexample} --- the two refutations that delimit the temporary-shock theorem: under the affine law a sanction held \emph{at} the pivot from state $0$ never reaches it, so held-sanction ascent cannot be weakened to $h=\bar x$; and for every preassigned duration and pivot in $(0,1]$ there is an affine rate at which that many periods of full sanction end below the pivot, so ``finite'' is not ``few''.
\item \path{qbCertMS_finite_episode_nonexit}, \path{qbCertMS_fullShock_episode_nonexit}, \path{qbCertMS_highEpisode_nonexit} (\path{Examples/QuadBumpFiniteEpisode.lean}) --- the three temporary-shock statements discharged at the non-affine quadratic-bump transition through the general-transition certificate route, pivot $1/10$; existence certificates only (Rule 9).
\item \path{temporary_relief_post_relief_low_persistence_generalF}, \path{temporary_relief_post_relief_anti_generalF}, \path{two_sided_regimes_generalF} (\path{Theory/LockIn/ReliefPersistence.lean}) --- the general-transition forms of the post-relief persistence results: after the relief horizon the path evolved by $\F$ under the dynamic rule stays in $[0,\bar x]$ and is non-increasing, given the low-basin level condition; and, for one dynamic rule satisfying both level conditions, the two opposite-direction crossings with conditional persistence on each side. Since round BH-2 the exit side needs only post-basin descent, $\F(x,0)<x$ on $[\bar x,1]$, which forces $\bar x>0$ (\path{postBasinDescent_of_strictAdjustment}, \path{fullReliefMinDrift_pos_attains_isMinOn_of_postBasinDescent}, \path{reliefPathF_uniform_fullReliefExitHorizon_of_postBasinDescent}, \path{temporary_relief_post_relief_low_persistence_generalF_of_postBasinDescent}, \path{temporary_relief_post_relief_anti_generalF_of_postBasinDescent}, \path{two_sided_regimes_generalF_of_thresholds_of_localTransition}; \path{Theory/LockIn/LocalisedReliefExit.lean}): \Cref{thm:relief-exit} and \Cref{prop:post-relief-persistence} as stated. \emph{Correction 2026-09-05:} the earlier \path{temporary_relief_post_relief_low_persistence}, \path{temporary_relief_post_relief_anti} and \path{two_sided_hysteresis} evolve their post-shock path by the affine law, and the last also uses the myopic best response on its high leg and two different rules; they are retained as affine specializations and a legacy template, and no manuscript claim rests on them.
\item \path{CubeData}, \path{CubeData.toParams}, \path{toParams_U_S_on_cube}, \path{toParams_U_T_on_cube}, \path{toParams_}$\sigma$\path{_on_square}, \path{toParams_F_on_square}, \path{toParams_strictAdjustment} (\path{Postulates/ParamsOfCube.lean}) --- the globality remark of \Cref{sec:model} completed: a structure collecting exactly the cube restrictions of the fields of a formal economy (payoffs continuous on the cube, a selector feasible, continuous and argmax on the square, a transition continuous, coordinatewise monotone and diagonal-fixing on the square, its feasibility derived as \path{CubeData.F_mem} and, for every formal economy, \path{Params.F_mem} since round BC), a constructor building a formal economy from it by clipping the payoffs and selector and extending the transition, agreement with the data on the cube and square, and the transfer of (F5). Since round BJ the adjustment-rate slot $\delta$ is not data (item BJ5.2): the constructor fills it with $1/2$, which no theorem about a cube-built economy reads, since its transition is the extension of $\F$; and the selector's continuity on the square, which \texttt{Params} no longer carries (item BJ5.1), supplies the constructed economy's static-selector instance (\path{CubeData.toParams_staticSelectorContinuous}).
\item \path{Params.StaticSelectorContinuous}, \texttt{Params.$\sigma$\_continuous} (\path{Postulates/AbstractGame.lean}), \path{bergeSigma_continuous}, \path{ofStrictConcaveUSF_staticSelectorContinuous} (\path{Theory/Equilibrium/SigmaParamsRefactor.lean}), \path{ChordEconomy.staticSelectorContinuous} --- the continuous-payoff paragraph of \Cref{app:moved-ad} (round BJ, item BJ5.1, 2026-10-07; checklist item 21 and decision 43): joint continuity of the sender's myopic best response is no longer a field of \texttt{Params} but the class \path{Params.StaticSelectorContinuous}, and the former field is a theorem under it. A constant-closure check over the compiled tree, run before the change, found that 57 of the 244 census roots used the field, each a myopic-equilibrium or myopic-sender result, and that none of the twelve labels the paragraph lists did; after the change exactly three census roots take the instance as a binder (headline rows 1 and 2 and \path{standardCompliance_three_goals_from_stage_flip_generalF_of_preBasinAscent}), while the concrete economies each carry the instance, so no other statement changed. Strictly concave economies derive it by Berge's theorem; the class keeps it as its own field and supplies the instance; cube data supply it from the selector's continuity on the square. No economy with a discontinuous myopic best response is constructed (Rule~9B): the generality is stated, not witnessed.
\item \path{two_sided_regimes_generalF_of_thresholds} (\path{Theory/LockIn/ReliefPersistence.lean}), \path{dualParams_two_sided_regimes}, \path{dualParams_two_sided_episode_sanctions} (\path{Examples/DualBasinTwoSidedRegimes.lean}) --- the two-threshold form of the general-transition two-crossing composition, and its discharge at the dual-basin economy of \Cref{cor:initial-condition-selection}: one equilibrium whose rule is at the top on $[19/20,1]$ and at zero on $[0,1/10]$, so that from every start a temporary full-sanction shock of the canonical length is followed by permanent residence in the high basin with the sanction at $1$, and a temporary suspension of sanctions of the canonical length by permanent residence in the low basin with the sanction at $0$ (\Cref{prop:two-sided-one-equilibrium}; affine transition; existence only).
\item \path{dbW_S}, \path{dbW_T}, \path{dbU_S_split}, \path{dbU_T_split}, \path{db_political_parts_zero_at_full_trade}, \path{db_political_parts_monotone}, \path{db}$\sigma$\path{_eq_zero_of_le}, \path{db_full_trade_payoffs_eq_welfare}, \path{db_entrenchment_loss_after_episode}, \texttt{dualRepParams}, \path{dualRepEquilibriumTop}, \path{dualRep_rule_eq_zero}, \path{dualRep_rule_eq_zero_of_carrier} (\path{Examples/DualBasinNationalWelfare.lean}); $\sigma$\path{_dyn_top_eq_zero_of_stateFree_of_alpha_zero}, \path{dynamicTopRealizedRule_eq_zero_of_stateFree} (\path{Theory/Welfare/RepresentativeSender.lean}) --- national welfare at the dual-basin economy (\Cref{prop:entrenchment-welfare}; Lean round D, 2026-10-12): each government's payoff is its nation's welfare plus a political part ($2xs$ for the sender, $\tfrac{19}{10}x$ for the target) that is zero at full trade and non-decreasing in the state; after a full-sanction shock of the canonical length, at one equilibrium, the state stays in $[19/20,1]$ with full sanctions and full non-compliance, each nation loses in every period (discounted losses $-1/(1-\beta)$ and $-(9/10)/(1-\beta)$ at every $\beta\in[0,1)$), the sender's government gains in every period, and the target's government, at its own discount factor, gets $10x$ from the episode-end state $x$, as under permanent relief; with the sender's payoff replaced by its nation's welfare, an exact equilibrium has zero sanctions, and so does every equilibrium whose target value lies in the dual-basin target value set; the general core: a sender whose stage payoff at full non-compliance is free of the state, zero at zero sanction and negative otherwise, has value zero and best response zero against full non-compliance (one economy and one equilibrium; the comparison within the stated value set).
\item \path{targetOnlyEquilibriumTop}, \path{targetObj10_strictAntiOn}, \path{targetOnlyRule_eq_vertex}, \path{targetOnlyPath_tendsto}, \path{targetOnly_steady_lt_dynamic_steady}, \path{target_only_regime_generalF} (\path{Examples/PlacableSenderRegime10.lean}) --- the target-only case of every economy in the class that satisfies the entrenchment condition (target forward-looking, sender myopic), constructed as an equilibrium in $\mathcal E$ at the profile view: the sender's sanction rule is the myopic sanction $\sigM(x,0)$ and its value the stage value; the target's unique optimum is non-compliance and its value a Banach fixed point against the myopic best response; the state obeys the lower bound on its build-up and converges from every start to a fixed point of the myopic sanction rule; on the family, under interiority (a hypothesis since round BJ, item BJ2.3, which \path{targetOnlyRule_eq_vertex} and the family form of \path{targetOnlyPath_tendsto} take), the rule is the vertex $(u_0+rx)/q$ and the limit is $u_0/(q-r)$; under the increasing-differences condition, a condition on the two-sided extension, and the myopic interiority $u_{\mathrm{hi}}<q$, the limit is strictly below the forward-looking one (\Cref{prop:target-only}). The case with a myopic target against a forward-looking sender is the next two entries.
\item \path{BoxCoeffs}, \path{placableBox}, \path{placableBox_margin}, \path{placableBox_Lloop_lt_one}, \path{placableBoxCs}, \path{placableBoxCs_Mtot}, \path{placableBox_relief}, \path{placableBox_four_cells}, \path{placableBox_robust} (\path{Examples/PlacableSenderOpenBox.lean}) --- robustness on a coefficient box at the affine member: every vector of the nine payoff coefficients within $1/20$ of the member's, with $n\in[0,1/20]$, is an economy in the class with the same affine transition and discounts, on which the increasing-differences condition is vacuous, both forward-looking incentives change behavior at the composed equilibrium, the target-only case converges to $u_0/(q-r)$, the closed loop is a contraction with one globally attracting steady state, and, the sanction-section modulus being zero, the forward-looking relief condition holds with $M_{\mathrm{tot}}=q$ and all four cases are ordered (\Cref{app:witness}; affine-restricted; the box is a neighborhood relative to $n\ge0$).
\item \path{SemiconvexArgmaxShift.argmax_shift_ge_of_semiconvex_linear} (\path{Foundations/SemiconvexArgmaxShift.lean}) --- a derivative-free lower bound on the movement of a maximizer: if $\psi$ is $M$-semiconvex on $[0,1]$ with an interior maximizer $s'$ and $s$ maximizes $\psi+\lambda u$ for $\lambda>0$, then $s-s'\ge\min(\lambda/(2M),(1-s')/2)$. The chord through $s'-\eta$ and $s'+\varepsilon$ over the maximizer gives $\psi(s'+\varepsilon)\ge\psi(s')-(M/2)\varepsilon(\varepsilon+\eta)$; optimality against $s'+\varepsilon$ and $\eta\to0$ finish it. Interiority is essential. The sharp form (round I) is \path{argmax_shift_ge_of_semiconvex_linear_sharp}: $s-s'\ge\min(\lambda/M,1-s')$, and $\lambda/M$ outright when $s<1$ (\path{argmax_shift_ge_linear_of_semiconvex_interior}), from the two-sided quadratic lower support at an interior maximizer (\path{quadratic_support_right_of_semiconvex_interior_max}, \path{quadratic_support_left_of_semiconvex_interior_max}); attained by a quadratic, so sharp on the class.
\item \path{PlacableEconomyFcs.Mcs}, \path{hFcs} (\path{Examples/PlacableSenderFamilyF.lean}), \path{affineF_half_cs}, \path{quadDiagF_cs} --- the sanction-section semiconvexity field (T6) of \Cref{app:proofs:placable}, on the enriched subclass, and its discharge at both members ($0$ and $1/5$).
\item \path{relaxOp}, \path{relaxOp_mem2}, \path{VS01}, \path{VS01_fix}, $\rho$\path{01_continuousOn}, \path{Mtot}, \path{senderObj_semiconvex}, $\rho$\path{01_shift_ge_linear}, \path{zeroLevel}, $\rho$\path{01_zeroLevel}, \path{myopicPolicy_lt_zeroLevel}, $\rho$\path{01_ge_of_lt_sharp}, \path{tObj01_lt_of_lt}, \path{zeroLevel_argmax_unique}, \path{zeroLevel_continuousOn}, \path{senderOnlyEquilibriumTop}, \path{senderOnlyRule_eq_zero}, \path{senderOnlyPath_tendsto_zero}, \path{sender_only_regime_generalF} (\path{Examples/PlacableSenderRegime01.lean}) --- the sender-only case of every economy in the class under the lifting condition (sender forward-looking, target myopic), constructed as an equilibrium in $\mathcal E$ at the profile view. The sender's value is the value of the zero-sanction path, the fixed point of $V\mapsto w\,x+\beta_SV(\F(x,0))$ on the two-sided invariant set, constructed before the target's policy is determined because along the constructed equilibrium the sanction is zero and at sanction zero the stage payoff and the transition do not depend on compliance; the response at that value is jointly continuous, and its relief rate in compliance is bounded below through the semiconvex shift lemma with $M_{\mathrm{tot}}=q+\beta_S(\sigma_Cd_{\mathrm{hi}}^2+K_SM_{\mathrm{cs}})$; the target's unique optimum, given the forward-looking relief condition (RL), is the zero level $a^{01}(x)$, the least compliance at which the response vanishes, which is continuous in the state and strictly above $a^{00}(x)$; the realized sanction is zero and the state converges to zero. Both subclass members discharge (RL) (\path{placableAffine_relief}, \path{placableQuadDiag_relief}). Since round I the relief rate is untruncated, $\rho\ge k\Delta/M_{\mathrm{tot}}$ ($\rho$\path{01_ge_of_lt_sharp}, from $\rho$\path{01_shift_ge_linear} and the sharp shift lemma), which is what lets (RL) take the form $\ell M_{\mathrm{tot}}<k(c-n)$ in \path{tObj01_lt_of_lt'}, \path{senderOnlyEquilibriumTop'} and \path{sender_only_regime_generalF'}; the two-branch declarations named above are their corollaries through \path{reliefSharp_of_relief}, and \path{reliefSharp_of_Mtot_eq_q} discharges (RL) from the one-period ranking condition wherever $M_{\mathrm{tot}}=q$, as does \path{reliefSharp_of_scaled_Mtot_le} wherever $u_0M_{\mathrm{tot}}\le kq$ (which covers \path{placableSharpRelief}, \path{placableSharpRelief_relief_disch}). \path{placableSharpRelief} (\path{Examples/PlacableSenderSharpReliefMember.lean}) is a non-affine member with $r,n>0$ inside the sharp class and outside the two-branch one (\path{exists_sharp_relief_member_not_old}); the zero level is continuous through the general continuity of the least zero under a linear relief rate (\path{ZeroLevel.level_continuousOn}, \path{Foundations/ZeroLevel.lean}), and the relief rate on the class uses the argmax shift under a nonlinear tilt (\path{argmax_shift_ge_tilt_of_semiconvex_interior}).
\item \path{four_cell_generalF} and \path{sender_only_regime_generalF} (both headline theorems), their forms under the sharp forward-looking relief condition \path{four_cell_generalF'} and \path{sender_only_regime_generalF'} (round I; the headline forms are their corollaries), \path{four_cell_longrun_difference'} and \path{four_cell_rule_difference_at_state'} (the two difference-in-differences corollaries under the sharp forward-looking relief condition; the unprimed forms are their corollaries), \path{placableAffine_four_cells}, \path{placableQuadDiag_four_cells}, \path{placableSharpRelief_four_cells}, \path{placableFloor_four_cells}, \path{placableRallyBinding_four_cells} --- the four cases (\Cref{thm:four-cells}): the state converges to zero in cases $00$ and $01$ and to a positive steady state in cases $10$ and $11$, $u_0/(q-r)$ when the sender is myopic and a strictly higher one when it is forward-looking; the target's forward-looking incentive produces entrenchment and the sender's raises it. The difference-in-differences corollary says $(L^{11}-L^{10})-(L^{01}-L^{00})=L^{11}-L^{10}>0$ for the limit of the state and nothing more. Instantiated at both subclass members, on the affine box, and at the sharp-relief member, the member that meets the entrenchment condition but not the strong one, and the member at which the increasing-differences condition binds.
\end{itemize}

\paragraph{Behavioral predicates, the strict sender comparison, the
two-sided curvature lemmas, and pinned rows not named elsewhere (promoted
2026-09-06).}
\begin{itemize}[leftmargin=1.4em,itemsep=1pt]
\item \path{ContinuationChangesSenderChoiceAt}, \path{ContinuationChangesTargetChoiceAt} (\path{Theory/Mechanism/ContinuationChangesChoice.lean}) --- the machine-checked anchors for ``mechanism-active'' in \Cref{app:witness}: no forward-looking maximizer is stage-optimal, stated over every maximizer so no selector is baked in, and over the structure's own fields. The value-level predicate that both discounted continuation values vary across the state is \path{ContinuationStakesNondegenerateOn} (\path{Theory/Mechanism/ContinuationStakes.lean}), with \path{betaT_pos}, \path{betaS_pos}; it is weaker and is not what the paper means by the phrase.
\item \path{sender_dynamic_sanctions_gt_myopic_at_same_action}, \path{senderSanctionContinuation_strictMono_of_separated_primitives} (\path{Theory/LockIn/SenderInvestment.lean}) --- \Cref{prop:strict-dominance}: every maximizer of $J_S$ lies strictly above every stage maximizer at a common target action, given a strictly increasing continuation value at the successor state (derived from primitives) and one sanction at which the discounted continuation gain outweighs the stage loss; the target-side mirror is \path{target_dynamic_complies_lt_myopic_at_same_response}.
\item \path{Foundations/SemiconcaveJointComposition.lean}, \path{Foundations/StrongConcaveArgmaxStability.lean}, \path{Foundations/CrossDifferenceSemiconvex.lean} (\path{abs_crossDiff_le_of_twoSided}), \path{Foundations/SemiconvexComposition.lean} --- the derivative-free lemmas behind panels D--E of \Cref{app:proofs:placable}: joint semiconcavity of $v\circ\F$ from slope bounds and a joint semiconcavity bound; the maximizer of a $\mu$-strongly concave objective moves at most $L/\mu$ under an $L$-Lipschitz perturbation (\path{argmax_dist_le_of_strongConcave_lipschitz_perturbation_sharp}, round L; the factor-four form of the first formalization is kept as a corollary), and at most $2\sqrt{\varepsilon/\mu}$ under a perturbation of sup-norm size $\varepsilon$ (\path{argmax_sq_dist_le_of_strongConcave_uniform_perturbation_sharp}, $\mu|a-b|^2\le4\varepsilon$, round S from the reader's 2026-09-12 candidate; the $16\varepsilon$ form of the first formalization is kept as a corollary; H\"older-$\tfrac12$, not Lipschitz); a lower bound on the mixed second difference of $v\circ\F$ from semiconvexity of $v$ (the other side of the curvature, which a semiconcave invariant set does not give); and semiconvexity carried through a state section regular in derivative-free form and a supremum.
\item \path{wobbleEconMS_three_goals_full} (\path{Examples/SymbolicCertifiedFamilyGeneralFRecovery.lean}), \path{eeCertB_VT_star_}$\beta$\path{S_invariant} (\path{Examples/FeedbackComparativeStatics.lean}), \path{equilibriumTop_attraction_on_derived_interval_of_invariant} (\path{Theory/LockIn/BasinAttraction.lean}), \path{qbCert_three_goals_generalF} (\path{Examples/QuadBumpCertWitnessMS.lean}) --- four pinned headline rows previously named nowhere in the paper: the three lock-in goals at the wobble economy, whose transition is not affine at any weight; the target's pinned Bellman value at the certificate-complete economy is invariant in the sender's discount factor; attraction of the equilibrium closed loop on an interval derived from the model; and the three goals discharged at \path{quadBumpF}, the row that makes the general-transition certificate's generality witnessed rather than stated (\Cref{app:witness:grading}).
\end{itemize}

\paragraph{Round D after the mathematical-validation package (2026-09-06):
scope regressions, a rate, a temporary-shock corollary, and the physical-state
modulus.}
\begin{itemize}[leftmargin=1.4em,itemsep=1pt]
\item \path{exists_greatest_stage_maximizer}, \path{all_stage_lt_dynamic_of_greatest_stage_improvement} (\path{Theory/LockIn/SenderInvestment.lean}) --- \Cref{prop:strict-dominance} as stated: a greatest stage maximizer exists on the compact action set, and with the stage loss measured from it every stage maximizer lies below every maximizer of $J_S$; the existing \path{sender_dynamic_sanctions_gt_myopic_at_same_action} orders only the supplied stage maximizer; \path{all_stage_lt_dynamic_of_greatest_stage_improvement_mono} (round L) reaches the same conclusion from a non-decreasing continuation value at the successor state with no positive-discount hypothesis; since round BH-2 \path{all_stage_lt_dynamic_of_greatest_stage_improvement_of_wedge_le} needs only one comparison of continuation values at the successor state per lower sanction, $\Gamma(d)\le\Gamma(m)$ for feasible $d\le m$, with no condition on $m$, and the two earlier roots are its corollaries.
\item \path{tied_stage_maximizers_counterexample} (\path{Examples/Counterexamples/TiedStageMaximizers.lean}) --- a full economy (affine transition, sender value $2y$, stage payoff $(3/2)x-1/2-s^2(1-s)^2$) with stage maximizers $\{0,1\}$, a strictly increasing continuation value at the successor state, the stage loss measured from $0$, and the maximizer $1$ of $J_S$: the all-maximizer version with the stage loss measured from an arbitrary stage maximizer is refuted.
\item \path{placableBoundary}, \path{placableBoundary_VS01}, \path{boundary_response}, \path{boundary_zeroLevel}, \path{placableBoundary_full_compliance}, \path{exists_sender_only_full_compliance} (\path{Examples/PlacableSenderBoundaryMember.lean}) --- an economy in the class from the enriched family ($q=2$, $u_0=w=c=1$, $k=4/3$, $r=n=0$, $\ell=1/10$, $v=4$, the lifting condition at equality) at which the constructed sender-only value is identified as $4x/3$ by fixed-point uniqueness, the response is $(2/3)(1-a)$, and the zero level is $a^{01}(x)=1$: relief is bought at full compliance while $a^{00}=3/4$. Existence only.
\item \path{lossOffsetParams}, \path{lossOffsetEquilibrium}, \path{lo_argmax_unique}, \path{loss_offset_converse_counterexample} (\path{Examples/LossOffsetConverseCounterexample.lean}) --- the same payoff shapes with $v=11/4$, so $c=1>11/12=\beta_Td_{\mathrm{lo}}g_T$ and the entrenchment condition fails; a constructed equilibrium in $\mathcal E$ with $V_S=8/9+4x/3$, $V_T=-1/9+11x/3$, response $(2/3)(1-a)$ and $\alpha\equiv0$ as the unique optimum (gap $(2/45)a$). The condition is sufficient, not necessary. Built as a \texttt{Params} directly, since the family structure cannot hold it.
\item \path{ssParams}, \path{ssEquilibrium}, \path{ss_not_IDO}, \path{stage_signs_not_IDO_counterexample} (\path{Examples/StageSignsNotIDOCounterexample.lean}) --- the stage signs do not imply \Cref{ass:sender}: an economy with a positive cross difference of $\Us$ in (sanction; state), a negative one in (sanction; compliance), strong concavity in the sanction and a positive cross difference of $\Ut$ in (sanction; state), the affine transition and interior discounts, at which a constructed equilibrium in $\mathcal E$ (values zero, response $\tau(x,a)=\tfrac12+\tfrac{x}{10}-\tfrac a5$, policy $a=x$, both proved unique optima) has the substituted sender objective $-(s-(\tfrac12-\tfrac{x}{10}))^2$, which is not interval-dominant in (sanction; state): at $x=0$ the sanction $\tfrac12$ beats $\tfrac25$ and dominates the interval between them, at $x=1$ it is strictly worse. Rule 9: it refutes one implication and says nothing about the class under the lifting condition, whose \Cref{ass:sender} is discharged from its slope bounds. Since round BH-5 (item BH5.3) the economy is built over any base economy, with its transition and both discount factors (\path{ssOver}, \path{stage_signs_not_IDO_generalF}), so no condition on the transition or the discounts closes the gap.
\item \path{CenteredContraction.step_abs_le}, \path{iterate_abs_le} (\path{Foundations/CenteredContraction.lean}); \path{vertex_sub_steady10}, \path{targetOnlyPath_error_le_geometric} (\path{Examples/PlacableSenderRegime10.lean}) --- a centered geometric bound: a diagonal-fixing, coordinatewise monotone transition with lower bound $d$ on its sanction slope and a rule affine about $z$ with slope $c_0\in[0,1)$ gives $|x_t-z|\le(1-d(1-c_0))^t|x_0-z|$; on the family, under interiority (a hypothesis since round BJ, item BJ2.3, which both family declarations take) and the entrenchment condition, the target-only path converges to $u_0/(q-r)$ at rate $1-d_{\mathrm{lo}}(1-r/q)$ ($7/12$ at the affine member). A centered error bound, not a pairwise contraction.
\item \texttt{PlacableEconomyF.WelfareCompletion}, \path{govS_eq}, \path{govT_eq}, \path{welfare_eq}, \path{welfare_zero_sanction}, \path{govS_add_govT_eq}, \path{flow_sub_eq}, \path{exists_welfareCompletion}, \path{tsum_geometric_lt_of_forall_lt}, \path{exists_completions_reversing_regime_ranking}, \path{VSstar_apply_eq_sSup}, \path{senderObj_le_VSstar}, \path{senderObj_responsePath_eq_VSstar}, \path{exists_planner_efficient_completion}, \path{placableBoundary_planner_efficient_completion}, and, from round AH, \path{n_eq_zero_of_bJ_eq_bK}, \path{exists_equalValuation_completion}, \path{exists_completion_bJ_eq_bK_iff} --- equal target valuations are available exactly at the members with $n=0$ (\path{Examples/PlacableSenderWelfareCompletion.lean}) --- \Cref{prop:family-welfare-nonidentification} (round AB in Lean, promoted in round AC, 2026-09-19, on the user's word; the reader's family-welfare challenge, nodes W0--W6). An interest-group microfoundation is a new accounting object, not a hypothesis on the game: positive-weight groups, two per country, whose weighted sums equal both government payoffs at every state and action, the boundaries $r=0$ and $n=0$ included; the existing \texttt{FourGroup} structure is untouched. The ranking reversal composes the closed-form premium bound and the ordered paths of the next entry with four flow inequalities and a bounded strict comparison under discounting; the quantifier order is for every member, there exist two microfoundations, for every start and planner discount. The boundary clause unfolds the sender's fixed point into a Bellman inequality over both actions and applies \path{stationary_optimal_against_markov_deviation_of_global_bounds}; it is the Rule 11 boundary of any universal strict-inefficiency sentence. Rule 9: existence statements; they license non-identification, not magnitudes.
\item \path{min_one_shifted_vertex_le_of_isMaxOn}, \path{shifted_vertex_le_of_isMaxOn_of_lt_one} (\path{Foundations/QuadraticArgmaxLowerBound.lean}); \path{wedge_eq}, \path{responsePath_ge_stage_add_wedge}, \path{targetOnlyRule_add_wedge_le_dynamicRule}, \path{targetOnlyPath_le_dynamicPath}, \path{targetOnlySanction_add_wedge_le_dynamicSanction}, \path{placable_resumption_after_episode}, \path{myopicPath_eq_senderOnlyPath'} (\path{Examples/PlacableSenderWedge.lean}); \path{reliefPathF_lt_of_fullReliefExitHorizon_succ} (\path{Theory/LockIn/ReliefExit.lean}) --- round AB (2026-09-19; the reader's family-welfare challenge, nodes P1--P4, and the two cheap theorems of the conflict-framing reviews). A maximizer on $[0,1]$ of a concave quadratic plus a term of lower slope $L$ is at least the vertex shifted by $L/q$, clipped at one, and since round AT (2026-09-25) a maximizer of a $\bar Q$-semiconvex function with an interior maximizer $v$ plus a term of lower slope $L$ is at least $\min\{1,v+L/\bar Q\}$ (\path{min_one_shifted_argmax_le_of_semiconvex}); on the class the lower slope $g_S=w_{\mathrm{lo}}/(1-\beta_S g_{\mathrm{lo}})$ of the sender invariant set and the lower bound $d_{\mathrm{lo}}$ on the sanction slope make $L=\beta_S g_S d_{\mathrm{lo}}$, the interior cap removes the clip, and the forward-looking response exceeds the myopic sanction by at least $\beta_S g_S d_{\mathrm{lo}}/\bar Q$ at every value in the invariant set and every state ($\beta_S w d_{\mathrm{lo}}/(q(1-\beta_S g_{\mathrm{lo}}))$ on the quadratic--linear family) --- the quantitative form of the strict foresight premium of \Cref{thm:regime-separation}, displayed as \eqref{eq:sender-wedge} since round AC (N89 promoted). From a common start the forward-looking path lies above the target-only path and the realized sanctions differ by at least the same foresight premium at every date; cases $00$ and $01$ share one state path. After any feasible temporary shock, relief included, the resumed forward-looking equilibrium sanctions at least the myopic sanction plus the foresight premium at every date and its state converges to a fixed point of the response above the lower bound $\underline s$, which is what the abstract's ``durable de-escalation requires a separate condition on resumed equilibrium play'' rests on, with the low-basin obstruction \path{dynamic_no_lowBasin_containing_zero}. One more zero-sanction period makes the weak crossing of \Cref{thm:relief-exit} strict.
\item \path{senderObj_state_perturbation_lipschitz_carrier}, \path{response_dist_le_state}, \path{response_sub_le_of_le}, \path{targetObj_strictAntiOn_of_sharp}, \path{targetArgmax_eq_zero_of_sharp}, \path{sharp_of_lossOffset}, \path{targetOp_mem_lowerSlopeCone}, \path{VTdef}, \path{VTdef_eq_VTstar}, \path{VTdef_lowerSlope}, \path{equilibriumTop_sharp}, \path{gT_balance}, \path{gEff}, \path{gSharp}, \path{gSharp_balance}, \path{equilibriumTop_gSharp}; the closed sides for existence (round BG-3), \path{targetObj_antitoneOn_of_sharp_le}, \path{gSharp_balance_lo}, \path{jointRel_le_of_lt}, \path{reliefBestReply_of_relief_side_le}; on the enriched subclass \path{zeroLevelAt}, \path{response_ge_of_lt_zeroLevelAt}, \path{Lalpha}, \path{zeroLevelAt_lipschitz}, \path{KTrel}, \path{reliefOp}, \path{reliefOp_lipschitz}, \path{VTrel}, \path{VTrel_lipschitz}, \path{jointRel_lt_of_lt}, \path{zeroLevel_argmax_unique_rel}, \path{reliefEquilibriumTop}, \path{reliefEquilibriumTop_rule_eq_zero}, \path{VTdef_lipschitz}, \path{defiance_not_maxOn_of_relief_side}, \path{no_defiance_equilibrium_of_relief_side}, \path{bracket_sides_disjoint}, \texttt{SameBut$\beta$T}, \texttt{sharp\_side\_mono\_$\beta$T}, \texttt{relief\_side\_anti\_$\beta$T}, \texttt{relief\_side\_notdef\_anti\_$\beta$T}, \path{collapse_sides} (\path{Examples/PlacableSenderBracket.lean}); \path{collapse_sides'} and \path{collapse_threshold'}, under fourteen equalities (\path{Examples/PlacableSenderCollapseEquilibrium.lean}, round BC); the members \path{placableAffineLo}, \path{placableQuadDiagLo} at $\beta_T=1/8$ with their relief equilibria and the sharp-side equilibria at the pinned members (\path{Examples/PlacableSenderBracketWitnesses.lean}); the satellites \path{lowerAppetite}, \path{nonplacable_floor_of_lowerAppetite}, \path{lowerAppetite_lt_of_}$\beta$\path{S_lt}, \texttt{A\_lt\_of\_$\beta$S\_lt} (round AY), \texttt{SameButRally}, \path{sharp_side_mono_nlo}, \path{relief_side_anti_plo}, \path{relief_side_notdef_anti}, \path{sharp_side_eq_of_target_side_eq}, \path{relief_side_eq_of_target_side_eq}, \path{sharpSideSet_isOpen}, \path{reliefSideSet_isOpen}, and from round AY (2026-09-27) the disjointness of the non-compliance side from the relief-existence side, \path{bracket_existence_sides_disjoint}, with \path{glo_add_dlo_le_one} and \path{vlo_nlo_floor_le}; every slope used in a sufficient condition for non-compliance is at most $g_T$, \path{le_gT_of_balance}, \path{sharp_gT_of_sharp_of_balance}; the sufficient condition for non-compliance is inherited by a more patient target, \texttt{certified\_sharp\_mono\_$\beta$T}; the balance holds at $g_T$ where the response does not move with the state, \path{gT_balance_of_Lresp_eq_zero}; and the threshold discount factor solved in closed form, \path{glo_add_dlo_eq_one_of_tight}, \path{betaStar}, \path{collapse_threshold} (\path{Examples/PlacableSenderBracketSatellites.lean}) --- rounds AQ and AS (2026-09-23): \Cref{prop:bracket}. General in $\F$ through the slope bounds and moduli. The affine family without state--sanction cross terms, \texttt{AffineThresholdEconomy} of rounds AF--AM, with its exact threshold, sharpness, partial-relief accounting, results without the lifting condition, extension with cross terms and affine-value uniqueness, is retired (doc 102 J1, J4): every result of it that has a general form is stated above, and its exact closed forms are \path{collapse_sides}.
\item \texttt{$\sigma$\_dyn\_top\_antitoneOn\_a\_of\_DD}, \path{V_T_eq_at_policy}, \path{bellmanS_BCF_congr_on} (\path{Theory/Equilibrium/CertifiedEquilibria.lean}); \path{gT}$_0$, \path{LossOffset}$_0$, \path{rule}, \path{V_S_monotone_of_antitone} (in the file but not read by the uniqueness chain, which reads no neutrality), \path{V_T_eq}, \path{deficit}, \path{deficit_ge}, \path{target_lowerSlope_of_monotone_class}, \path{policy_eq_zero_of_monotone_class}, \path{unique_equilibrium_of_monotone_class}, \path{equilibriumTop_mem_monotone_class} (\path{Examples/PlacableSenderUniqueness.lean}) --- round AR (2026-09-23): \Cref{prop:monotone-uniqueness}; since round BC its class is the monotone sanction rule alone, through \path{slopeDeficit_step}, \path{lowerSlope_of_rule_monotone} (\path{Theory/Equilibrium/DefianceCore.lean}), \path{values_eq_of_policy_zero} and \path{target_lowerSlope_of_rule_monotone}, \path{policy_eq_zero_of_rule_monotone}, \path{unique_equilibrium_of_rule_monotone} (\path{Examples/PlacableSenderDefianceCore.lean}).
\item \path{summable_geometric_mul_of_abs_le}, \path{withSender}, \texttt{WeightBounded}, \texttt{HawkBounded}, \path{weightBounded_iff}, \path{hawkBounded_iff}, \path{etaLo}, \path{etaHi}, \path{hawkBounded_eta_mem}, \path{exists_withSender_hawkBounded_eta}, \path{exists_hawkBounded_eta}, \path{hawkBounded_nonempty_iff}, \path{dX}, \path{dS}, \path{dS2}, \path{regimeValue_flow_eq}, \path{regimeValue11_sub_10_eq}, \path{dS_ge}, \path{dS_pos}, \path{gap}, \path{gapAt}, \path{gap_mem_interval}, \path{gap_attained_withSender}, \path{gap_attained}, \path{gap_pos_of_lower_pos}, \path{gap_neg_of_upper_neg}, \path{exists_both_signs_of_straddle}, \path{exists_weightBounded_eta_gt}, \path{exists_weightBounded_eta_lt}, \path{exists_weightBounded_reversal}, and, from round AH, \path{gapAt_lo_le_hi}, \path{forall_gap_pos_iff}, \path{forall_gap_neg_iff}, \path{exists_both_signs_iff} --- what the interval identifies, as three equivalences that differ at a zero endpoint --- and the \texttt{withSender} forms \path{exists_withSender_both_signs_of_straddle}, \path{exists_withSender_weightBounded_eta_gt}, \path{exists_withSender_weightBounded_eta_lt}, \path{exists_withSender_weightBounded_reversal}, of which the four older statements are now corollaries (\path{Examples/PlacableSenderWelfareInterval.lean}) --- round AG: \Cref{prop:welfare-interval}. \texttt{HawkBounded} is a definition over the microfoundation's own data, not a hypothesis on the game; the two \texttt{withSender} statements name the attaining microfoundation, so ``target side unchanged'' is in the statement and not only in the proof.
\item \path{HistoryPlan}, \path{historyRun}, \path{historyDiscountedPayoff}, \path{replayPlan}, \path{historyDiscountedPayoff_eq_replay}, \path{historyState_toHistoryPlan}, \path{stationary_optimal_against_history_deviation_of_global_bounds}, \path{leader_optimal_against_history_deviation}, \path{follower_optimal_against_history_deviation}, \path{follower_node_optimal_against_history_deviation} (\path{Foundations/HistoryDependentDeviations.lean}), with the three \path{_top} statements of \path{Theory/Equilibrium/TopLifetimeVerification.lean} --- round AE: in a deterministic discounted problem a pure history-dependent plan and the dated plan replaying its own action sequence have the same states, actions and payoff, so every optimality statement against time-dependent Markov plans holds against pure history-dependent plans, for both players of an abstract Stackelberg feedback equilibrium and at the follower's node after an arbitrary observed leader action (\Cref{prop:history-deviations}). Scope: pure strategies, deterministic transition; no characterization of other equilibria.
\item \path{sum_add_terminal_add_residual_le}, \path{sum_add_terminal_add_gap_le}, \path{one_le_geom_sum} (\path{Foundations/FiniteDiscountedTelescope.lean}); \path{reliefState}, \path{senderReliefValue}, \path{targetReliefValue}, \path{sender_relief_step}, \path{target_relief_step}, \path{senderReliefValue_add_loss_le}, \path{targetReliefValue_add_cost_le}, \path{targetReliefValue_le_defiance}, \path{relief_joint_loss}, \path{no_date_zero_transfer_supports_relief}, \path{relief_participation_of_margin}, \path{lossOffset_relief_period_benefits_target} (\path{Examples/PlacableSenderReliefParticipation.lean}) --- round AE: \Cref{prop:relief-participation}, general in the transition, with its refutation outside the entrenchment condition at the affine economy \path{lossOffsetParams}; its statement at the target's own optimum, with no condition at zero sanction, is \path{relief_joint_loss_at_defiance} and \path{no_date_zero_transfer_supports_relief_at_defiance} (\path{Examples/PlacableSenderReliefAtDefiance.lean}, round BC); at every full-non-compliance equilibrium, with a non-decreasing response in place of the increasing-differences condition, \path{relief_participation_of_policy_zero_at_defiance} and \path{relief_participation_of_policy_zero} (\path{Examples/PlacableSenderDefianceEquilibriaWelfare.lean}, round BH-3), of which the forms under that condition are corollaries.
\item \path{senderObj_zero_eq_of_sender_side_eq}, \path{bellmanS_zero_eq_of_sender_side_eq}, \path{VSstar_eq_of_sender_side_eq}, \path{responsePath_eq_of_sender_side_eq}, \path{dynamicPath_eq_of_sender_side_eq}, \path{dynamicSanction_eq_of_sender_side_eq} (\path{Examples/PlacableSenderBranchInvariance.lean}) --- round AE: \Cref{prop:branch-invariance} for the constructed equilibria, through uniqueness of the Bellman fixed point on the ambient value space; for any full-non-compliance equilibria of any two games, \path{bellmanS_BCF_zero_eq_of_defiance_section_eq}, \path{V_S_eq_V_S_zero_of_policy_zero}, \path{senderValue_eq_of_defiance_section_eq}, \path{equilibriumRule_eq_of_defiance_section_eq}, \path{equilibriumPath_eq_of_defiance_section_eq}, \path{defiance_branch_invariance_generalF} (\path{Theory/Equilibrium/BranchInvarianceParams.lean}, round BC).
\item \path{fixedPoint_ge_steady10_add_gap}, \path{steadyGap_eq}, \path{dynamicPath_limit_ge_steady10_add_gap} (\path{Examples/PlacableSenderWedge.lean}); \path{boundary_bellmanS_loVS}, \path{placableBoundary_VSstar}, \path{placableBoundary_responsePath}, \path{placableBoundary_wedge_attained}, \path{wedge_coefficient_cannot_be_raised} (\path{Examples/PlacableSenderWedgeSharp.lean}) --- round AE: \Cref{prop:steady-gap} and \Cref{rem:wedge-sharp}.
\item \path{exists_planner_efficient_completion_of_policy_zero} (\path{Examples/PlacableSenderDefianceEquilibriaWelfare.lean}), \path{flat_of_margin_of_r_eq_zero}, \path{exists_planner_efficient_completion_at_equilibrium}, \path{exists_planner_efficient_completion_of_flat}, \path{placableBoundary_planner_efficient_at_equilibrium} (\path{Examples/PlacableSenderWelfareCompletion.lean}) --- \Cref{cor:planner-efficient-boundary} at every full-non-compliance equilibrium, with no increasing-differences condition (round BH-3); the round-AE form at the constructed government equilibrium is its corollary; at $r=0$ the increasing-differences condition forces $M_{\mathrm{mix}}=0$ and $M_{\mathrm{cx}}g_{\mathrm{hi}}=0$, so through that condition the boundary reaches only flat transitions.
\item \path{rule_ge_sfloor_of_policy_zero}, \path{placable_post_episode_nonexit_of_policy_zero}, \path{placable_highEpisode_nonexit_of_policy_zero}, \path{no_lowBasin_containing_zero_of_policy_zero}, \path{four_cell_rule_chain_of_policy_zero} (\path{Examples/PlacableSenderDefianceEpisode.lean}, round BH-3); \path{dynamicRule_ge_sfloor}, \path{toParams_strictAdjustment}, \path{placable_post_episode_nonexit_generalF}, \path{placable_highEpisode_nonexit_generalF} (\path{Examples/PlacableSenderEpisode.lean}) --- \Cref{thm:finite-episode} on the class at every full-non-compliance equilibrium, with pivot and floor $\sigM(0,0)$ ($u_0/q$ on the family): every such equilibrium realizes the constructed response, which exceeds the myopic best response, the diagonal gives $\F(b,b)=b$, and strict adjustment follows from the positive lower bound on the sanction slope; the low-basin obstruction and the case-$11$ link of the rule chain across the four cases hold at every such equilibrium too, without the increasing-differences condition. The forms at the constructed equilibrium, which read the entrenchment and increasing-differences conditions only to name it, are corollaries; a full-non-compliance equilibrium exists exactly when full non-compliance is a target best reply at the constructed values. Weak invariance and a sanction floor; not the ratchet, not hysteresis.
\item \path{TransitionStateLipschitzOnWith}, \path{TransitionStateLipschitzWith.toOn} (\path{Theory/Basin/PrimitiveHighBasin.lean}); \path{V_S_valueLipschitzWith_of_transitionStateLipschitzOn} and the \path{_on} chain (\path{Theory/Bellman/ValueLipschitz.lean}) --- the inherited-state modulus (F6-Q) quantified over feasible states only, which is the form the sender value-Lipschitz route consumes; the global predicate and its theorems remain as specializations.
\item \path{MonotoneDifferenceArgmax.isMaxOn_max_of_monotone_diff} (\path{Foundations/MonotoneDifferenceArgmax.lean}) --- the Milgrom--Shannon greatest-argmax preservation lemma, factored out of \path{Theory/LockIn/SigmaDynComparison.lean}, which keeps its name as the $[0,1]$ instance.
\item \path{InterestGroupRep}, \path{payoffs_add_eq}, \path{rescaleS}, \path{rescaleS_welfare}, \path{tsum_pow_mul_lt_of_forall_lt}, \path{exists_scale_ranks}, \path{exists_rescaleS_ranks} (\path{Theory/Welfare/InterestGroup.lean}); \path{zeroI}, \path{statePath_stationary_succ}, \path{discountedPayoff_stationary_eq}, \path{abs_discountedPayoff_stationary_le}, \path{plannerFlow}, \path{plannerStep}, \path{zeroSanctionPlan}, \path{zeroPath}, \path{zeroPath_ordered}, \path{zeroSanction_planner_optimal} (\path{Theory/Welfare/PlannerZeroSanction.lean}); \path{discountedWelfare}, \path{U_S_le_defiance}, \path{senderWelfareRep}, \path{senderWelfareRep_welfare}, \path{repFlow}, \path{cubeStep}, \path{defiancePlan}, \path{exists_planner_efficient_rep}, \path{planner_efficient_rep_at_policy_zero}, \path{etaRep}, \path{etaRep_welfare}, \path{repFlow00}, \path{repFlow01}, \path{repFlow10}, \path{repFlow11}, \path{exists_reps_reversing_regime_ranking}, \path{repFlowMyopic}, \path{repFlowForward}, \path{exists_reps_reversing_myopic_forward_ranking} (round W2, 2026-10-16; \path{Examples/ChordEconomyWelfareGeneral.lean}) --- welfare in general abstract form (round W, 2026-10-15): interest-group representations of arbitrary payoffs, the accounting identity, rescaling one group, a sufficient condition for zero sanctions to be socially optimal, and the class-wide non-identification and planner-efficiency results of \Cref{sec:welfare}.
\end{itemize}

\endgroup
